\documentclass[11pt, letterpaper]{article}
\usepackage[utf8]{inputenc}
\usepackage[T1]{fontenc}
\usepackage{lmodern}
\renewcommand{\familydefault}{\sfdefault}
\usepackage[spanish]{babel}
\usepackage[letterpaper, top=1.5cm, bottom=1.5cm, left=1.27cm, right=1.27cm]{geometry}
\usepackage{titlesec}
\usepackage{enumitem} 
\usepackage{xcolor}
\usepackage{booktabs}
\usepackage{tabularx}
\usepackage{array}
\usepackage{float}
\usepackage{amsmath}
\usepackage{longtable}
\usepackage{multirow}
\usepackage{graphicx}
\usepackage[hidelinks]{hyperref}

\linespread{1.1}

\titleformat{\section}{\Large\bfseries}{\thesection.}{0.5em}{}
\titleformat{\subsection}{\large\bfseries}{\thesubsection.}{0.5em}{}
\titleformat{\subsubsection}{\normalsize\bfseries}{\thesubsubsection.}{0.5em}{}
\titlespacing*{\section}{0pt}{18pt}{6pt}
\titlespacing*{\subsection}{0pt}{12pt}{4pt}
\titlespacing*{\subsubsection}{0pt}{8pt}{2pt}

\setlength{\parindent}{0pt}
\setlength{\parskip}{4pt}
\renewcommand{\arraystretch}{1.3}
\pagestyle{plain}

\begin{document}

\begin{titlepage}
    \centering
    \vspace*{1.5cm}

    {\Huge \textbf{SkyAnalytics Inc.}}\\[1cm]
    {\footnotesize\textbf{Drive:} \texttt{https://drive.google.com/drive/folders/1f9FfqWsVkM63CCfhutV755ouBE61dDS6?usp=sharing}}\\[0.2cm]
    {\footnotesize \textbf{Github:} \texttt{https://github.com/khrizzcoronel/SkyAnalytics.git}}\\[0.8cm]
    {\Large \textbf{Perfil Estratégico y Corporativo}}\\[0.3cm]

    {\normalsize \textit{Análisis Documental Estratégico}}\\[0.3cm]
    {\normalsize \textit{Khriz Coronel}}\\[0.3cm]
    {\normalsize \textit{Versión 3.0}}\\[0.3cm]

    

\end{titlepage}

\tableofcontents
\newpage

%=====================================================================
% I. DESARROLLO EMPRESARIAL Y NIVELES ORGANIZACIONALES
%=====================================================================

\section{Desarrollo y Organización Empresarial}

\subsection{Nivel Estratégico}

Define el rumbo de la organización a largo plazo. La alta dirección (CEO, CTO) traza la misión, visión, alianzas estratégicas y objetivos globales. Se materializa en 10 Objetivos Estratégicos (OE1--OE10). Las decisiones en este nivel determinan el posicionamiento competitivo de SkyAnalytics como referente mundial en inteligencia aeronáutica.

\subsection{Nivel Táctico}

Traduce la estrategia en planes concretos por área a mediano plazo. Los líderes de área (VP Marketing, CTO, SRE Lead, Data Engineer Lead, ML Engineer Lead) definen campañas, presupuestos, configuración de infraestructura y metodologías. Se concreta en 22 Objetivos Tácticos (OT1--OT22).

\subsection{Nivel Operativo}

Ejecución de tareas diarias y semanales a corto plazo. Ingenieros, analistas y SREs sostienen la operación: pipelines de datos, entrenamiento de modelos, monitoreo de infraestructura, atención de tickets y documentación. Se especifica en 34 Objetivos Operativos (OP1--OP34).




Los actores del sistema se dividen en los niveles Estratégico, Táctico, Operativo y Externos. Véase la tabla detallada de Actores del Sistema en la Sección 3.


%=====================================================================
% II. ESTRATEGIA DE SkyAnalytics Inc.
%=====================================================================

\section{Estrategia de SkyAnalytics Inc.}
Actividad: ingesta y análisis profundo de macrodatos (Big Data) de vuelos históricos y en tiempo real, APIs predictivas REST/GraphQL, modelos de Machine Learning para predicción de retrasos, clima severo y optimización de rutas, dashboards self-service para aerolíneas, agencias de viaje, operadores logísticos y entidades gubernamentales. Plataforma SaaS con modelo de suscripción corporativa.


\subsection{Misión}

Proveer inteligencia accionable, precisa y oportuna al ecosistema de la aviación global, permitiendo a nuestros clientes optimizar drásticamente sus operaciones comerciales, minimizar la incertidumbre en sus procesos logísticos y elevar la calidad de la experiencia que ofrecen a sus propios pasajeros mediante decisiones puramente guiadas por datos.

\subsection{Visión}

Posicionarnos para el año 2028 como la firma líder y el referente tecnológico mundial en el análisis predictivo del mercado aeronáutico. Aspiramos a ser el ecosistema de datos subyacente y el aliado estratégico indispensable para cualquier actor de la industria que busque operar con la máxima eficiencia y absoluta confiabilidad técnica.

\subsection{Objetivo Estratégico General}

Convertir datos aeronáuticos globales en inteligencia accionable y rentable para cada actor de la industria, mediante una plataforma SaaS de APIs predictivas, infraestructura cloud de alta disponibilidad (99.99\%), modelos de Machine Learning auditables con explainability (SHAP) y una cultura Remote-First de innovación continua, asegurando crecimiento financiero sostenido, excelencia operativa y confianza normativa absoluta (SOC 2, ISO 27001, GDPR, IATA).

%---------------------------------------------------------------------
\subsection{Balanced Scorecard por Perspectiva}

\subsubsection{Perspectiva Financiera}

\begin{table}[H]
\small
\centering
\begin{tabularx}{\linewidth}{>{\raggedright\arraybackslash}p{2.2cm} >{\raggedright\arraybackslash}p{1.8cm} >{\raggedright\arraybackslash}p{2.5cm} >{\raggedright\arraybackslash}p{2cm} >{\raggedright\arraybackslash}X >{\raggedright\arraybackslash}p{2cm}}
\toprule
\multicolumn{1}{>{\raggedright\arraybackslash}p{2.2cm}}{\textbf{Objetivo}} &
\multicolumn{1}{>{\raggedright\arraybackslash}p{1.8cm}}{\textbf{Indicador}} &
\multicolumn{1}{>{\raggedright\arraybackslash}p{2.5cm}}{\textbf{Fórmula}} &
\multicolumn{1}{>{\raggedright\arraybackslash}p{2cm}}{\textbf{Meta}} &
\multicolumn{1}{>{\raggedright\arraybackslash}X}{\textbf{Iniciativa Estratégica}} &
\multicolumn{1}{>{\raggedright\arraybackslash}p{2cm}}{\textbf{Responsable}} \\
\midrule
Incrementar ARR & ARR & Suma suscripciones activas $\times$ 12 & Duplicar anualmente & Growth Hacking, automatización marketing, SDKs & CEO \\
\midrule
Reducir CAC & Costo Adquisición & Gasto marketing / clientes nuevos & \textless\ \$500 enterprise & SEO, contenido técnico, comunidad & VP Marketing \\
\midrule
Mejorar LTV/CAC & Ratio LTV/CAC & LTV / CAC & \textgreater\ 3x & Customer Success automatizado, expansión & VP Customer Success \\
\midrule
Sostener margen bruto & Gross Margin & (ARR - costo cloud) / ARR $\times$ 100 & \textgreater\ 75\% & Optimización infra, serverless mixta & CTO \\
\midrule
Expandir clientes & NRR & (ARR inicial + expansión - churn) / ARR inicial $\times$ 100 & \textgreater\ 110\% & Planes escalonados, upselling & VP Ventas \\
\bottomrule
\end{tabularx}
\caption{KPIs de la Perspectiva Financiera}
\end{table}

%---------------------------------------------------------------------
\subsubsection{Perspectiva del Cliente}

\begin{table}[H]
\small
\centering
\begin{tabularx}{\linewidth}{>{\raggedright\arraybackslash}p{2.2cm} >{\raggedright\arraybackslash}p{1.8cm} >{\raggedright\arraybackslash}p{2.5cm} >{\raggedright\arraybackslash}p{2cm} >{\raggedright\arraybackslash}X >{\raggedright\arraybackslash}p{2cm}}
\toprule
\multicolumn{1}{>{\raggedright\arraybackslash}p{2.2cm}}{\textbf{Objetivo}} &
\multicolumn{1}{>{\raggedright\arraybackslash}p{1.8cm}}{\textbf{Indicador}} &
\multicolumn{1}{>{\raggedright\arraybackslash}p{2.5cm}}{\textbf{Fórmula}} &
\multicolumn{1}{>{\raggedright\arraybackslash}p{2cm}}{\textbf{Meta}} &
\multicolumn{1}{>{\raggedright\arraybackslash}X}{\textbf{Iniciativa Estratégica}} &
\multicolumn{1}{>{\raggedright\arraybackslash}p{2cm}}{\textbf{Responsable}} \\
\midrule
Mejorar satisfacción & NPS & \% promotores - \% detractores & \textgreater\ 50 enterprise / \textgreater\ 30 SMB & Encuestas trimestrales, feedback loop & VP Customer Success \\
\midrule
Reducir churn & Churn Rate & Clientes cancelados / total $\times$ 100 & \textless\ 0.5\% enterprise / \textless\ 2\% SMB & Health scoring, alertas riesgo, onboarding & Customer Success Mgr \\
\midrule
Acelerar valor & TTFV & Tiempo registro $\rightarrow$ primer insight & \textless\ 15 min & Sandbox instantáneo, datos sintéticos & Developer Advocate \\
\midrule
Calidad soporte & CSAT & Satisfechos / encuestados $\times$ 100 & \textgreater\ 90\% & Chatbot IA nivel 1, portal autoservicio & VP Customer Success \\
\bottomrule
\end{tabularx}
\caption{KPIs de la Perspectiva del Cliente}
\end{table}

%---------------------------------------------------------------------
\subsubsection{Perspectiva de Procesos Internos}

\begin{table}[H]
\small
\centering
\begin{tabularx}{\linewidth}{>{\raggedright\arraybackslash}p{2.2cm} >{\raggedright\arraybackslash}p{1.8cm} >{\raggedright\arraybackslash}p{2.5cm} >{\raggedright\arraybackslash}p{2cm} >{\raggedright\arraybackslash}X >{\raggedright\arraybackslash}p{2cm}}
\toprule
\multicolumn{1}{>{\raggedright\arraybackslash}p{2.2cm}}{\textbf{Objetivo}} &
\multicolumn{1}{>{\raggedright\arraybackslash}p{1.8cm}}{\textbf{Indicador}} &
\multicolumn{1}{>{\raggedright\arraybackslash}p{2.5cm}}{\textbf{Fórmula}} &
\multicolumn{1}{>{\raggedright\arraybackslash}p{2cm}}{\textbf{Meta}} &
\multicolumn{1}{>{\raggedright\arraybackslash}X}{\textbf{Iniciativa Estratégica}} &
\multicolumn{1}{>{\raggedright\arraybackslash}p{2cm}}{\textbf{Responsable}} \\
\midrule
Garantizar disponibilidad & Uptime & (Tiempo total - caído) / total $\times$ 100 & $\geq$ 99.99\% & Multi-región activo-activo, CDN, chaos engineering & SRE Lead \\
\midrule
Gestionar riesgo & Error Budget & Tiempo indisponible / tiempo total mensual & $\leq$ 80\% consumido (4.38 min/mes) & Congelar deploys si \textgreater\ 80\% & SRE Lead \\
\midrule
Frescura de datos & Data Freshness & \% datos en API \textless\ 5 min & \textgreater\ 95\% & Streaming pipelines, monitoreo lag & Data Engineer Lead \\
\midrule
Recuperación rápida & MTTR & Tiempo recuperación / \# incidentes Sev1/2 & \textless\ 15 min & Runbooks automatizados, on-call, post-mortems & SRE Lead \\
\bottomrule
\end{tabularx}
\caption{KPIs de la Perspectiva de Procesos Internos}
\end{table}

%---------------------------------------------------------------------
\subsubsection{Perspectiva de Aprendizaje y Crecimiento}

\begin{table}[H]
\small
\centering
\begin{tabularx}{\linewidth}{>{\raggedright\arraybackslash}p{2.2cm} >{\raggedright\arraybackslash}p{1.8cm} >{\raggedright\arraybackslash}p{2.5cm} >{\raggedright\arraybackslash}p{2cm} >{\raggedright\arraybackslash}X >{\raggedright\arraybackslash}p{2cm}}
\toprule
\multicolumn{1}{>{\raggedright\arraybackslash}p{2.2cm}}{\textbf{Objetivo}} &
\multicolumn{1}{>{\raggedright\arraybackslash}p{1.8cm}}{\textbf{Indicador}} &
\multicolumn{1}{>{\raggedright\arraybackslash}p{2.5cm}}{\textbf{Fórmula}} &
\multicolumn{1}{>{\raggedright\arraybackslash}p{2cm}}{\textbf{Meta}} &
\multicolumn{1}{>{\raggedright\arraybackslash}X}{\textbf{Iniciativa Estratégica}} &
\multicolumn{1}{>{\raggedright\arraybackslash}p{2cm}}{\textbf{Responsable}} \\
\midrule
Satisfacción equipo & eNPS & \% promotores - \% detractores (interna) & \textgreater\ 40 & Cultura Remote-First, feedback 360 & CEO / VP People \\
\midrule
Onboarding rápido & Time-to-Productivity & Días hasta primer deploy & \textless\ 5 días hábiles & Documentación interna, mentor, entorno automatizado & CTO \\
\midrule
Retener talento & Tasa Retención & Empleados permanecen / total inicio $\times$ 100 & \textgreater\ 90\% anual & Plan carrera IC, formación, revisiones & VP People \\
\midrule
Innovación abierta & Publicaciones Técnicas & Artículos + charlas + contribuciones / trimestre & $\geq$ 4 & Presupuesto conferencias, blog técnico & CTO \\
\bottomrule
\end{tabularx}

\vspace{0.15cm}
{\footnotesize
\textbf{Nota metodológica eNPS:} Medición trimestral mediante encuesta interna anónima. \\
Umbral mínimo de participación: 70\,\% del equipo. \\[0.1cm]
Interpretación: \quad $> 40$ = zona de excelencia \quad|\quad $10$--$40$ = zona de mejora \quad|\quad $< 10$ = zona crítica.
}
\caption{KPIs de la Perspectiva de Aprendizaje y Crecimiento}
\end{table}

\subsection{Cuadro Resumen del Balanced Scorecard}

\begin{table}[H]
\small
\centering
\begin{tabularx}{\linewidth}{>{\raggedright\arraybackslash}p{2.5cm} >{\raggedright\arraybackslash}p{3.5cm} >{\raggedright\arraybackslash}p{3cm} >{\raggedright\arraybackslash}X}
\toprule
\multicolumn{1}{>{\raggedright\arraybackslash}p{2.5cm}}{\textbf{Perspectiva}} &
\multicolumn{1}{>{\raggedright\arraybackslash}p{3.5cm}}{\textbf{Objetivo Principal}} &
\multicolumn{1}{>{\raggedright\arraybackslash}p{3cm}}{\textbf{Indicador Clave}} &
\multicolumn{1}{>{\raggedright\arraybackslash}X}{\textbf{Meta Estratégica}} \\
\midrule
Financiera & Incrementar ARR & ARR & Duplicar anualmente \\
\midrule
Financiera & Reducir adquisición & CAC & \textless\ \$500 enterprise \\
\midrule
Financiera & Eficiencia comercial & LTV/CAC & \textgreater\ 3x \\
\midrule
Financiera & Margen bruto & Gross Margin & \textgreater\ 75\% \\
\midrule
Financiera & Expandir clientes & NRR & \textgreater\ 110\% \\
\midrule
Cliente & Satisfacción & NPS & \textgreater\ 50 enterprise \\
\midrule
Cliente & Reducir churn & Churn Rate & \textless\ 0.5\% enterprise \\
\midrule
Cliente & Velocidad valor & TTFV & \textless\ 15 min \\
\midrule
Cliente & Soporte & CSAT & \textgreater\ 90\% \\
\midrule
Procesos & Disponibilidad & Uptime & $\geq$ 99.99\% \\
\midrule
Procesos & Riesgo & Error Budget & $\leq$ 80\% \\
\midrule
Procesos & Frescura & Data Freshness & \textgreater\ 95\% \\
\midrule
Procesos & Recuperación & MTTR & \textless\ 15 min \\
\midrule
Aprendizaje & Equipo & eNPS & \textgreater\ 40 \\
\midrule
Aprendizaje & Onboarding & Time-to-Productivity & \textless\ 5 días \\
\midrule
Aprendizaje & Retención & Retención anual & \textgreater\ 90\% \\
\midrule
Aprendizaje & Innovación & Publicaciones & $\geq$ 4/trim \\
\bottomrule
\end{tabularx}
\caption{Cuadro resumen del Balanced Scorecard de SkyAnalytics}
\end{table}

%=====================================================================
% III. GESTIÓN Y CASOS DE USO DEL SISTEMA
%=====================================================================

\section{Gestión y Casos de Uso del Sistema}

En esta sección se detallan los objetivos y casos de uso que dan vida operativa a la estrategia del sistema propuesto.

\subsection{Actores del Sistema}

\begin{table}[H]
\small
\centering
\begin{tabularx}{\linewidth}{>{\raggedright\arraybackslash}p{3.5cm} >{\raggedright\arraybackslash}X >{\raggedright\arraybackslash}p{2.5cm}}
\toprule
\multicolumn{1}{>{\raggedright\arraybackslash}p{3.5cm}}{\textbf{Actor}} &
\multicolumn{1}{>{\raggedright\arraybackslash}X}{\textbf{Descripción}} &
\multicolumn{1}{>{\raggedright\arraybackslash}p{2.5cm}}{\textbf{Nivel}} \\
\midrule
CEO/Founder & Define visión, BSC, alianzas & Estratégico \\
\midrule
CTO & Arquitectura, stack, decisiones técnicas, presupuesto cloud & Estratégico \\
\midrule
VP Customer Success & Estrategia retención, churn, NPS, expansión & Táctico \\
\midrule
VP Marketing & Campañas, growth hacking, contenido, comunidad & Táctico \\
\midrule
VP People & Cultura Remote-First, retención talento, eNPS interno, onboarding & Táctico \\
\midrule
Data Engineer & Pipelines ETL, data warehouse, calidad, observabilidad & Táctico/Operativo \\
\midrule
ML Engineer & Entrena, versiona, despliega, monitorea modelos ML & Táctico/Operativo \\
\midrule
SRE/DevOps & Infra cloud, CI/CD, monitoreo, alertas, seguridad & Táctico/Operativo \\
\midrule
Developer Advocate & Developer Portal, SDKs, docs, Sandbox, comunidad & Táctico/Operativo \\
\midrule
Cliente B2B (Aerolínea/Agencia) & Consume APIs, dashboards, reporta bugs & Externo \\
\midrule
Proveedor de Datos Aeronáuticos & Suministra datos crudos de vuelos, clima & Externo \\
\midrule
Auditor de Compliance & Verifica SOC 2, ISO 27001, GDPR, IATA & Externo \\
\bottomrule
\end{tabularx}
\caption{Actores del ecosistema SkyAnalytics por nivel organizacional}
\end{table}

\subsection{Objetivos Estratégicos, Tácticos y Operativos}

\subsubsection{Objetivos Estratégicos (10)}

Los Objetivos Estratégicos definen la dirección a largo plazo de SkyAnalytics. Cada OE se alinea con una o más perspectivas del BSC y se descompone en objetivos tácticos y operativos para su ejecución en cascada.

\begin{enumerate}[leftmargin=1.5cm, label=\textbf{OE\arabic*}, itemsep=2pt, topsep=4pt]
    \item \textbf{Penetración Digital (Growth Hacking):} capturar clientes internacionales mediante automatización de marketing en HubSpot, posicionamiento SEO, contenido técnico de alto valor y construcción de comunidad de desarrolladores. La estrategia se apoya en un funnel de adquisición que combina tráfico orgánico, campañas segmentadas por industria aeronáutica y un programa de referidos. El objetivo es alcanzar 50 Beta Testers activos en Q4 2026 y crecer a 500+ clientes para 2028.

    \item \textbf{Escalabilidad vía Ecosistemas (APIs/SDKs):} adoptar Specification-Driven Development con contratos OpenAPI 3.1 para cada endpoint. Publicar SDKs auto-generados en Python, JavaScript y Java. Listar la API en marketplaces como RapidAPI y construir un Developer Portal con documentación interactiva (Mintlify), sandbox con datos sintéticos y gestión de API Keys. Esto permite que un desarrollador externo se integre en menos de 3 minutos y acelera la adopción orgánica.

    \item \textbf{Infraestructura Cloud de Alta Disponibilidad:} garantizar uptime $\geq$ 99.99\,\% mediante despliegue multi-región activo-activo sobre Kubernetes, con Terraform como Infrastructure as Code. Implementar CDNs multi-proveedor para reducir latencia global, auto-scaling horizontal basado en métricas de tráfico y Disaster Recovery con RPO $\leq$ 5 min y RTO $\leq$ 15 min. La infraestructura se versiona en Git y todo cambio pasa por PR review antes del apply.

    \item \textbf{Inteligencia de Negocio Centralizada:} consolidar todos los datos operativos y financieros en un Data Warehouse columnar (MonetDB) con pipelines ETL orquestados por dbt y validados con Great Expectations. Exponer dashboards self-service en PowerBI para clientes Enterprise, con datos actualizados diariamente y segmentación por aerolínea, ruta y período. Implementar precios dinámicos basados en patrones de consumo.

    \item \textbf{Innovación en IA Aeronáutica (I+D):} desarrollar y mantener modelos de Machine Learning para predicción de retrasos, impacto climático en rutas y optimización de trayectorias. Implementar ciclo MLOps completo con MLflow (tracking y registry), Feast (feature store), SHAP (explainability), A/B testing framework y monitoreo de drift (PSI). Extender la capacidad con técnicas de Deep Learning para OCR de documentos de compliance y Speech-to-Text para análisis de feedback cualitativo (OT21, OT22).

    \item \textbf{Gobernanza y Compliance:} implementar un marco DevSecOps con Zero Trust Architecture, cifrado end-to-end (TLS 1.3 en tránsito, AES-256 en reposo), escaneo SAST/DAST en cada PR y gestión de secretos con HashiCorp Vault. Obtener certificación SOC 2 Tipo II en el año 2 e ISO 27001 en el año 3. Cumplir con GDPR mediante residencia de datos por región y con los estándares IATA para el manejo de datos aeronáuticos.

    \item \textbf{Retención mediante Customer Success:} implementar un programa de Customer Success que incluye onboarding guiado con sandbox instantáneo (TTFV $\leq$ 15 min), health scoring automatizado por tenant, alertas tempranas de riesgo de churn, chatbot IA de nivel 1 con respaldo de knowledge base y portal de autoservicio. Medir NPS trimestralmente con meta $>$ 50 en segmento Enterprise y ejecutar ciclos de feedback para mejora continua.

    \item \textbf{Consolidación de Talento Remote-First:} construir una cultura de trabajo asíncrono con ceremonias ágiles optimizadas para equipos distribuidos globalmente. Documentar todo proceso mediante ADRs (Architecture Decision Records), wikis internas y runbooks de guardia. Implementar career track IC (Individual Contributor) con revisiones de desempeño 360° semestrales. Meta: retención anual $>$ 90\,\% y eNPS $>$ 40.

    \item \textbf{Calidad y Observabilidad de Datos:} establecer data contracts formales entre productores y consumidores de datos, monitorear frescura ($<$ 5 min), volumen, schema drift y linaje mediante dashboards de observabilidad. Validar cada ingesta con suites de Great Expectations y alertar proactivamente ante degradaciones. Implementar monitoreo de extremo a extremo con trazas distribuidas (Tempo) y correlación de logs (Loki).

    \item \textbf{Gestión de Producto Data-Driven:} mantener un roadmap público y transparente para clientes y comunidad. Priorizar iniciativas con metodología RICE (Reach, Impact, Confidence, Effort). Realizar investigación continua de usuarios mediante entrevistas, encuestas y análisis de comportamiento en la plataforma. Validar pricing con A/B testing y entrevistas de willingness-to-pay antes de cada ajuste de precios.
\end{enumerate}

\subsubsection{Objetivos Tácticos (22)}

Los Objetivos Tácticos traducen la estrategia en iniciativas concretas por área funcional a mediano plazo (6--18 meses). Cada OT se vincula al menos a un OE y se descompone en objetivos operativos para su ejecución diaria.

\begin{enumerate}[leftmargin=1.5cm, label=\textbf{OT\arabic*}, itemsep=2pt, topsep=4pt]
    \item \textbf{Automatizar flujos HubSpot para Beta Testers.} Configurar secuencias de email automation, formularios de registro segmentados por industria (aerolínea, agencia, logística) y scoring de leads en HubSpot. \textit{Responsable: VP Marketing.} \textit{Plazo: Q3 2026.}

    \item \textbf{Gestionar programa Early Adopters con registro simplificado y Freemium.} Diseñar flujo de registro self-service con email verification, autoprovisionamiento de tenant y plan Freemium con 1.000 API calls/mes. \textit{Responsable: VP Marketing + Developer Advocate.} \textit{Plazo: Q3--Q4 2026.}

    \item \textbf{Estandarizar contratos OpenAPI 3.1 y GraphQL.} Definir y mantener especificaciones OpenAPI 3.1 para todos los endpoints REST, con schemas GraphQL complementarios. Validar automáticamente en CI/CD con Spectral linter antes de cada merge. \textit{Responsable: CTO.} \textit{Plazo: Q3 2026.}

    \item \textbf{Publicar API Beta en RapidAPI y Developer Portal.} Configurar listing en RapidAPI con planes Freemium/Pro/Enterprise, OAuth2 client\_credentials y rate limits por plan. Publicar documentación interactiva auto-generada en el Developer Portal (Mintlify). \textit{Responsable: CTO + Developer Advocate.} \textit{Plazo: Q3 2026.}

    \item \textbf{Desplegar infraestructura con Terraform (IaC) y CI/CD GitHub Actions.} Versionar todos los recursos cloud como módulos Terraform. Implementar pipeline CI/CD con GitHub Actions para plan/apply automatizado tras PR review. Separar entornos dev/staging/prod en cuentas cloud diferentes. \textit{Responsable: SRE Lead.} \textit{Plazo: Q3 2026.}

    \item \textbf{Configurar CDNs multi-proveedor.} Desplegar CDNs en al menos dos proveedores (CloudFront + Cloudflare) con failover automático. Configurar caché de respuestas de API con TTL por endpoint y purga programada. \textit{Responsable: SRE Lead.} \textit{Plazo: Q4 2026.}

    \item \textbf{Construir pipelines ETL con Great Expectations y dbt.} Diseñar flujos de ingesta incremental desde SFTP de proveedores, transformación con dbt y validación con Great Expectations. Orquestar con Airflow. Monitorear duración, volumen y estado de cada DAG. \textit{Responsable: Data Engineer Lead.} \textit{Plazo: Q3--Q4 2026.}

    \item \textbf{Implementar dashboards PowerBI y self-service para clientes.} Construir dashboards ejecutivos internos (BSC, financieros, operativos) y dashboards embebidos para clientes Enterprise con segmentación por aerolínea, ruta y período. Programar refreshes nocturnos. \textit{Responsable: Data Engineer Lead.} \textit{Plazo: Q4 2026.}

    \item \textbf{Gestionar ciclo de vida de modelos ML: MLflow, Feast, SHAP, A/B testing.} Implementar MLOps pipeline con MLflow para tracking de experimentos y model registry, Feast como feature store centralizada, SHAP para explainability y framework de A/B testing en producción. \textit{Responsable: ML Engineer Lead.} \textit{Plazo: Q4 2026--Q2 2027.}

    \item \textbf{Implementar seguridad perimetral, E2E encryption, SAST/DAST en CI/CD.} Configurar WAF, DDoS protection, TLS 1.3 en todas las comunicaciones, cifrado AES-256 en reposo. Integrar SonarQube (SAST) y OWASP ZAP (DAST) en el pipeline de CI/CD. Bloquear merge ante vulnerabilidades críticas. \textit{Responsable: SRE Lead.} \textit{Plazo: Q4 2026--Q1 2027.}

    \item \textbf{Mantener Developer Portal: documentación interactiva, Sandbox, SDKs, bugs.} Construir y mantener portal con Mintlify, sandbox con datos sintéticos renovados semanalmente, SDKs auto-generados (Python, JS, Java) desde specs OpenAPI y sistema de reporte de bugs con triaging NLP. \textit{Responsable: Developer Advocate.} \textit{Plazo: Q4 2026--Q2 2027.}

    \item \textbf{Gestionar trabajo en Linear con sprints de 2 semanas.} Configurar workspace por equipo, backlog priorizado con RICE, ceremonias ágiles asíncronas (daily standup en Slack, planning y retro en videollamada). Medir velocity y cycle time por sprint. \textit{Responsable: CTO + todos los leads.} \textit{Plazo: continuo desde Q3 2026.}

    \item \textbf{Mantener documentación interna: ADRs, wiki, runbooks.} Documentar toda decisión arquitectónica en ADRs versionados en Git. Mantener wiki con onboarding guides, runbooks de incidentes por tipo de alerta, y playbooks de operaciones. \textit{Responsable: CTO.} \textit{Plazo: continuo desde Q3 2026.}

    \item \textbf{Monitorear y optimizar costos cloud.} Configurar dashboards de costos por servicio y equipo, alertas de anomalías de gasto diario, recomendaciones automáticas de rightsizing y detección de recursos ociosos. Revisión semanal con CTO y SRE Lead. \textit{Responsable: SRE Lead.} \textit{Plazo: continuo desde Q4 2026.}

    \item \textbf{Implementar stack de observabilidad: Loki, Prometheus, Tempo, PagerDuty.} Desplegar Grafana LGTM stack (Loki para logs, Prometheus para métricas, Tempo para trazas) con dashboards por servicio y alertas configuradas en PagerDuty. Umbrales: CPU $>$ 80\%, memoria $>$ 85\%, tasa errores 5xx $>$ 1\%. \textit{Responsable: SRE Lead.} \textit{Plazo: Q4 2026.}

    \item \textbf{Ejecutar pruebas de carga (k6) y chaos engineering.} Diseñar escenarios de carga con 10K usuarios concurrentes simulados en k6. Ejecutar chaos experiments mensuales en staging (kill pod aleatorio, network partition, DB failover) con Litmus Chaos. Documentar resultados y MTTR. \textit{Responsable: SRE Lead.} \textit{Plazo: Q1--Q2 2027.}

    \item \textbf{Establecer fundación legal: ToS, Privacy Policy, DPA, ruta SOC 2.} Redactar y publicar Términos de Servicio, Política de Privacidad y Data Processing Agreement (DPA) con firma de abogados especializados en SaaS. Diseñar roadmap de certificación SOC 2 Tipo II para auditoría en año 2. \textit{Responsable: CEO + CTO.} \textit{Plazo: Q4 2026.}

    \item \textbf{Validar pricing: entrevistas, análisis competitivo, A/B testing.} Realizar entrevistas de willingness-to-pay con 15+ prospectos del sector aeronáutico. Analizar precios de competidores (FlightAware, OAG, Cirium). Ejecutar A/B test en landing page con 2--3 variantes de pricing. \textit{Responsable: VP Marketing + CEO.} \textit{Plazo: Q1--Q2 2027.}

    \item \textbf{Ejecutar ciclo de feedback continuo: entrevistas, NPS, sesiones, RICE.} Implementar encuestas NPS trimestrales automatizadas, entrevistas cualitativas mensuales con 5 clientes, sesiones de co-creación de features y priorización de backlog con metodología RICE. \textit{Responsable: VP Customer Success + VP Marketing.} \textit{Plazo: continuo desde Q1 2027.}

    \item \textbf{Publicar 3 artículos técnicos y participar en 2 conferencias en fase Beta.} Redactar artículos técnicos sobre arquitectura de la plataforma, modelos ML para aviación y lecciones de DevSecOps. Someter propuestas de charlas a conferencias de aviación y tecnología (IATA, Airline Tech Summit, KubeCon). \textit{Responsable: CTO + Developer Advocate.} \textit{Plazo: Q1--Q2 2027.}

    \item \textbf{Explorar técnicas de Deep Learning para compliance y soporte.} Evaluar modelos de OCR con redes neuronales para procesamiento automatizado de documentos de auditoría (facturas de proveedores, certificados, logs escaneados). Experimentar con NLP avanzado (transformers) para clasificación de tickets y análisis de sentimiento en feedback de clientes. \textit{Responsable: ML Engineer Lead.} \textit{Plazo: Q3--Q4 2027.}

    \item \textbf{Implementar Speech-to-Text para transcripción de feedback cualitativo.} Desplegar pipeline de transcripción automática con Whisper o DeepSpeech para entrevistas de feedback de clientes. Integrar con sistema de análisis cualitativo para extraer temas recurrentes y sentimiento. Medir Word Error Rate (WER) objetivo $\leq$ 10\%. \textit{Responsable: ML Engineer Lead.} \textit{Plazo: Q1--Q3 2028.}
\end{enumerate}

\subsubsection{Objetivos Operativos (34)}

Los Objetivos Operativos ejecutan la acción diaria, semanal y mensual que sostiene la operación del ecosistema SkyAnalytics. Cada OP describe una tarea concreta con su herramienta, cadencia y KPI de éxito, y se mapea a al menos un Caso de Uso operativo para su trazabilidad.

\begin{enumerate}[leftmargin=1.5cm, label=\textbf{OP\arabic*}, itemsep=2pt, topsep=4pt]
    \item \textbf{Gestionar registro automatizado de usuarios.} Monitorear el pipeline de registro self-service: verificar email confirmations, detectar cuentas duplicadas o fraudulentas, validar que el aprovisionamiento de tenant y API Keys se complete en $\leq$ 30 segundos. \textit{Herramienta: HubSpot + API Gateway.} \textit{Cadencia: diario.} \textit{KPI: $\geq$ 99\,\% registros exitosos.} \textit{CU: CU-O01.}

    \item \textbf{Monitorear rate limiting y auto-scaling de la API.} Supervisar thresholds de rate limit por tenant y plan, configurar políticas de throttling graduado (warning $\rightarrow$ soft limit $\rightarrow$ hard limit), verificar que el auto-scaling de pods responda en $\leq$ 2 minutos ante picos de tráfico. \textit{Herramienta: K8s HPA + API Gateway + Prometheus.} \textit{Cadencia: continuo.} \textit{KPI: $\leq$ 1 incidente de rate limit breach por semana.} \textit{CU: CU-O02.}

    \item \textbf{Supervisar ejecución de pipelines ETL diarios.} Monitorear duración, volumen ingerido, row count y estado de cada DAG de Airflow. En caso de fallo: investigar logs, corregir y re-ejecutar. Mantener un scorecard semanal de salud de pipelines. \textit{Herramienta: Airflow + dbt + Great Expectations.} \textit{Cadencia: diario.} \textit{KPI: 100\,\% pipelines completados.} \textit{CU: CU-O03.}

    \item \textbf{Validar calidad de datos post-ingesta.} Ejecutar suite de Great Expectations sobre datos en staging. Verificar: nulos en flight\_id y departure\_time, duplicados, rangos plausibles (retraso $\in [0, 1440]$ min), frescura $\leq$ 5 min y conformidad de schema contra data contract. Escalar hallazgos al Data Engineer Lead. \textit{Herramienta: Great Expectations + Slack.} \textit{Cadencia: tras cada ingesta.} \textit{KPI: $\geq$ 95\,\% datos con frescura $\leq$ 5 min.} \textit{CU: CU-O04.}

    \item \textbf{Reentrenar modelos ML programados.} Ejecutar pipeline semanal de reentrenamiento con features desde Feast. Entrenar XGBoost/LightGBM para retrasos y demanda. Loggear métricas en MLflow, comparar MAPE contra baseline. Auto-promover a staging si mejora $>$ 2\,\%. \textit{Herramienta: MLflow + Feast + XGBoost.} \textit{Cadencia: semanal.} \textit{KPI: MAPE $\leq$ 15\,\%.} \textit{CU: CU-O05.}

    \item \textbf{Verificar drift de modelos en producción.} Monitorear PSI (Population Stability Index) y data drift diariamente sobre modelos productivos. Alertar al ML Engineer si PSI $>$ 0.2. Disparar reentrenamiento automático si se supera el umbral. Documentar causas de drift en runbook. \textit{Herramienta: Evidently AI + MLflow + PagerDuty.} \textit{Cadencia: diario.} \textit{KPI: modelos con PSI $\leq$ 0.2 en 95\,\% del tiempo.} \textit{CU: CU-O05.}

    \item \textbf{Refrescar dashboards BI para clientes.} Ejecutar refresh programado nocturno de todos los dashboards PowerBI. Validar que los KPIs carguen correctamente (verificación visual de widgets clave), verificar consistencia numérica contra fuente, loggear duración total del refresh. \textit{Herramienta: PowerBI + dbt.} \textit{Cadencia: nocturno diario.} \textit{KPI: dashboards actualizados antes de las 06:00 UTC.} \textit{CU: CU-O06.}

    \item \textbf{Monitorear telemetría de contenedores en tiempo real.} Recolectar métricas de CPU, memoria, red y tasa de errores por pod cada 30 segundos. Visualizar en dashboards de Grafana por servicio. Responder a alertas de PagerDuty según severidad en el runbook correspondiente. \textit{Herramienta: Prometheus + Grafana + PagerDuty.} \textit{Cadencia: continuo.} \textit{KPI: detección de anomalías $\leq$ 2 min.} \textit{CU: CU-O07.}

    \item \textbf{Configurar y mantener reglas de alerta.} Definir thresholds por servicio (CPU $>$ 80\,\%, memoria $>$ 85\,\%, errores 5xx $>$ 1\,\%, latencia p95 $>$ 500ms), probar disparo de cada alerta en staging, mantener runbooks actualizados, revisar y eliminar falsos positivos semanalmente. \textit{Herramienta: Prometheus AlertManager + PagerDuty.} \textit{Cadencia: semanal (revisión).} \textit{KPI: $\leq$ 5\,\% falsos positivos.} \textit{CU: CU-O07.}

    \item \textbf{Ejecutar escaneo SAST/DAST en pipeline de CI/CD.} Integrar SonarQube (SAST) y OWASP ZAP (DAST) en GitHub Actions. Bloquear merge ante vulnerabilidades con severidad critical o high. Generar reporte de seguridad por release. Mantener baseline de falsos positivos aceptados documentado. \textit{Herramienta: GitHub Actions + SonarQube + OWASP ZAP.} \textit{Cadencia: por cada PR.} \textit{KPI: 0 vulnerabilidades críticas en producción.} \textit{CU: CU-O18.}

    \item \textbf{Aplicar parches de seguridad críticos.} Monitorear CVEs de dependencias del stack con Dependabot y Snyk. Parchear vulnerabilidades críticas en $\leq$ 24 horas, altas en $\leq$ 72 horas. Actualizar inventory de versiones de paquetes en la wiki. \textit{Herramienta: Dependabot + Snyk + GitHub Security Alerts.} \textit{Cadencia: continuo.} \textit{KPI: tiempo medio de parche crítico $\leq$ 12 horas.} \textit{CU: CU-O18.}

    \item \textbf{Atender tickets de desarrolladores externos.} Triaging automático con NLP para clasificar severidad, categoría y prioridad del ticket. Responder primera respuesta en $\leq$ 2 horas. Investigar y resolver en $\leq$ 24 horas, o escalar a engineering con contexto completo. Notificar resolución al reportante. \textit{Herramienta: Linear + chatbot IA + Slack.} \textit{Cadencia: continuo.} \textit{KPI: tiempo primera respuesta $\leq$ 2 horas.} \textit{CU: CU-O08.}

    \item \textbf{Documentar FAQs en base de conocimientos del chatbot.} Detectar temas recurrentes cuando se acumulan $\geq$ 3 tickets similares en 7 días. Generar borrador automático de FAQ con NLP. Revisar, editar, añadir screenshots y publicar. Actualizar embeddings del chatbot. Medir tasa de deflexión resultante. \textit{Herramienta: NLP embeddings + KB portal + chatbot IA.} \textit{Cadencia: semanal.} \textit{KPI: tasa de deflexión $\geq$ 30\,\%.} \textit{CU: CU-O14.}

    \item \textbf{Gestionar sprints en Linear.} Crear tickets desde el roadmap y el backlog priorizado. Actualizar estado diario en Slack standup. Cerrar sprint con retrospectiva documentada. Medir velocity, cycle time y burndown. Ajustar capacidad del próximo sprint según datos históricos. \textit{Herramienta: Linear + Slack bot.} \textit{Cadencia: daily + cierre de sprint (2 semanas).} \textit{KPI: $\geq$ 85\,\% tickets completados por sprint.} \textit{CU: CU-O15.}

    \item \textbf{Revisar y optimizar costos cloud semanalmente.} Auditar factura cloud desglosada por servicio. Comparar gasto real contra forecast mensual. Detectar anomalías de gasto diario ($>$ 20\,\% desviación). Identificar recursos ociosos o sobre-aprovisionados. Generar recomendaciones de ahorro y rightsizing. \textit{Herramienta: AWS Cost Explorer + Grafana.} \textit{Cadencia: semanal.} \textit{KPI: desviación $\leq$ 5\,\% vs presupuesto mensual.} \textit{CU: CU-O19.}

    \item \textbf{Ejecutar backup diario automatizado.} Programar pg\_dump de base de datos transaccional a las 03:00 UTC. Encriptar backup con KMS (AES-256). Subir a cold storage (S3 Glacier Deep Archive). Loggear hash SHA-256 de integridad. Verificar tamaño dentro del rango esperado. \textit{Herramienta: pg\_dump + AWS KMS + S3 Glacier.} \textit{Cadencia: diario.} \textit{KPI: backup exitoso 100\,\% de los días.} \textit{CU: CU-O09.}

    \item \textbf{Ejecutar prueba de restauración semanal.} Restaurar último backup en instancia de prueba aislada. Validar integridad de datos (row counts por tabla, checksums, foreign key consistency). Ejecutar smoke tests de API sobre la instancia restaurada. Loggear resultado y alertar si falla. \textit{Herramienta: pg\_restore + script de validación + k6 smoke test.} \textit{Cadencia: semanal.} \textit{KPI: restore exitoso 100\,\% de las semanas.} \textit{CU: CU-O09.}

    \item \textbf{Calcular SLI/SLO y error budget de uptime.} Computar SLI de uptime real de los últimos 30 días desde métricas de Prometheus. Comparar contra SLO de 99.99\,\%. Proyectar consumo de error budget a fin de mes (target: 4.38 min/mes). Recomendar congelación de deploys si consumo $>$ 80\,\%. \textit{Herramienta: Prometheus + script SLI/SLO.} \textit{Cadencia: semanal.} \textit{KPI: uptime real $\geq$ 99.99\,\%.} \textit{CU: CU-O10.}

    \item \textbf{Ejecutar congelación de deploys si error budget $>$ 80\,\%.} Si el umbral se supera, bloquear pipeline de CI/CD automáticamente vía branch protection en GitHub. Notificar a CTO, tech leads y SRE team por Slack y PagerDuty con la justificación y el plan de remediación. Documentar la congelación en el runbook de incidentes. \textit{Herramienta: GitHub Actions + PagerDuty + Slack.} \textit{Cadencia: bajo demanda.} \textit{KPI: tiempo detección $\rightarrow$ congelación $\leq$ 5 min.} \textit{CU: CU-O10.}

    \item \textbf{Publicar changelog semanal con Semantic Versioning.} Recolectar commits desde el último release tag. Categorizar automáticamente por tipo (feat, fix, breaking, docs, chore). Bump version según reglas de SemVer. Publicar en Developer Portal. Enviar email a suscriptores y publicar en canal \#changelog de Slack. \textit{Herramienta: git log + release-it + DevPortal CMS.} \textit{Cadencia: semanal (viernes).} \textit{KPI: changelog publicado 100\,\% de los viernes.} \textit{CU: CU-O11.}

    \item \textbf{Rotar credenciales de servicio mensualmente.} Identificar secrets próximos a expirar (API Keys, DB passwords, TLS certs) desde Vault. Generar nuevas credenciales. Actualizar K8s secrets y reiniciar pods. Verificar que todos los servicios funcionen con las nuevas credenciales vía smoke tests. Desactivar credenciales antiguas tras período de gracia de 24 horas. \textit{Herramienta: HashiCorp Vault + K8s + script smoke test.} \textit{Cadencia: mensual.} \textit{KPI: 0 incidentes por credencial expirada.} \textit{CU: CU-O12.}

    \item \textbf{Auditar accesos IAM trimestralmente.} Revisar todos los roles, políticas y accesos de servicio en la cuenta cloud. Identificar permisos sobre-asignados (principle of least privilege). Revocar accesos de ex-empleados y servicios deprecados. Generar reporte de auditoría con hallazgos y acciones tomadas. \textit{Herramienta: AWS IAM Access Analyzer + reporte manual.} \textit{Cadencia: trimestral.} \textit{KPI: 0 accesos no justificados al cierre de auditoría.} \textit{CU: CU-O12.}

    \item \textbf{Realizar post-mortem blameless $\leq$ 72 horas post-incidente.} Generar timeline automático desde logs y métricas del incidente. Facilitar reunión blameless con 5 whys, documentar aprendizajes y acciones preventivas. Publicar reporte en wiki y crear tickets en Linear. \textit{Herramienta: PagerDuty + Linear + wiki.} \textit{Cadencia: $\leq$ 72 horas.} \textit{KPI: 100\,\% incidentes Sev1/Sev2 con post-mortem.} \textit{CU: CU-O13.}

    \item \textbf{Validar especificaciones OpenAPI 3.1 en CI/CD.} Ejecutar linter Spectral sobre cada spec en cada PR. Verificar que no haya breaking changes sin bump de versión major. Validar ejemplos de request/response incluidos. Asegurar que todos los endpoints tengan descripción y tags. \textit{Herramienta: Spectral + OpenAPI Generator + GitHub Actions.} \textit{Cadencia: por cada PR.} \textit{KPI: 0 errores de lint en specs mergeadas.} \textit{CU: CU-O16.}

    \item \textbf{Monitorear presencia en marketplace RapidAPI.} Verificar uptime y disponibilidad del listing. Responder a reviews y preguntas de desarrolladores en $\leq$ 24 horas. Actualizar descripciones de planes trimestralmente. Monitorear métricas de adopción: page views, subscriptions por plan, churn del marketplace. \textit{Herramienta: RapidAPI Dashboard.} \textit{Cadencia: continuo.} \textit{KPI: rating $\geq$ 4.5 estrellas en marketplace.} \textit{CU: CU-T02.}

    \item \textbf{Mantener documentación interactiva del Developer Portal.} Sincronizar documentación con specs OpenAPI tras cada release. Verificar que todos los endpoints tengan ejemplos request/response funcionales. Actualizar guías de migración ante breaking changes. Revisar analíticas de páginas para identificar endpoints con documentación deficiente. \textit{Herramienta: Mintlify + CI pipeline.} \textit{Cadencia: tras cada release.} \textit{KPI: docs sincronizadas en $\leq$ 24 horas post-release.} \textit{CU: CU-T07.}

    \item \textbf{Ejecutar pruebas de carga semanales en staging.} Simular tráfico de 10K usuarios concurrentes con k6. Medir: latencia p95 y p99, throughput (req/s), tasa de errores y uso de recursos durante la prueba. Comparar contra baseline. Documentar resultados en reporte con gráficos. \textit{Herramienta: k6 + Grafana k6 Cloud.} \textit{Cadencia: semanal (lunes staging).} \textit{KPI: latencia p95 $\leq$ 500ms bajo 10K usuarios concurrentes.} \textit{CU: CU-T09.}

    \item \textbf{Ejecutar chaos engineering mensualmente.} Inyectar fallos controlados en staging: kill de pod aleatorio, network partition entre servicios, DB primary failover, latencia inyectada de 500ms. Medir MTTR de cada experimento. Documentar hallazgos, tiempo de recuperación y recomendaciones en reporte. \textit{Herramienta: Litmus Chaos + k6.} \textit{Cadencia: mensual (primer viernes del mes).} \textit{KPI: MTTR $\leq$ RTO de 15 minutos en todos los escenarios.} \textit{CU: CU-T09.}

    \item \textbf{Mantener documentación legal y compliance actualizada.} Revisar y actualizar ToS, Privacy Policy y DPA ante cambios regulatorios o de producto. Publicar registros de procesamiento de datos por región (GDPR Art. 30). Mantener evidencias de controles SOC 2 organizadas por trust service criteria. \textit{Herramienta: compliance platform + wiki + auditor externo.} \textit{Cadencia: trimestral.} \textit{KPI: 100\,\% documentos legales vigentes y actualizados.} \textit{CU: CU-E04.}

    \item \textbf{Recolectar y analizar feedback de clientes.} Enviar encuestas NPS trimestrales automatizadas a todos los usuarios activos (Typeform). Recolectar CSAT post-interacción con soporte. Transcribir entrevistas cualitativas con Speech-to-Text. Clasificar comentarios por tema y sentimiento para alimentar backlog de producto. \textit{Herramienta: Typeform + HubSpot + STT engine.} \textit{Cadencia: trimestral (NPS) + continuo (CSAT).} \textit{KPI: tasa de respuesta NPS $\geq$ 30\,\%.} \textit{CU: CU-T10.}

    \item \textbf{Publicar contenido técnico regularmente.} Redactar y publicar artículos técnicos en el blog corporativo. Compartir en Dev.to, HackerNews y LinkedIn. Preparar y someter propuestas de charlas a conferencias (KubeCon, IATA Summit, AI in Aviation). Medir engagement y conversiones a registro. \textit{Herramienta: blog CMS + Buffer + Google Analytics.} \textit{Cadencia: semanal (artículo) + trimestral (conferencia).} \textit{KPI: $\geq$ 1 artículo técnico por semana, $\geq$ 2 charlas por año.} \textit{CU: CU-T01.}

    \item \textbf{Ejecutar experimentos de Deep Learning para OCR.} Explorar modelos de OCR con redes neuronales (CRNN + CTC, TrOCR) para procesamiento automatizado de: facturas de proveedores de datos, certificados de cumplimiento, logs de auditoría escaneados. Comparar precisión contra baseline de entrada manual. Registrar experimentos en MLflow. \textit{Herramienta: PyTorch / TensorFlow + MLflow.} \textit{Cadencia: sprint de experimentación (4 semanas).} \textit{KPI: precisión OCR $\geq$ 95\,\% en documentos de compliance.} \textit{CU: CU-O20.}

    \item \textbf{Implementar y mantener pipeline de Speech-to-Text.} Desplegar modelo STT (Whisper / DeepSpeech). Integrar con sistema de recolección de feedback para transcripción automática de entrevistas. Medir Word Error Rate (WER). Reentrenar con vocabulario de dominio aeronáutico para mejorar precisión. Extraer temas y sentimiento del texto transcrito. \textit{Herramienta: Whisper / DeepSpeech + MLflow.} \textit{Cadencia: continuo.} \textit{KPI: WER $\leq$ 10\,\%.} \textit{CU: CU-O20.}

    \item \textbf{Validar data contracts entre productores y consumidores.} Definir y versionar en Git los schema contracts para cada fuente de datos externa. Validar en CI que los datos entrantes cumplan el contrato (tipos, nulabilidad, rangos). Alertar al Data Engineer Lead ante schema drift del proveedor. Mantener changelog de cambios de contrato. \textit{Herramienta: Great Expectations + dbt contracts + Git.} \textit{Cadencia: por cada ingesta.} \textit{KPI: 0 ingestion failures por schema mismatch.} \textit{CU: CU-O17.}
\end{enumerate}

\medskip

\textbf{Trazabilidad OE $\rightarrow$ Perspectiva BSC $\rightarrow$ KPI $\rightarrow$ Meta}

\begin{table}[H]
\small
\centering
\begin{tabularx}{\linewidth}{>{\raggedright\arraybackslash}p{0.8cm} >{\raggedright\arraybackslash}p{3cm} >{\raggedright\arraybackslash}p{2.8cm} >{\raggedright\arraybackslash}p{2.5cm} >{\raggedright\arraybackslash}p{2.5cm} >{\raggedright\arraybackslash}X}
\toprule
\multicolumn{1}{>{\raggedright\arraybackslash}p{0.8cm}}{\textbf{OE}} &
\multicolumn{1}{>{\raggedright\arraybackslash}p{3cm}}{\textbf{Perspectiva BSC}} &
\multicolumn{1}{>{\raggedright\arraybackslash}p{2.8cm}}{\textbf{KPI Principal}} &
\multicolumn{1}{>{\raggedright\arraybackslash}p{2.5cm}}{\textbf{Meta 2026}} &
\multicolumn{1}{>{\raggedright\arraybackslash}p{2.5cm}}{\textbf{Meta 2028}} &
\multicolumn{1}{>{\raggedright\arraybackslash}X}{\textbf{OT Asociados}} \\
\midrule
OE1 & Financiera + Cliente & Beta Testers activos & 50 & $\ge$ 500 & OT1, OT2, OT4, OT19 \\
\midrule
OE2 & Cliente + Procesos & TTFV (min) & $\le$ 3 & $\le$ 1 & OT3, OT4, OT11 \\
\midrule
OE3 & Procesos & Uptime (\%) & 99.9\% & $\ge$ 99.99\% & OT5, OT6, OT14, OT15, OT16 \\
\midrule
OE4 & Financiera + Procesos & Data Freshness & $\ge$ 95\% & $\ge$ 98\% & OT7, OT8, OT15 \\
\midrule
OE5 & Aprendizaje + Procesos & MAPE retrasos & $\le$ 15\% & $\le$ 8\% & OT9, OT13, OT21, OT22 \\
\midrule
OE6 & Procesos & Controles cumplidos (\%) & 70\% (SOC 2) & 100\% (SOC 2 + ISO) & OT10, OT17 \\
\midrule
OE7 & Cliente & NPS Enterprise & --- & $>$ 50 & OT2, OT11, OT19 \\
\midrule
OE8 & Aprendizaje & eNPS interno & $>$ 40 & $>$ 50 & OT12, OT13, OT20 \\
\midrule
OE9 & Procesos & Data Contracts válidos (\%) & 90\% & 100\% & OT7, OT15, OT18 \\
\midrule
OE10 & Cliente + Financiera & Initiatives on track (\%) & 80\% & $\ge$ 90\% & OT18, OT19, OT20 \\
\bottomrule
\end{tabularx}
\caption{Trazabilidad de cada Objetivo Estratégico con su perspectiva BSC, KPI principal, metas a corto y largo plazo, y objetivos tácticos asociados}
\end{table}

\subsection{Catálogo General de Casos de Uso}

\small
\setlength{\LTleft}{0pt}
\setlength{\LTright}{0pt}
\begin{longtable}{>{\raggedright\arraybackslash}p{1.6cm} >{\raggedright\arraybackslash}p{6.2cm} >{\raggedright\arraybackslash}p{1cm} >{\raggedright\arraybackslash}p{1.8cm} >{\raggedright\arraybackslash}p{2.7cm}}
\caption{Catálogo de 36 Casos de Uso del Ecosistema SkyAnalytics}\\
\toprule
\textbf{Código} & \textbf{Nombre} & \textbf{Prior.} & \textbf{Nivel} & \textbf{Actor Principal} \\
\midrule
\endfirsthead
\multicolumn{5}{c}{\small \textit{...continuación...}}\\
\toprule
\textbf{Código} & \textbf{Nombre} & \textbf{Prior.} & \textbf{Nivel} & \textbf{Actor Principal} \\
\midrule
\endhead
\midrule
\endfoot
\bottomrule
\endlastfoot
\multicolumn{5}{>{\raggedright\arraybackslash}p{\linewidth}}{\textbf{Estratégicos (6)}} \\
\midrule
CU-E01 & Consultar Tablero BSC & 10 & Estratégico & CEO \\
CU-E02 & Analizar ARR y Rentabilidad & 9 & Estratégico & CEO \\
CU-E03 & Evaluar Uptime Global y SLA & 10 & Estratégico & CTO \\
CU-E04 & Revisar Cumplimiento Normativo & 8 & Estratégico & CTO / Auditor \\
CU-E05 & Definir Metas Estratégicas Trimestrales & 9 & Estratégico & CEO / CTO \\
CU-E06 & Analizar Retención de Talento y eNPS & 7 & Estratégico & CEO / VP People \\
\midrule
\multicolumn{5}{>{\raggedright\arraybackslash}p{\linewidth}}{\textbf{Tácticos (10)}} \\
\midrule
CU-T01 & Gestionar Campañas de Growth Hacking & 8 & Táctico & VP Marketing \\
CU-T02 & Configurar y Publicar API en RapidAPI & 9 & Táctico & CTO / Developer Advocate \\
CU-T03 & Gestionar Infraestructura con Terraform & 10 & Táctico & SRE / DevOps \\
CU-T04 & Monitorear Pipelines ETL y Calidad & 9 & Táctico & Data Engineer \\
CU-T05 & Entrenar y Versionar Modelos ML & 9 & Táctico & ML Engineer \\
CU-T06 & Configurar Alertas de Seguridad & 9 & Táctico & SRE / DevOps \\
CU-T07 & Mantener Developer Portal y SDKs & 8 & Táctico & Developer Advocate \\
CU-T08 & Analizar Costos Cloud y Optimizar & 8 & Táctico & CTO / SRE \\
CU-T09 & Ejecutar Pruebas de Carga y Chaos & 8 & Táctico & SRE / DevOps \\
CU-T10 & Validar Estrategia de Pricing & 8 & Táctico & VP Marketing / CEO \\
\midrule
\multicolumn{5}{>{\raggedright\arraybackslash}p{\linewidth}}{\textbf{Operativos (20)}} \\
\midrule
CU-O01 & Registrar Tenant y Autogenerar API Keys & 10 & Operativo & Cliente B2B / Sistema \\
CU-O02 & Consultar Datos de Vuelo vía API & 10 & Operativo & Cliente B2B \\
CU-O03 & Ejecutar Ingesta Diaria de Datos & 10 & Operativo & Data Engineer / Sistema \\
CU-O04 & Validar Calidad de Datos en ETL & 10 & Operativo & Data Engineer / Sistema \\
CU-O05 & Reentrenar Modelo de Retrasos & 9 & Operativo & ML Engineer \\
CU-O06 & Refrescar Dashboards BI & 9 & Operativo & Data Engineer / Sistema \\
CU-O07 & Monitorear Telemetría de Contenedores & 10 & Operativo & SRE / DevOps \\
CU-O08 & Atender Ticket de Bug & 9 & Operativo & Developer Advocate \\
CU-O09 & Ejecutar Backup y Prueba de Restauración & 10 & Operativo & SRE / DevOps \\
CU-O10 & Revisar Error Budget y Congelar Deploys & 9 & Operativo & SRE / DevOps \\
CU-O11 & Publicar Changelog Semanal & 7 & Operativo & Developer Advocate \\
CU-O12 & Rotar Credenciales y Auditar IAM & 9 & Operativo & SRE / DevOps \\
CU-O13 & Realizar Post-Mortem de Incidente & 8 & Operativo & SRE / DevOps \\
CU-O14 & Documentar FAQ en Base de Conocimientos & 7 & Operativo & Developer Advocate \\
CU-O15 & Gestionar Sprint en Linear & 8 & Operativo & Todo el equipo \\
CU-O16 & Validar Especificaciones OpenAPI en CI/CD & 9 & Operativo & CTO / SRE \\
CU-O17 & Validar Data Contracts en Pipeline ETL & 10 & Operativo & Data Engineer \\
CU-O18 & Ejecutar Pipeline Seguridad SAST/DAST en CI/CD & 9 & Operativo & SRE / DevOps \\
CU-O19 & Analizar y Optimizar Costos Cloud & 8 & Operativo & CTO / SRE \\
CU-O20 & Ejecutar Experimentos de Deep Learning & 7 & Operativo & ML Engineer \\
\bottomrule
\end{longtable}

\medskip
\textbf{Nota:} El detalle completo de cada Caso de Uso —incluyendo propósito, descripción detallada, precondiciones, flujo principal, postcondiciones y 180 historias de usuario— se documenta en el \textbf{Apéndice: Casos de Uso Detallados} al final de este documento.

%---------------------------------------------------------------------


\subsection{Módulos Principales del Sistema}

\begin{table}[H]
\small
\centering
\begin{tabularx}{\linewidth}{>{\raggedright\arraybackslash}p{3.5cm} >{\raggedright\arraybackslash}X}
\toprule
\multicolumn{1}{>{\raggedright\arraybackslash}p{3.5cm}}{\textbf{Módulo}} &
\multicolumn{1}{>{\raggedright\arraybackslash}X}{\textbf{Descripción}} \\
\midrule
API Gateway & Autenticación OAuth2, Rate Limiting, versionado APIs REST/GraphQL, logging consumo \\
\midrule
Data Ingestion \& ETL & Ingesta proveedores, validación Great Expectations, transformación dbt, carga incremental a DW. \textbf{BD transaccional:} PocketBase con escalado horizontal planificado según crecimiento de tenants. \\
\midrule
Data Warehouse \& BI & Almacenamiento columnar (MonetDB), catálogo datos, dashboards PowerBI, dashboards embebidos clientes \\
\midrule
ML Pipeline & Entrenamiento, registro (MLflow), feature store (Feast), explainability (SHAP), despliegue, monitoreo drift, A/B testing \\
\midrule
Developer Portal & Documentación interactiva (Mintlify), Sandbox, generación SDKs, gestión API Keys, changelog, reporte bugs \\
\midrule
Billing \& Subscriptions & Planes (Freemium, Pro, Enterprise), facturación por uso, pasarela pago \\
\midrule
Observability \& Alerting & Logs (Loki), métricas (Prometheus), trazas (Tempo), dashboards (Grafana), alertas PagerDuty/Slack \\
\midrule
Compliance \& Security Hub & Escaneo SAST/DAST, secretos, auditoría IAM, registros acceso, reportes cumplimiento, residencia datos \\
\bottomrule
\end{tabularx}
\caption{Módulos principales del ecosistema SkyAnalytics}
\end{table}

%---------------------------------------------------------------------
\subsubsection{Stack Tecnológico por Módulo}

\begin{table}[H]
\small
\centering
\begin{tabularx}{\linewidth}{>{\raggedright\arraybackslash}p{2.8cm} >{\raggedright\arraybackslash}p{5.2cm} >{\raggedright\arraybackslash}X}
\toprule
\multicolumn{1}{>{\raggedright\arraybackslash}p{2.8cm}}{\textbf{Módulo}} &
\multicolumn{1}{>{\raggedright\arraybackslash}p{5.2cm}}{\textbf{Tecnologías Principales}} &
\multicolumn{1}{>{\raggedright\arraybackslash}X}{\textbf{Interfaces / Dependencias}} \\
\midrule
API Gateway & Kong / AWS API GW, OAuth2, JWT, OpenAPI 3.1, Spectral linter & ML Pipeline (predicciones), Data WH (consultas), Billing (planes), Observability (logs) \\
\midrule
Data Ingestion \& ETL & Apache Airflow, dbt, Great Expectations, SFTP/S3, PocketBase & Data WH (salida), Observability (logs de DAGs), ML Pipeline (features) \\
\midrule
Data Warehouse \& BI & MonetDB (columnar), dbt, PowerBI, catálogo de datos & API Gateway (servicio datos), Compliance Hub (auditoría de acceso) \\
\midrule
ML Pipeline & MLflow, Feast (feature store), SHAP, XGBoost, LightGBM, Evidently AI & Data WH (features), API Gateway (predicciones), Observability (métricas de drift) \\
\midrule
Developer Portal & Mintlify, OpenAPI Generator, Sandbox, GitHub Actions & API Gateway (specs), Observability (estado sandbox) \\
\midrule
Billing \& Subscriptions & Stripe / Paddle, PocketBase (tenants), plan engine & API Gateway (cuotas por plan), Observability (consumo) \\
\midrule
Observability \& Alerting & Prometheus, Loki, Tempo, Grafana, PagerDuty, Slack & Todos los módulos (consumidor de métricas, logs y trazas) \\
\midrule
Compliance \& Security Hub & HashiCorp Vault, SonarQube, OWASP ZAP, Dependabot, AWS IAM & CI/CD (pipeline de seguridad), Observability (registros de acceso) \\
\bottomrule
\end{tabularx}
\caption{Stack tecnológico completo por módulo del ecosistema SkyAnalytics}
\end{table}

%---------------------------------------------------------------------
\subsubsection{Diagrama de Dependencias entre Módulos}

La siguiente tabla describe las relaciones de dependencia directa entre módulos. Una dependencia indica que el módulo origen consume servicios, datos o eventos del módulo destino para su funcionamiento.

\begin{table}[H]
\small
\centering
\begin{tabularx}{\linewidth}{>{\raggedright\arraybackslash}p{3.5cm} >{\raggedright\arraybackslash}p{3.5cm} >{\raggedright\arraybackslash}X}
\toprule
\multicolumn{1}{>{\raggedright\arraybackslash}p{3.5cm}}{\textbf{Módulo Origen}} &
\multicolumn{1}{>{\raggedright\arraybackslash}p{3.5cm}}{\textbf{Módulo Destino}} &
\multicolumn{1}{>{\raggedright\arraybackslash}X}{\textbf{Tipo de Dependencia}} \\
\midrule
API Gateway & ML Pipeline & Solicita predicciones en tiempo real (REST) \\
\midrule
API Gateway & Data WH \& BI & Consulta datos históricos de vuelos y KPIs \\
\midrule
API Gateway & Billing & Valida plan activo y cuotas por tenant \\
\midrule
Data Ingestion \& ETL & Data WH \& BI & Carga transformada incremental al DW \\
\midrule
Data Ingestion \& ETL & ML Pipeline & Provee features al feature store (Feast) \\
\midrule
ML Pipeline & Data WH \& BI & Lee features históricas y escribe predicciones \\
\midrule
Developer Portal & API Gateway & Consume specs OpenAPI para documentación \\
\midrule
Compliance Hub & CI/CD (GitHub Actions) & Integra SAST/DAST en cada Pull Request \\
\midrule
Observability & Todos los módulos & Recolecta métricas, logs y trazas (consumidor universal) \\
\midrule
Billing & API Gateway & Aplica límites de rate per plan en tiempo real \\
\bottomrule
\end{tabularx}
\caption{Dependencias directas entre módulos del ecosistema SkyAnalytics}
\end{table}

%---------------------------------------------------------------------
\subsection{Plan de Acción Estratégico}

El Plan de Acción detalla las iniciativas concretas requeridas para ejecutar la estrategia en tres fases. Cada acción incluye dependencias, responsable y un KPI de éxito medible. Las acciones se organizan por fase del roadmap: MVP (2026), Growth (2027) y Scale (2028).

\small
\setlength{\LTleft}{0pt}
\setlength{\LTright}{0pt}
\begin{longtable}{>{\raggedright\arraybackslash}p{0.5cm} >{\raggedright\arraybackslash}p{3.1cm} >{\raggedright\arraybackslash}p{1.8cm} >{\raggedright\arraybackslash}p{1cm} >{\raggedright\arraybackslash}p{2.2cm} >{\raggedright\arraybackslash}p{6.5cm}}
\caption{Plan de Acción Estratégico con 32 iniciativas distribuidas en 3 fases}\\
\toprule
\textbf{\#} & \textbf{Acción} & \textbf{Plazo} & \textbf{Dep.} & \textbf{Responsable} & \textbf{KPI de Éxito} \\
\midrule
\endfirsthead
\multicolumn{6}{c}{\small \textit{...continuación...}}\\
\toprule
\textbf{\#} & \textbf{Acción} & \textbf{Plazo} & \textbf{Dep.} & \textbf{Responsable} & \textbf{KPI de Éxito} \\
\midrule
\endhead
\midrule
\endfoot
\bottomrule
\endlastfoot
\multicolumn{6}{>{\raggedright\arraybackslash}p{\linewidth}}{\textbf{Fase 1 --- MVP (Q3--Q4 2026)}} \\
\midrule
1 & Publicar API Beta en RapidAPI & Q3 2026 & --- & CTO, Developer Advocate & Listing activo con $\ge$ 3 endpoints funcionales \\
\midrule
2 & Automatizar flujos HubSpot para captación & Q3 2026 & --- & VP Marketing & $\ge$ 200 registros calificados \\
\midrule
3 & Desplegar infra K8s con Terraform (IaC) & Q3 2026 & --- & SRE Lead & 100\,\% de recursos cloud versionados en Git \\
\midrule
4 & Construir pipeline ETL con Great Expectations + dbt & Q3--Q4 2026 & 3 & Data Engineer Lead & Datos en staging con $\le$ 5 min de lag \\
\midrule
5 & Implementar stack de observabilidad (LGTM) & Q4 2026 & 3 & SRE Lead & Dashboards operativos + alertas PagerDuty funcionales \\
\midrule
6 & Desarrollar modelo de predicción de retrasos v1 & Q4 2026 & 4 & ML Engineer Lead & MAPE $\le$ 15\,\% en datos de prueba \\
\midrule
7 & Lanzar Developer Portal + Sandbox & Q4 2026 & 1 & Developer Advocate & Desarrollador externo puede hacer primera API call en $\le$ 3 min \\
\midrule
8 & Establecer fundación legal (ToS, DPA, ruta SOC 2) & Q4 2026 & --- & CEO, CTO & ToS + Privacy + DPA publicados y revisados por abogados \\
\midrule
9 & Configurar autenticación OAuth2 + API Keys & Q4 2026 & 1, 3 & CTO, SRE & Flujo client\_credentials + API Key generación $<$ 30s \\
\midrule
10 & Implementar rate limiting por plan & Q4 2026 & 1, 3 & SRE Lead & Rate limits funcionales con headers X-RateLimit-Remaining \\
\midrule
11 & Lanzar plan Freemium con registro self-service & Q4 2026 & 1, 2, 9 & VP Marketing, Developer Advocate & Primeros 50 Beta Testers onboardeados \\
\midrule
12 & Configurar CI/CD con GitHub Actions & Q3 2026 & 3 & SRE Lead & Pipeline funcional: test $\rightarrow$ build $\rightarrow$ deploy automático \\
\midrule
\multicolumn{6}{>{\raggedright\arraybackslash}p{\linewidth}}{\textbf{Fase 2 --- Growth (Q1--Q4 2027)}} \\
\midrule
13 & Ejecutar pruebas de carga (k6) y chaos engineering & Q1 2027 & 3, 5 & SRE Lead & Latencia p95 $\le$ 500ms con 10K usuarios concurrentes \\
\midrule
14 & Validar estrategia de pricing con A/B testing & Q1--Q2 2027 & 7 & VP Marketing, CEO & 2--3 variantes testadas, pricing final definido \\
\midrule
15 & Publicar 3 artículos técnicos + 2 conferencias & Q1--Q2 2027 & 6 & CTO, Developer Advocate & $\ge$ 3 artículos en blog + $\ge$ 2 charlas aceptadas \\
\midrule
16 & Implementar MLOps pipeline completo (Feast, SHAP) & Q1--Q2 2027 & 6 & ML Engineer Lead & Feature store operativa + explainability en predicciones \\
\midrule
17 & Optimizar PocketBase para alta concurrencia & Q2 2027 & 4 & CTO, DataEng & PocketBase optimizado para 1.000+ tenants simultáneos \\
\midrule
18 & Implementar monitoring de costos cloud & Q2 2027 & 5 & SRE Lead & Dashboard de costos + alertas de anomalía \\
\midrule
19 & Lanzar SDKs auto-generados (Python, JS, Java) & Q2 2027 & 7 & Developer Advocate & 3 SDKs publicados en npm, PyPI, Maven Central \\
\midrule
20 & Implementar data contracts con validación automática & Q2--Q3 2027 & 4 & Data Engineer Lead & 100\,\% de fuentes con contrato definido y validado \\
\midrule
21 & Iniciar programa SOC 2 Tipo II (monitoreo de controles) & Q3 2027 & 8 & CEO, CTO & Evidencias de controles recolectadas por 6+ meses \\
\midrule
22 & Implementar pipeline SAST/DAST en CI/CD & Q3 2027 & 12 & SRE Lead & 0 vulnerabilidades críticas en producción \\
\midrule
23 & Lanzar dashboards self-service para clientes & Q3 2027 & 4, 17 & Data Engineer Lead & Clientes Enterprise pueden ver sus propios dashboards \\
\midrule
24 & Implementar NPS trimestral automatizado + feedback loop & Q3 2027 & 11 & VP Customer Success & $\ge$ 30\,\% tasa de respuesta en encuestas \\
\midrule
25 & Explorar técnicas de Deep Learning (OCR, NLP avanzado) & Q3--Q4 2027 & 6, 16 & ML Engineer Lead & $\ge$ 2 experimentos documentados en MLflow \\
\midrule
\multicolumn{6}{>{\raggedright\arraybackslash}p{\linewidth}}{\textbf{Fase 3 --- Scale (Q1--Q4 2028)}} \\
\midrule
26 & Desplegar infraestructura multi-región activo-activo & Q1 2028 & 13, 17 & SRE Lead & Uptime $\ge$ 99.99\,\% con failover automático \\
\midrule
27 & Implementar Speech-to-Text para feedback cualitativo & Q1--Q2 2028 & 25 & ML Engineer Lead & WER $\le$ 10\,\% en transcripciones de entrevistas \\
\midrule
28 & Obtener certificación SOC 2 Tipo II & Q2 2028 & 21, 22 & CEO, CTO & Informe SOC 2 Tipo II emitido por auditor externo \\
\midrule
29 & Lanzar marketplace de integraciones & Q2 2028 & 7, 19 & Developer Advocate, CTO & $\ge$ 5 integraciones listadas por terceros \\
\midrule
30 & Iniciar roadmap ISO 27001 & Q3 2028 & 28 & CTO, SRE & Gap analysis completado + plan de remediación \\
\midrule
31 & Implementar auto-scaling predictivo con ML & Q3 2028 & 16, 18 & ML Engineer, SRE & 0 incidentes por saturación de capacidad \\
\midrule
32 & Expandir cobertura de datos a sustentabilidad (emisiones CO$_2$) & Q4 2028 & 4, 20 & Data Engineer, CEO & Métricas de carbono por ruta disponibles vía API \\
\bottomrule
\end{longtable}

%=====================================================================
% IV. ANÁLISIS ESTRATÉGICO Y DE MERCADO
%=====================================================================

\section{Análisis Estratégico y de Mercado}

\subsection{Análisis FODA (SWOT)}

El análisis FODA sintetiza las condiciones internas y externas que enmarcan la estrategia de SkyAnalytics. Sirve como base para la priorización de iniciativas y la gestión de riesgos.

\begin{table}[H]
\small
\centering
\begin{tabularx}{\linewidth}{>{\raggedright\arraybackslash}p{6.5cm} >{\raggedright\arraybackslash}X}
\toprule
\multicolumn{1}{>{\raggedright\arraybackslash}p{6.5cm}}{\textbf{Fortalezas (Internas)}} &
\multicolumn{1}{>{\raggedright\arraybackslash}X}{\textbf{Debilidades (Internas)}} \\
\midrule
\textbf{F1.} Stack tecnológico moderno y cloud-native: K8s, Terraform, MLflow, Feast, SHAP. Infraestructura 100\,\% como código desde el día 1. & \textbf{D1.} Equipo inicial reducido (fundadores + contratistas) con riesgo de desgaste en fases de alto crecimiento. \\
\textbf{F2.} Modelo de negocio SaaS recurrente con margen bruto proyectado $>$ 75\,\% y baja dependencia de CAPEX inicial. & \textbf{D2.} Dependencia de un número limitado de proveedores de datos aeronáuticos; cualquier cambio de condiciones contractuales impacta la operación. \\
\textbf{F3.} Enfoque en explainability (SHAP) y transparencia de modelos, diferenciador clave frente a competidores que operan como cajas negras. & \textbf{D3.} Marca desconocida en un mercado dominado por jugadores establecidos (OAG, Cirium, FlightAware) con relaciones comerciales de décadas. \\
\textbf{F4.} Cultura Remote-First desde el origen, permitiendo contratar talento global sin restricciones geográficas y con costos competitivos. & \textbf{D4.} PocketBase como base transaccional MVP limita la escalabilidad; requiere optimización para soportar crecimiento de tenants. \\
\textbf{F5.} Arquitectura orientada a ecosistema (APIs, SDKs, marketplace) que habilita adopción orgánica y efectos de red entre desarrolladores. & \textbf{D5.} Sin certificaciones de seguridad vigentes (SOC 2, ISO 27001) en el lanzamiento, lo cual puede ser bloqueante para clientes Enterprise. \\
\textbf{F6.} Pipeline de ML maduro con MLOps completo, feature store, A/B testing y monitoreo de drift — capacidades que los competidores legacy no tienen. & \textbf{D6.} Ausencia de datos propios; la propuesta de valor depende de la calidad y disponibilidad de datos de terceros. \\
\midrule
\multicolumn{1}{>{\raggedright\arraybackslash}p{6.5cm}}{\textbf{Oportunidades (Externas)}} &
\multicolumn{1}{>{\raggedright\arraybackslash}X}{\textbf{Amenazas (Externas)}} \\
\midrule
\textbf{O1.} Mercado global de aviation analytics en crecimiento: se proyecta USD 4.3B en 2024 a USD 7.8B en 2030 (CAGR 10.3\,\%). Demanda impulsada por digitalización post-COVID y presión por eficiencia operativa. & \textbf{A1.} Competidores establecidos (FlightAware, FlightGlobal, OAG, Cirium) con más recursos, datos propios y relaciones comerciales consolidadas pueden replicar funcionalidades de ML. \\
\textbf{O2.} Vacío de mercado en el segmento mid-market: aerolíneas regionales, agencias de viaje medianas y operadores logísticos sin acceso a herramientas avanzadas de predictive analytics. & \textbf{A2.} Cambios regulatorios en protección de datos aeronáuticos (IATA, ICAO, GDPR) pueden imponer restricciones sobre el uso y la distribución de datos de vuelo. \\
\textbf{O3.} Adopción acelerada de APIs como modelo de consumo B2B; las aerolíneas modernizan sus stacks y buscan integraciones vía API en lugar de reportes estáticos. & \textbf{A3.} Proveedores de datos aeronáuticos pueden lanzar sus propias plataformas de analytics compitiendo directamente, eliminando el intermediario. \\
\textbf{O4.} Regulaciones de sostenibilidad (CORSIA, EU ETS) que exigen a las aerolíneas reportar y optimizar emisiones; SkyAnalytics puede incorporar métricas de carbono por ruta. & \textbf{A4.} Concentración de clientes: si los primeros 3--5 clientes Enterprise concentran $>$ 50\,\% del ARR, la pérdida de uno solo impacta severamente la viabilidad financiera. \\
\textbf{O5.} Crecimiento del ecosistema de marketplaces de APIs (RapidAPI, AWS Marketplace, Google Cloud Marketplace) que reducen el costo de distribución y adquisición de clientes. & \textbf{A5.} Avance de modelos LLM generalistas (GPT, Claude, Gemini) que podrían ofrecer predicciones de vuelo básicas como commodity, erosionando la diferenciación técnica si no se profundiza en el dominio. \\
\textbf{O6.} Alianzas estratégicas con aerolíneas como design partners para co-desarrollar el producto, generando credibilidad y referencias en la industria. & \textbf{A6.} Riesgo cambiario y de jurisdicción al operar con clientes y equipo en múltiples países; exposición a fluctuaciones de divisas y complejidad legal. \\
\bottomrule
\end{tabularx}
\caption{Matriz FODA de SkyAnalytics Inc.}
\end{table}

\subsection{Análisis de Mercado}

\subsubsection{Tamaño del Mercado (TAM / SAM / SOM)}

El mercado global de aviation analytics se estima en USD 4.3 mil millones en 2024, con proyección de alcanzar USD 7.8 mil millones en 2030 (CAGR 10.3\,\%). SkyAnalytics se posiciona en la intersección de aviation data, predictive analytics y plataformas API-first.

\begin{table}[H]
\small
\centering
\begin{tabularx}{\linewidth}{>{\raggedright\arraybackslash}p{3cm} >{\raggedright\arraybackslash}p{3.5cm} >{\raggedright\arraybackslash}p{3.5cm} >{\raggedright\arraybackslash}X}
\toprule
\multicolumn{1}{>{\raggedright\arraybackslash}p{3cm}}{\textbf{Nivel}} &
\multicolumn{1}{>{\raggedright\arraybackslash}p{3.5cm}}{\textbf{Alcance}} &
\multicolumn{1}{>{\raggedright\arraybackslash}p{3.5cm}}{\textbf{Estimación (2030)}} &
\multicolumn{1}{>{\raggedright\arraybackslash}X}{\textbf{Supuestos}} \\
\midrule
\textbf{TAM} & Mercado global de datos y analytics de aviación (software, servicios, consultoría) & USD 7.8B & Todas las aerolíneas, aeropuertos, agencias de viaje, operadores logísticos y entes reguladores a nivel global \\
\textbf{SAM} & Plataformas SaaS de predictive analytics aeronáutico (API-first, modelos ML) para aerolíneas y agencias de viaje & USD 2.1B & Aerolíneas comerciales (Top 200 globales) + Top 500 agencias de viaje y operadores logísticos \\
\textbf{SOM} & Mercado alcanzable por SkyAnalytics en 5 años con estrategia de penetración orgánica (SEO, APIs, marketplace) & USD 35--50M ARR & Captura de 2--3\,\% del SAM. Meta aspiracional: 500 clientes activos, ARR promedio \$70--100K por cuenta Enterprise \\
\bottomrule
\end{tabularx}
\caption{Estimación TAM / SAM / SOM para SkyAnalytics}
\end{table}

\subsubsection{Segmentos de Clientes Objetivo}

\begin{table}[H]
\small
\centering
\begin{tabularx}{\linewidth}{>{\raggedright\arraybackslash}p{2.5cm} >{\raggedright\arraybackslash}p{3cm} >{\raggedright\arraybackslash}p{2.2cm} >{\raggedright\arraybackslash}p{2.2cm} X}
\toprule
\multicolumn{1}{>{\raggedright\arraybackslash}p{2.5cm}}{\textbf{Segmento}} &
\multicolumn{1}{>{\raggedright\arraybackslash}p{3cm}}{\textbf{Perfil}} &
\multicolumn{1}{>{\raggedright\arraybackslash}p{2.2cm}}{\textbf{Necesidad Clave}} &
\multicolumn{1}{>{\raggedright\arraybackslash}p{2.2cm}}{\textbf{Plan Recomendado}} &
\multicolumn{1}{>{\raggedright\arraybackslash}X}{\textbf{Estrategia de Captación}} \\
\midrule
Aerolíneas Enterprise (Top 50) & Flotas $>$ 100 aeronaves, múltiples hubs, equipos de datos internos & Predicción de retrasos, optimización de rutas, dashboards ejecutivos, SLA 99.99\,\% & Enterprise (USD 5K--20K/mes) & Venta directa con CTO/VP Operations, prueba de concepto con datos propios, referencias \\
Aerolíneas Regionales & Flotas 20--100 aeronaves, 1--2 hubs, equipos técnicos reducidos & Predicción de retrasos, monitoreo de clima en rutas, API simple & Pro (USD 500--2K/mes) & Inbound vía contenido técnico, marketplace RapidAPI, webinars sectoriales \\
Agencias de Viaje & B2B/B2C, alto volumen de consultas de vuelos, necesitan datos en tiempo real & Consulta de disponibilidad, precios y retrasos vía API para integrar en sus propios sistemas & Pro / API-Only (basado en volumen) & Documentación de integración excepcional, SDKs, sandbox instantáneo \\
Operadores Logísticos & Carga aérea, cadena de frío, time-critical shipments & Predicción de retrasos en rutas de carga, optimización de conexiones & Enterprise / Pro & Casos de uso de ROI cuantificable, integración con TMS existentes \\
Entes Gubernamentales & Autoridades de aviación civil, reguladores de transporte & Monitoreo de puntualidad de la industria, reportes de cumplimiento, dashboards sectoriales & Enterprise & Licitaciones, alianzas con consultoras del sector público \\
\bottomrule
\end{tabularx}
\caption{Segmentos de clientes objetivo de SkyAnalytics}
\end{table}

\subsection{Mapa Competitivo}

SkyAnalytics compite en el espacio de aviation data intelligence, diferenciándose por su enfoque API-first, explainability de modelos ML y cultura Remote-First que permite una estructura de costos más eficiente.

\begin{table}[H]
\small
\centering
\begin{tabularx}{\linewidth}{>{\raggedright\arraybackslash}p{2.2cm} >{\raggedright\arraybackslash}p{2.5cm} >{\raggedright\arraybackslash}p{2.5cm} >{\raggedright\arraybackslash}p{2.5cm} >{\raggedright\arraybackslash}p{2.5cm} >{\raggedright\arraybackslash}X}
\toprule
\multicolumn{1}{>{\raggedright\arraybackslash}p{2.2cm}}{\textbf{Competidor}} &
\multicolumn{1}{>{\raggedright\arraybackslash}p{2.5cm}}{\textbf{Tipo}} &
\multicolumn{1}{>{\raggedright\arraybackslash}p{2.5cm}}{\textbf{Precio Aprox.}} &
\multicolumn{1}{>{\raggedright\arraybackslash}p{2.5cm}}{\textbf{Fortalezas}} &
\multicolumn{1}{>{\raggedright\arraybackslash}p{2.5cm}}{\textbf{Debilidades}} &
\multicolumn{1}{>{\raggedright\arraybackslash}X}{\textbf{Diferenciador SkyAnalytics}} \\
\midrule
FlightAware & Tracking en tiempo real, APIs & \$500--5K/mes & Cobertura global, datos propios, marca consolidada & APIs legacy (SOAP/XML), sin modelos ML avanzados, explainability nula & Modelos ML con SHAP, APIs REST/GraphQL modernas, SDKs auto-generados \\
FlightGlobal / FlightStats & Datos históricos + analytics & \$1K--10K/mes & Datos históricos extensos (10+ años), relaciones con aerolíneas & Plataforma legacy, UX pobre, integración compleja & Dashboards self-service, sandbox instantáneo, TTFV $\leq$ 3 min \\
OAG & Datos de horarios, analytics enterprise & \$2K--15K/mes & Base de datos de horarios más completa del mundo, estándar de industria & Foco en horarios no en predicción, costo prohibitivo para mid-market & Predicción de retrasos en tiempo real, pricing accesible para aerolíneas regionales \\
Cirium (RELX) & Analytics enterprise, data as a service & \$5K--25K/mes & Respaldo de RELX, datos de flota, integración con banca de inversión & Mercado primario: finanzas, no operaciones. Sin self-service real. Precio extremadamente alto. & Self-service, APIs para desarrolladores, cultura DevRel \\
Xweather (anterior AerisWeather) & Datos meteorológicos para aviación & \$200--2K/mes & Meteorología especializada, buenas APIs & Solo clima, sin predicción integrada de retrasos ni analytics de vuelo & Predicción integrada clima + retrasos en un solo endpoint \\
\bottomrule
\end{tabularx}
\caption{Mapa competitivo de SkyAnalytics vs principales actores del mercado}
\end{table}

\textbf{Ventaja Competitiva Principal:} SkyAnalytics es el único competidor que combina APIs modernas (REST/GraphQL), modelos ML con explainability (SHAP), dashboards self-service y una cultura Remote-First que se traduce en precios 30--50\,\% menores que los incumbentes. La ausencia de sistemas legacy permite iterar 3--5 veces más rápido en features de producto.

\subsection{Modelo de Revenue}

\subsubsection{Planes de Suscripción}

SkyAnalytics adopta un modelo SaaS con planes escalonados por volumen de consumo y funcionalidades. La facturación es mensual con descuento por pago anual. El modelo está diseñado para que los clientes crezcan orgánicamente de un plan a otro conforme aumenta su dependencia de la plataforma.

\begin{table}[H]
\small
\centering
\begin{tabularx}{\linewidth}{>{\raggedright\arraybackslash}p{2cm} >{\raggedright\arraybackslash}p{2.8cm} >{\raggedright\arraybackslash}p{2.8cm} >{\raggedright\arraybackslash}p{2.8cm} >{\raggedright\arraybackslash}X}
\toprule
\multicolumn{1}{>{\raggedright\arraybackslash}p{2cm}}{\textbf{Plan}} &
\multicolumn{1}{>{\raggedright\arraybackslash}p{2.8cm}}{\textbf{Precio}} &
\multicolumn{1}{>{\raggedright\arraybackslash}p{2.8cm}}{\textbf{Límites}} &
\multicolumn{1}{>{\raggedright\arraybackslash}p{2.8cm}}{\textbf{Features}} &
\multicolumn{1}{>{\raggedright\arraybackslash}X}{\textbf{Soporte}} \\
\midrule
\textbf{Freemium} & \$0/mes & 1.000 API calls/mes, 1 tenant, datos sintéticos en sandbox & Acceso a 3 endpoints básicos, documentación, sandbox & Comunidad (foro), documentación, sin SLA \\
\textbf{Pro} & \$500/mes (facturación mensual) o \$5.000/año & 50.000 API calls/mes, 5 endpoints, 3 usuarios, 1 tenant & APIs REST + GraphQL, dashboards predefinidos, exportación CSV, webhooks & Email en 8h hábiles, SLA 99.9\,\% \\
\textbf{Enterprise} & \$2.000--\$20.000/mes (según volumen) & API calls ilimitadas (fair use), usuarios ilimitados, multi-tenant & Todo Pro + dashboards custom, datos reales, SDKs prioritarios, feature store compartido, A/B testing de modelos & Slack compartido, SLA 99.99\,\%, gerente de cuenta dedicado \\
\textbf{API-Only} & \$0.01--\$0.05 por API call & Volumen variable, pago por uso & Acceso a todos los endpoints vía API, sin dashboards ni UI & Email en 24h, SLA 99.9\,\% \\
\bottomrule
\end{tabularx}
\caption{Planes de suscripción de SkyAnalytics}
\end{table}

\subsubsection{Estrategia de Pricing}

\begin{itemize}[leftmargin=1.5cm, label=\textbullet, itemsep=2pt, topsep=4pt]
    \item \textbf{Penetración inicial (2026--2027):} Freemium generoso + Pro a precio por debajo del mercado para construir base de usuarios y referencias.
    \item \textbf{Expansión (2027--2028):} Ajuste de pricing validado con A/B testing y entrevistas de willingness-to-pay. Introducción de add-ons (más endpoints, feature store, entrenamiento custom).
    \item \textbf{Madurez (2028+):} Precios alineados al valor entregado (value-based pricing). Descuentos por volumen y contratos anuales. Revenue expansion vía upselling de Pro a Enterprise.
\end{itemize}

\subsection{Registro de Riesgos}

La gestión de riesgos es un componente esencial del gobierno corporativo de SkyAnalytics. Cada riesgo identificado se evalúa en probabilidad e impacto (escala 1--5), y se define una estrategia de mitigación con responsable asignado.

\begin{table}[H]
\small
\centering
\begin{tabularx}{\linewidth}{>{\raggedright\arraybackslash}p{0.6cm} >{\raggedright\arraybackslash}p{2.6cm} >{\raggedright\arraybackslash}p{1.2cm} >{\raggedright\arraybackslash}p{0.9cm} >{\raggedright\arraybackslash}p{0.8cm} >{\raggedright\arraybackslash}p{3cm} >{\raggedright\arraybackslash}X}
\toprule
\multicolumn{1}{>{\raggedright\arraybackslash}p{0.6cm}}{\textbf{\#}} &
\multicolumn{1}{>{\raggedright\arraybackslash}p{2.6cm}}{\textbf{Riesgo}} &
\multicolumn{1}{>{\raggedright\arraybackslash}p{1.2cm}}{\textbf{Probab.}} &
\multicolumn{1}{>{\raggedright\arraybackslash}p{0.9cm}}{\textbf{Impacto}} &
\multicolumn{1}{>{\raggedright\arraybackslash}p{0.8cm}}{\textbf{Score}} &
\multicolumn{1}{>{\raggedright\arraybackslash}p{3cm}}{\textbf{Mitigación}} &
\multicolumn{1}{>{\raggedright\arraybackslash}X}{\textbf{Owner}} \\
\midrule
R01 & Proveedor de datos interrumpe el servicio o modifica condiciones & Media & 5 & 15 & Contratos con 2+ proveedores alternativos, cláusula de continuidad, caché de datos históricos & CTO \\
R02 & Vulnerabilidad de seguridad crítica en producción & Baja & 5 & 10 & SAST/DAST en CI/CD, parches $\leq$ 24h, WAF, Zero Trust, auditorías trimestrales & SRE Lead \\
R03 & Salida de CTO o ML Engineer Lead (key person risk) & Media & 4 & 12 & Documentación exhaustiva (ADRs, runbooks, wiki), bus factor $\geq$ 2 en cada área, plan de sucesión & CEO \\
R04 & No alcanzar product-market fit en 18 meses & Media & 5 & 15 & Ciclo continuo de feedback, métricas de engagement, pivot temprano si los KPIs de adopción no se alcanzan & CEO \\
R05 & Competidor grande lanza producto similar con pricing predatorio & Alta & 3 & 12 & Diferenciación en explainability y DevRel, nicho mid-market desatendido, switching cost vía integración profunda & CEO \\
R06 & Ingesta de datos supera capacidad del pipeline y genera lag & Media & 3 & 9 & Auto-scaling de workers ETL, alertas de lag, arquitectura de streaming para fuentes críticas & Data Engineer Lead \\
R07 & Modelo ML degrada (drift) y entrega predicciones erróneas sin ser detectado & Media & 4 & 12 & Monitoreo PSI diario, A/B testing en staging, rollback automático, alertas & ML Engineer Lead \\
R08 & Costos cloud superan el presupuesto por error de configuración o ataque & Baja & 4 & 8 & Alertas de anomalía de gasto diario, budget caps por cuenta, revisión semanal & SRE Lead \\
R09 & Cambio regulatorio restringe el uso de datos aeronáuticos en ciertas jurisdicciones & Baja & 3 & 6 & Monitoreo regulatorio proactivo, residencia de datos configurable por región, asesoría legal continua & CEO \\
R10 & Incidente Sev1 no resuelto en $\leq$ 15 minutos por falta de cobertura on-call & Baja & 4 & 8 & Rotación on-call 24/7, runbooks actualizados, escalation automática, post-mortem obligatorio & SRE Lead \\
R11 & Cliente Enterprise concentra $>$ 40\,\% del ARR & Media & 4 & 12 & Diversificación activa de cartera, límite interno de concentración por cliente, expansión del pipeline de ventas & VP Ventas \\
R12 & Falla en backup o restore durante un incidente real de pérdida de datos & Baja & 5 & 10 & Backup diario encriptado + restore semanal verificado, alerta de integridad de backup, DR multi-región & SRE Lead \\
\bottomrule
\end{tabularx}
\caption{Registro de riesgos con probabilidad, impacto, score, mitigación y responsable}
\end{table}

\subsection{Proyección Financiera 3 Años}

La proyección financiera presenta un escenario base asumiendo ejecución exitosa del plan de acción. Las cifras son estimaciones para fines de planificación estratégica y no constituyen compromisos financieros vinculantes.

\begin{table}[H]
\small
\centering
\begin{tabularx}{\linewidth}{>{\raggedright\arraybackslash}p{2cm} >{\raggedright\arraybackslash}p{2.5cm} >{\raggedright\arraybackslash}p{2.5cm} >{\raggedright\arraybackslash}p{2.5cm} >{\raggedright\arraybackslash}p{2.5cm} >{\raggedright\arraybackslash}X}
\toprule
\multicolumn{1}{>{\raggedright\arraybackslash}p{2cm}}{\textbf{Indicador}} &
\multicolumn{1}{>{\raggedright\arraybackslash}p{2.5cm}}{\textbf{2026 (MVP)}} &
\multicolumn{1}{>{\raggedright\arraybackslash}p{2.5cm}}{\textbf{2027 (Growth)}} &
\multicolumn{1}{>{\raggedright\arraybackslash}p{2.5cm}}{\textbf{2028 (Scale)}} &
\multicolumn{1}{>{\raggedright\arraybackslash}p{2.5cm}}{\textbf{Fórmula / Supuesto}} &
\multicolumn{1}{>{\raggedright\arraybackslash}X}{\textbf{Notas}} \\
\midrule
Clientes activos & 50 (Beta) & 100 -- 200 & 500+ & Crecimiento compuesto vía marketing + marketplace & Beta testers no generan revenue en 2026 \\
\midrule
ARR & \$0 (pre-revenue) & \$100K -- \$200K & \$1.5M -- \$3M & ARR = clientes $\times$ ARPU. ARPU crece de \$500 a \$5K+ & Primeros pagos estimados Q1 2027 \\
\midrule
Costo Cloud Mensual & \$2K -- \$5K & \$8K -- \$15K & \$25K -- \$50K & Escala con volumen de datos y tráfico API & Mayor driver: K8s nodos + DB columnar + egress \\
\midrule
Costo Personas & Variable (equipo inicial) & 5--12 personas & 15--30 personas & Remote-First contractors + full-time hires & Principal partida de costo operativo \\
\midrule
Gross Margin & No aplica & 60\,\% -- 75\,\% & $>$ 75\,\% & (ARR - costo cloud) / ARR & Mejora con optimización cloud y economías de escala \\
\midrule
Runway & 18 -- 24 meses (fondeo inicial) & Break-even o Serie A & Rentable o Serie B & Cash flow: ingresos - egresos operativos & Hito: alcanzar break-even mensual en Q3 2027 \\
\midrule
LTV & No aplica & \$5K -- \$15K & \$20K -- \$60K & ARPU $\times$ vida media del cliente en meses & Crece con upselling y reducción de churn \\
\midrule
CAC & No aplica & \$300 -- \$500 & \$200 -- \$400 & Gasto marketing + ventas / clientes nuevos & Decrece con inbound orgánico y referidos \\
\midrule
LTV/CAC & No aplica & $>$ 3x & $>$ 5x & LTV / CAC & Indicador de eficiencia comercial \\
\bottomrule
\end{tabularx}
\caption{Proyección financiera estimada a 3 años (escenario base)}
\end{table}

%=====================================================================
% V. VISIÓN ARQUITECTÓNICA Y ANALÍTICA
%=====================================================================

\section{Visión Arquitectónica y Analítica}

\subsection{Resumen Analítico por Nivel Organizacional}

\begin{table}[H]
\small
\centering
\begin{tabularx}{\linewidth}{>{\raggedright\arraybackslash}p{1.8cm} >{\raggedright\arraybackslash}p{2.5cm} >{\raggedright\arraybackslash}p{3cm} >{\raggedright\arraybackslash}p{2cm} >{\raggedright\arraybackslash}X}
\toprule
\textbf{Nivel} & \textbf{Decisión} & \textbf{Técnicas} & \textbf{Casos de Uso} & \textbf{Resultado} \\
\midrule
Estratégico & Gerencial (largo plazo) & BI, ML, simulación & CU-E01 a E06 & Visibilidad 360°, alineación de metas \\
\midrule
Táctico & Planeación por área & Forecasting, optimización & CU-T01 a T10 & Eficiencia, control de costos \\
\midrule
Operativo & Ejecución diaria & Reglas, validaciones, alertas & CU-O01 a O20 & Confiabilidad, ejecución de pipelines \\
\bottomrule
\end{tabularx}
\caption{Resumen analítico integrado por nivel organizacional}
\end{table}

%---------------------------------------------------------------------
\subsection{Modelo Analítico General del Sistema}

\begin{table}[H]
\small
\centering
\begin{tabularx}{\linewidth}{>{\centering\arraybackslash}p{2.0cm} c >{\centering\arraybackslash}p{2.0cm} c >{\centering\arraybackslash}p{1.8cm} c >{\centering\arraybackslash}p{1.8cm} c >{\centering\arraybackslash}p{1.8cm} c >{\centering\arraybackslash}p{1.8cm}}
\toprule
\textbf{APIs + Dev Portal} & $\rightarrow$ & \textbf{BD Transaccional} & $\rightarrow$ & \textbf{ETL + Validación} & $\rightarrow$ & \textbf{Data Warehouse} & $\rightarrow$ & \textbf{BI + ML Pipeline} & $\rightarrow$ & \textbf{Reportes} \\
\midrule
REST/GraphQL, OAuth2, Rate Limits, SDKs & & PocketBase, registro tenants, consumo API & & Great Expectations, dbt, Airflow & & MonetDB columnar, catálogo datos & & PowerBI, MLflow, Feast, SHAP & & Dashboards, predicciones, alertas \\
\bottomrule
\end{tabularx}
\caption{Flujo de datos y componentes del ecosistema SkyAnalytics}
\end{table}

%---------------------------------------------------------------------
\subsection{Tablas de Hechos Principales (Fact)}

\begin{table}[H]
\small
\centering
\begin{tabularx}{\linewidth}{>{\raggedright\arraybackslash}p{3.5cm} >{\raggedright\arraybackslash}p{3.5cm} >{\raggedright\arraybackslash}X}
\toprule
\multicolumn{1}{>{\raggedright\arraybackslash}p{3.5cm}}{\textbf{Tabla Fact}} &
\multicolumn{1}{>{\raggedright\arraybackslash}p{3.5cm}}{\textbf{Qué Mide}} &
\multicolumn{1}{>{\raggedright\arraybackslash}X}{\textbf{Ejemplos de Métricas}} \\
\midrule
Fact\_API\_Call & Consumo API por cliente & Total llamadas, latencia, errores, banda \\
\midrule
Fact\_Vuelo & Datos de vuelos procesados & Vuelos registrados, puntualidad, cancelaciones \\
\midrule
Fact\_Retraso & Predicciones de retraso & Retraso predicho vs real, MAPE, precisión \\
\midrule
Fact\_Suscripcion & Ingresos recurrentes & ARR, MRR, expansión, downgrades, churn \\
\midrule
Fact\_Ingesta\_ETL & Ejecución pipelines & Volumen ingerido, tiempo ejecución, errores \\
\midrule
Fact\_Modelo\_ML & Desempeño modelos ML & Precisión, recall, drift, tiempo entrenamiento \\
\midrule
Fact\_Satisfaccion & Encuestas clientes & NPS, CSAT, comentarios, recomendación \\
\midrule
Fact\_Incidente & Incidentes operativos & Tiempo detección, MTTR, severidad \\
\bottomrule
\end{tabularx}
\caption{Tablas de hechos principales del Data Warehouse}
\end{table}

%---------------------------------------------------------------------
\subsection{Dimensiones Principales (Dim)}

\begin{table}[H]
\small
\centering
\begin{tabularx}{\linewidth}{>{\raggedright\arraybackslash}p{4cm} >{\raggedright\arraybackslash}X}
\toprule
\multicolumn{1}{>{\raggedright\arraybackslash}p{4cm}}{\textbf{Dimensión}} &
\multicolumn{1}{>{\raggedright\arraybackslash}X}{\textbf{Uso}} \\
\midrule
Dim\_Tiempo & Día, semana, mes, trimestre, año \\
\midrule
Dim\_Aerolinea & Código IATA/ICAO, país, alianza, tipo \\
\midrule
Dim\_Aeropuerto & Código IATA/ICAO, ciudad, país, región, zona horaria \\
\midrule
Dim\_Ruta & Origen-destino, distancia, tipo ruta \\
\midrule
Dim\_Cliente & Empresa B2B, plan, industria, región \\
\midrule
Dim\_Plan & Tipo suscripción, límites API, features, precio \\
\midrule
Dim\_Endpoint & Endpoint API, versión, método, módulo \\
\midrule
Dim\_Empleado & Rol, equipo, seniority, ubicación \\
\bottomrule
\end{tabularx}
\caption{Dimensiones principales del modelo de datos}
\end{table}

%---------------------------------------------------------------------
\subsection{Técnicas Usadas por Nivel Organizacional}

\subsubsection{Nivel Estratégico}

\begin{table}[H]
\small
\centering
\begin{tabularx}{\linewidth}{>{\raggedright\arraybackslash}p{4cm} >{\raggedright\arraybackslash}X}
\toprule
\multicolumn{1}{>{\raggedright\arraybackslash}p{4cm}}{\textbf{Técnica}} &
\multicolumn{1}{>{\raggedright\arraybackslash}X}{\textbf{Uso en el Nivel Estratégico}} \\
\midrule
BI y agregaciones & Consolidar KPIs en dashboards ejecutivos con comparativas históricas \\
\midrule
ML predictivo & Proyectar ARR, churn y demanda de API con modelos de series temporales \\
\midrule
Detección de anomalías & Identificar desviaciones en revenue, uptime y satisfacción \\
\midrule
Segmentación de clientes & Clasificar cuentas por valor, riesgo y potencial de expansión \\
\midrule
Recomendación & Sugerir planes y endpoints a clientes según su perfil de consumo \\
\midrule
Dashboards BSC & Visualizar semáforos y tendencias de todos los KPIs estratégicos \\
\bottomrule
\end{tabularx}
\caption{Técnicas analíticas del nivel estratégico}
\end{table}

\subsubsection{Nivel Táctico}

\begin{table}[H]
\small
\centering
\begin{tabularx}{\linewidth}{>{\raggedright\arraybackslash}p{4cm} >{\raggedright\arraybackslash}X}
\toprule
\multicolumn{1}{>{\raggedright\arraybackslash}p{4cm}}{\textbf{Técnica}} &
\multicolumn{1}{>{\raggedright\arraybackslash}X}{\textbf{Uso en el Nivel Táctico}} \\
\midrule
Agregaciones por área & Reportes de consumo API por endpoint, costos cloud por servicio \\
\midrule
Forecasting & Prever picos de tráfico, necesidades de capacidad y presupuesto \\
\midrule
Alertas inteligentes & Disparar notificaciones ante brechas de SLA, calidad o seguridad \\
\midrule
Optimización de recursos & Identificar recursos cloud ociosos y redimensionar instancias \\
\midrule
Análisis de campañas & Medir conversión, ROI y segmentación de campañas de growth \\
\bottomrule
\end{tabularx}
\caption{Técnicas analíticas del nivel táctico}
\end{table}

\subsubsection{Nivel Operativo}

\begin{table}[H]
\small
\centering
\begin{tabularx}{\linewidth}{>{\raggedright\arraybackslash}p{4cm} >{\raggedright\arraybackslash}X}
\toprule
\multicolumn{1}{>{\raggedright\arraybackslash}p{4cm}}{\textbf{Técnica}} &
\multicolumn{1}{>{\raggedright\arraybackslash}X}{\textbf{Uso en el Nivel Operativo}} \\
\midrule
Reglas de negocio & Validar rate limits, permisos de plan, autenticación OAuth2 \\
\midrule
Validaciones automáticas & Ejecutar suites de Great Expectations tras cada ingesta \\
\midrule
Alertas operativas & Notificar fallos de pipeline, saturación de pods, errores 500 \\
\midrule
Recomendaciones simples & Sugerir endpoints relacionados según historial de consumo \\
\midrule
NLP para chatbot & Clasificar tickets, buscar FAQs semánticamente, detectar sentimiento \\
\bottomrule
\end{tabularx}
\caption{Técnicas analíticas del nivel operativo}
\end{table}

%---------------------------------------------------------------------
\subsection{Matriz de Casos de Uso Estratégicos}

\begin{table}[H]
\small
\centering
\begin{tabularx}{\linewidth}{>{\raggedright\arraybackslash}p{2cm} >{\raggedright\arraybackslash}p{3.5cm} >{\raggedright\arraybackslash}p{3.5cm} >{\raggedright\arraybackslash}X}
\toprule
\multicolumn{1}{>{\raggedright\arraybackslash}p{2cm}}{\textbf{CU}} &
\multicolumn{1}{>{\raggedright\arraybackslash}p{3.5cm}}{\textbf{Técnicas}} &
\multicolumn{1}{>{\raggedright\arraybackslash}p{3.5cm}}{\textbf{Modelo Fact-Dim}} &
\multicolumn{1}{>{\raggedright\arraybackslash}X}{\textbf{Reportes}} \\
\midrule
CU-E01 BSC & BI, agregaciones, semáforos, ML & Consultar sección Tablas de Hechos & Dashboard BSC \\
\midrule
CU-E02 ARR & Agregaciones financieras, anomalías & Consultar sección Tablas de Hechos & Reporte rentabilidad \\
\midrule
CU-E03 Uptime & Agregaciones, MTTR, error budget & Consultar sección Tablas de Hechos & Dashboard SLA/SLO \\
\midrule
CU-E04 Compliance & Agregaciones documentales & Consultar sección Tablas de Hechos & Reporte cumplimiento \\
\midrule
CU-E05 Metas & Simulación, predicción escenarios & Consultar sección Tablas de Hechos & Reporte metas trimestrales \\
\midrule
CU-E06 Talento & Agregaciones, retención & Consultar sección Tablas de Hechos & Dashboard People Analytics \\
\bottomrule
\end{tabularx}
\caption{Matriz de casos de uso estratégicos: técnicas, modelo Fact-Dim y reportes}
\end{table}

%---------------------------------------------------------------------
\subsection{Matriz de Casos de Uso Tácticos}

\begin{table}[H]
\small
\centering
\begin{tabularx}{\linewidth}{>{\raggedright\arraybackslash}p{2cm} >{\raggedright\arraybackslash}p{3.5cm} >{\raggedright\arraybackslash}p{3.5cm} >{\raggedright\arraybackslash}X}
\toprule
\multicolumn{1}{>{\raggedright\arraybackslash}p{2cm}}{\textbf{CU}} &
\multicolumn{1}{>{\raggedright\arraybackslash}p{3.5cm}}{\textbf{Técnicas}} &
\multicolumn{1}{>{\raggedright\arraybackslash}p{3.5cm}}{\textbf{Modelo Fact-Dim}} &
\multicolumn{1}{>{\raggedright\arraybackslash}X}{\textbf{Reportes}} \\
\midrule
CU-T01 Growth & Segmentación, conversión, scoring & Consultar sección Tablas de Hechos & Reporte campañas \\
\midrule
CU-T02 API & API analytics, monitoreo & Consultar sección Tablas de Hechos & Dashboard uso API \\
\midrule
CU-T03 Infra & Monitoreo tiempo real, alertas & Consultar sección Tablas de Hechos & Dashboard cloud \\
\midrule
CU-T04 ETL & Validación, alertas frescura & Consultar sección Tablas de Hechos & Dashboard calidad \\
\midrule
CU-T05 MLflow & Experimentos, tracking métricas & Consultar sección Tablas de Hechos & Dashboard MLflow \\
\midrule
CU-T06 Security & SAST/DAST, detección vulns & Consultar sección Tablas de Hechos & Reporte vulnerabilidades \\
\midrule
CU-T07 DevPortal & Analíticas tráfico, engagement & Consultar sección Tablas de Hechos & Dashboard portal \\
\midrule
CU-T08 Costos & Análisis costos, anomalías & Consultar sección Tablas de Hechos & Dashboard costos \\
\midrule
CU-T09 Carga & Métricas rendimiento, k6 & Consultar sección Tablas de Hechos & Reporte pruebas \\
\midrule
CU-T10 Pricing & A/B testing, WTP & Consultar sección Tablas de Hechos & Reporte validación \\
\bottomrule
\end{tabularx}
\caption{Matriz de casos de uso tácticos: técnicas, modelo Fact-Dim y reportes}
\end{table}

%---------------------------------------------------------------------
\subsection{Matriz de Casos de Uso Operativos}

\begin{table}[H]
\small
\centering
\begin{tabularx}{\linewidth}{>{\raggedright\arraybackslash}p{2cm} >{\raggedright\arraybackslash}p{3.5cm} >{\raggedright\arraybackslash}p{3.5cm} >{\raggedright\arraybackslash}X}
\toprule
\multicolumn{1}{>{\raggedright\arraybackslash}p{2cm}}{\textbf{CU}} &
\multicolumn{1}{>{\raggedright\arraybackslash}p{3.5cm}}{\textbf{Técnicas Usadas}} &
\multicolumn{1}{>{\raggedright\arraybackslash}p{3.5cm}}{\textbf{Modelo Fact-Dim}} &
\multicolumn{1}{>{\raggedright\arraybackslash}X}{\textbf{Registros / Salidas}} \\
\midrule
CU-O01 Registro & Validación duplicados, auto-gen & Consultar sección Tablas de Hechos & API Key, tenant activo \\
\midrule
CU-O02 Consulta & Rate limiting, caching, auth & Consultar sección Tablas de Hechos & JSON response \\
\midrule
CU-O03 Ingesta & Descarga programada, schema & Consultar sección Tablas de Hechos & Datos en staging \\
\midrule
CU-O04 Calidad & Great Expectations checks & Consultar sección Tablas de Hechos & Reporte pass/fail \\
\midrule
CU-O05 Reentrenar & MLflow tracking, comparación & Consultar sección Tablas de Hechos & Modelo registrado \\
\midrule
CU-O06 Dashboards & Refresh programado, verificación & Consultar sección Tablas de Hechos & Dashboards actualizados \\
\midrule
CU-O07 Telemetría & Recolección métricas, alertas & Consultar sección Tablas de Hechos & Alertas PagerDuty \\
\midrule
CU-O08 Tickets & Triaging NLP, clasificación & Consultar sección Tablas de Hechos & Ticket resuelto \\
\midrule
CU-O09 Backup & Backup automático, verificación & Consultar sección Tablas de Hechos & Backup + log restore \\
\midrule
CU-O10 ErrorBudget & Cálculo SLI/SLO, umbrales & Consultar sección Tablas de Hechos & Decisión congelar/liberar \\
\midrule
CU-O11 Changelog & Generación commits, SemVer & Consultar sección Tablas de Hechos & Changelog portal \\
\midrule
CU-O12 IAM & Rotación automática, verificación & Consultar sección Tablas de Hechos & Credenciales rotadas \\
\midrule
CU-O13 PostMortem & Documentación, timeline auto & Consultar sección Tablas de Hechos & Post-mortem + tickets \\
\midrule
CU-O14 FAQ & Clasificación NLP, KB & Consultar sección Tablas de Hechos & FAQ publicado \\
\midrule
CU-O15 Sprints & Gestión backlog, standups, retrospectivas & Consultar sección Tablas de Hechos & Sprint cerrado con métricas \\
\midrule
CU-O16 OpenAPI & Spectral linter, detección breaking changes & Consultar sección Tablas de Hechos & Spec validada, merge aprobado/bloqueado \\
\midrule
CU-O17 DataContracts & Validación schema, detección drift & Consultar sección Tablas de Hechos & Contrato validado o ingesta pausada \\
\midrule
CU-O18 SAST/DAST & SonarQube + OWASP ZAP, bloqueo críticos & Consultar sección Tablas de Hechos & Reporte de vulnerabilidades por release \\
\midrule
CU-O19 Costos & Análisis anomalías, derechosizing & Consultar sección Tablas de Hechos & Reporte de optimización semanal \\
\midrule
CU-O20 DeepLearning & OCR (CRNN/TrOCR), STT (Whisper), MLflow & Consultar sección Tablas de Hechos & Experimentos documentados con recomendación go/no-go \\
\bottomrule
\end{tabularx}
\caption{Matriz de casos de uso operativos: técnicas, modelo Fact-Dim y registros}
\end{table}

%---------------------------------------------------------------------
\subsection{Agregaciones Usadas en el Sistema}

\begin{table}[H]
\small
\centering
\begin{tabularx}{\linewidth}{>{\raggedright\arraybackslash}p{3.5cm} >{\raggedright\arraybackslash}p{5cm} >{\raggedright\arraybackslash}X}
\toprule
\multicolumn{1}{>{\raggedright\arraybackslash}p{3.5cm}}{\textbf{Agregación}} &
\multicolumn{1}{>{\raggedright\arraybackslash}p{5cm}}{\textbf{Fórmula o Cálculo}} &
\multicolumn{1}{>{\raggedright\arraybackslash}X}{\textbf{Uso}} \\
\midrule
Total API Calls & SUM(llamadas) GROUP BY día, endpoint, cliente & Facturación, dashboards \\
\midrule
Latencia p95 & PERCENTILE(latencia, 0.95) GROUP BY endpoint, hora & Monitoreo SLA \\
\midrule
ARR & SUM(mrr) WHERE activo $\times$ 12 & Reporte financiero \\
\midrule
Churn Rate & COUNT(cancelados) / COUNT(total) $\times$ 100 & Retención \\
\midrule
NRR & (ARR inicial + expansión - churn) / ARR inicial $\times$ 100 & Crecimiento orgánico \\
\midrule
NPS & (\% promotores - \% detractores) $\times$ 100 & Satisfacción \\
\midrule
MAPE & AVG(ABS(real - predicho) / real) $\times$ 100 & Evaluación ML \\
\midrule
Data Freshness & COUNT(datos \textless\ 5 min) / COUNT(total) $\times$ 100 & Calidad datos \\
\midrule
Uptime Real & (tiempo total - caído) / tiempo total $\times$ 100 & SLA compliance \\
\midrule
Error Budget Consumido & SUM(tiempo caído mensual) / 4.38 min $\times$ 100 & Gestión riesgo \\
\midrule
Tasa Errores API & COUNT(5xx) / COUNT(total) $\times$ 100 & Confiabilidad \\
\midrule
TTFV Promedio & AVG(timestamp\_primer\_insight - timestamp\_registro) & Onboarding \\
\bottomrule
\end{tabularx}
\caption{Agregaciones utilizadas en el sistema SkyAnalytics}
\end{table}

%---------------------------------------------------------------------
\subsection{Técnicas de IA y Machine Learning Aplicables}

\subsubsection{Segmentación de Clientes}

Clustering (K-Means, DBSCAN) para agrupar aerolíneas por patrón de consumo; RFM para clasificar por recencia, frecuencia y monto; Árboles de decisión para identificar probabilidad de expansión. Entrada: historial de API calls, plan, industria. Salida: segmento (high-value, growth, at-risk).

\subsubsection{Predicción de Retrasos de Vuelo}

XGBoost/LightGBM con features: aerolínea, aeropuerto, hora, meteorología, historial de ruta, tráfico del sector, temporada. Salida: probabilidad y magnitud de retraso, clasificación de severidad, ventana óptima de salida.

\subsubsection{Predicción de Demanda de API}

Prophet / ARIMA / LSTM para prever picos de tráfico por endpoint con ventanas de 15 minutos, permitiendo auto-scaling proactivo y dimensionamiento de capacidad.

\subsubsection{Detección de Anomalías}

\begin{table}[H]
\small
\centering
\begin{tabularx}{\linewidth}{>{\raggedright\arraybackslash}p{4cm} >{\raggedright\arraybackslash}X}
\toprule
\multicolumn{1}{>{\raggedright\arraybackslash}p{4cm}}{\textbf{Anomalía}} &
\multicolumn{1}{>{\raggedright\arraybackslash}X}{\textbf{Posible Causa}} \\
\midrule
Errores 500 aumentan & Bug en código, saturación de DB, ataque DDoS \\
\midrule
Consumo API cae & Migración a competidor, fallo de autenticación OAuth2 \\
\midrule
Precisión ML degrada & Data drift, feature no disponible en fuente \\
\midrule
Volumen datos anómalo & Error del proveedor, duplicación en ingesta \\
\midrule
Tickets soporte aumentan & Bug no detectado, cambio de API sin documentación \\
\bottomrule
\end{tabularx}
\caption{Anomalías operativas y sus posibles causas}
\end{table}

\subsubsection{NLP para Chatbot de Soporte}

Clasificación de tickets para asignar categoría y prioridad automáticamente; búsqueda semántica con embeddings de documentación; generación de FAQs agrupando tickets similares; análisis de sentimiento para detectar clientes frustrados y escalar a un agente humano.

\subsubsection{Recomendación de Endpoints y Planes}

Filtrado colaborativo combinado con content-based filtering. Entrada: historial de API, industria, tamaño de flota. Salida: endpoints sugeridos y plan recomendado según perfil de consumo.

%---------------------------------------------------------------------
\subsection{Técnicas de Deep Learning Aplicables}

\begin{table}[H]
\small
\centering
\begin{tabularx}{\linewidth}{>{\raggedright\arraybackslash}p{3.5cm} >{\raggedright\arraybackslash}p{5cm} >{\raggedright\arraybackslash}X}
\toprule
\multicolumn{1}{>{\raggedright\arraybackslash}p{3.5cm}}{\textbf{Técnica}} &
\multicolumn{1}{>{\raggedright\arraybackslash}p{5cm}}{\textbf{Uso Posible}} &
\multicolumn{1}{>{\raggedright\arraybackslash}X}{\textbf{CU Relacionado}} \\
\midrule
OCR con redes neuronales & Procesar documentos de compliance, facturas de proveedores & CU-E04 \\
\midrule
Análisis de sentimiento & Interpretar comentarios de encuestas NPS y tickets de soporte & CU-O08, CU-O13 \\
\midrule
Clasificación de imágenes & Verificar evidencias de auditoría (screenshots, logs escaneados) & CU-E04 \\
\midrule
Detección de anomalías visuales & Monitorear dashboards de Grafana para detectar patrones anómalos & CU-T03 \\
\midrule
Speech-to-Text & Transcribir entrevistas de feedback con clientes para análisis cualitativo & CU-O20 \\
\bottomrule
\end{tabularx}
\caption{Técnicas de Deep Learning aplicables al ecosistema SkyAnalytics}
\end{table}

%---------------------------------------------------------------------
\subsection{Relación entre Registros Operativos y Reportes Gerenciales}

\begin{center}
\begin{minipage}{0.95\linewidth}
\textbf{Registro Operativo} (API call / Pipeline ETL / Encuesta / Incidente) $\rightarrow$ \textbf{Agregación} (API calls por endpoint, satisfacción promedio, uptime mensual) $\rightarrow$ \textbf{Modelo Fact-Dim} (Fact\_API\_Call + Dim\_Tiempo + Dim\_Cliente + Dim\_Endpoint) $\rightarrow$ \textbf{Reporte Táctico o Estratégico} (rentabilidad, uptime, fidelización, calidad de datos) $\rightarrow$ \textbf{Decisión} (promoción de plan, optimización de infraestructura, capacitación, mejora de proceso)
\end{minipage}
\end{center}

%---------------------------------------------------------------------
\subsection{Ejemplo Completo de Trazabilidad}

\begin{description}[leftmargin=*,style=nextline]
    \item[\textbf{Objetivo Estratégico:}] OE5 --- Innovación en IA Aeronáutica.
    \item[\textbf{Objetivo Táctico:}] OT9 --- Entrenar modelos ML con MLflow, Feast, SHAP y A/B testing.
    \item[\textbf{Objetivo Operativo:}] OP5 --- Reentrenar modelos con datos frescos, verificar drift.
    \item[\textbf{Casos de Uso:}] Estratégico CU-E01 (consultar BSC con KPIs de IA), Táctico CU-T05 (entrenar modelos), Operativo CU-O05 (reentrenar modelo).
    \item[\textbf{Modelo Fact-Dim:}] Fact\_Modelo\_ML, Fact\_Retraso, Fact\_Vuelo + Dim\_Tiempo, Dim\_Aerolinea, Dim\_Ruta.
    \item[\textbf{Agregaciones:}] MAPE ($\text{AVG}(\lvert\text{real} - \text{predicho}\rvert / \text{real}) \times 100$), Precisión por ruta, Drift detection rate, Tiempo de entrenamiento, Experimentos A/B por trimestre.
    \item[\textbf{Técnicas:}] XGBoost/LightGBM, SHAP explainability, MLflow tracking, Feast feature store, A/B testing framework, PSI drift monitoring.
    \item[\textbf{Flujo Completo:}] Datos de vuelo (Fact\_Vuelo) $\rightarrow$ Features en Feast $\rightarrow$ Entrenamiento en MLflow $\rightarrow$ Predicciones (Fact\_Retraso) $\rightarrow$ Métricas (Fact\_Modelo\_ML) $\rightarrow$ Dashboard BSC $\rightarrow$ Decisión de mejora del modelo.
\end{description}

%---------------------------------------------------------------------


%---------------------------------------------------------------------
\subsection{Hoja de Ruta 2026--2028}

El roadmap se despliega en siete sub-fases trimestrales con entregables específicos, hitos de validación y metas intermedias de crecimiento. Cada sub-fase construye sobre la anterior, siguiendo una progresión lógica de MVP $\rightarrow$ Growth $\rightarrow$ Scale.

\subsubsection{Fase 1 --- MVP (2026): Validación de Producto}

\textbf{Q3 2026 --- Fundación Técnica y Legal}
\begin{itemize}[leftmargin=1.5cm, label=\textbullet, itemsep=2pt, topsep=4pt]
    \item Despliegue de infraestructura K8s con Terraform (IaC) y CI/CD con GitHub Actions.
    \item Publicación de API Beta con 3 endpoints en RapidAPI, autenticación OAuth2 y rate limiting por plan.
    \item Configuración de flujos HubSpot para captación de Beta Testers.
    \item Redacción y publicación de ToS, Privacy Policy y DPA.
    \item Lanzamiento de landing page con registro self-service.
    \item \textbf{Hito:} infraestructura versionada, API pública, documentación legal lista.
\end{itemize}

\textbf{Q4 2026 --- Lanzamiento del MVP}
\begin{itemize}[leftmargin=1.5cm, label=\textbullet, itemsep=2pt, topsep=4pt]
    \item Construcción del pipeline ETL completo con Great Expectations, dbt y MonetDB.
    \item Entrenamiento del primer modelo de predicción de retrasos (XGBoost) con MAPE $\le$ 15\,\%.
    \item Implementación del stack de observabilidad (Grafana LGTM + PagerDuty).
    \item Lanzamiento del Developer Portal interactivo con sandbox de datos sintéticos.
    \item Onboarding de primeros 50 Beta Testers activos.
    \item \textbf{Hito:} MVP funcional, 50 Beta Testers, uptime $\ge$ 99.9\,\%.
\end{itemize}

\subsubsection{Fase 2 --- Growth (2027): Escalamiento y Validación Comercial}

\textbf{Q1 2027 --- Validación de Robustez}
\begin{itemize}[leftmargin=1.5cm, label=\textbullet, itemsep=2pt, topsep=4pt]
    \item Ejecución de pruebas de carga con k6 (10K usuarios concurrentes) y primeros chaos experiments.
    \item Inicio de A/B testing de pricing en landing page (2--3 variantes).
    \item Publicación de primeros 2 artículos técnicos y envío de propuestas a conferencias.
    \item Configuración de CDNs multi-proveedor con failover automático.
    \item \textbf{Hito:} uptime 99.95\,\% validado, pricing en experimentación.
\end{itemize}

\textbf{Q2 2027 --- Profesionalización del Producto}
\begin{itemize}[leftmargin=1.5cm, label=\textbullet, itemsep=2pt, topsep=4pt]
    \item Optimización de PocketBase para escalado horizontal.
    \item Implementación de MLOps pipeline completo: Feast feature store, SHAP explainability, A/B testing.
    \item Lanzamiento de SDKs auto-generados en Python, JavaScript y Java.
    \item Dashboard de costos cloud con alertas de anomalía activas.
    \item \textbf{Hito:} SDKs publicados, BD migrada, primeros 100 clientes activos.
\end{itemize}

\textbf{Q3--Q4 2027 --- Crecimiento y Primeros Ingresos}
\begin{itemize}[leftmargin=1.5cm, label=\textbullet, itemsep=2pt, topsep=4pt]
    \item Implementación de data contracts con validación automática para todas las fuentes.
    \item Inicio del programa SOC 2 Tipo II: recolección de evidencias de controles.
    \item Lanzamiento de dashboards self-service para clientes Enterprise.
    \item Implementación de pipeline SAST/DAST en CI/CD.
    \item Implementación de encuestas NPS trimestrales automatizadas.
    \item Inicio de experimentos de Deep Learning (OCR para compliance, NLP avanzado).
    \item \textbf{Hito:} primeros \$100K ARR, NPS $>$ 30, controles SOC 2 en marcha.
\end{itemize}

\subsubsection{Fase 3 --- Scale (2028): Liderazgo Global}

\textbf{Q1--Q2 2028 --- Expansión Geográfica y Certificación}
\begin{itemize}[leftmargin=1.5cm, label=\textbullet, itemsep=2pt, topsep=4pt]
    \item Despliegue de infraestructura multi-región activo-activo (al menos 3 regiones cloud).
    \item Implementación de pipeline Speech-to-Text para feedback cualitativo (WER $\le$ 10\,\%).
    \item Obtención de certificación SOC 2 Tipo II (auditoría externa).
    \item Lanzamiento de marketplace de integraciones con SDK para partners.
    \item \textbf{Hito:} SOC 2 certificado, uptime $\ge$ 99.99\,\%, presencia multi-región.
\end{itemize}

\textbf{Q3--Q4 2028 --- Consolidación y Diversificación}
\begin{itemize}[leftmargin=1.5cm, label=\textbullet, itemsep=2pt, topsep=4pt]
    \item Inicio del roadmap de certificación ISO 27001 (gap analysis + plan de remediación).
    \item Implementación de auto-scaling predictivo con modelos ML.
    \item Expansión de cobertura de datos a métricas de sustentabilidad (emisiones CO$_2$ por ruta).
    \item Construcción de comunidad de desarrolladores con programa de embajadores.
    \item Evaluación de expansión a verticales adyacentes (aviación ejecutiva, drones, espacio aéreo urbano).
    \item \textbf{Hito:} $\ge$ 500 clientes activos, ARR duplicado anualmente, NPS enterprise $>$ 50.
\end{itemize}

\medskip
\textbf{Metas Intermedias por Fase}

\begin{table}[H]
\small
\centering
\begin{tabularx}{\linewidth}{>{\raggedright\arraybackslash}p{2.5cm} >{\raggedright\arraybackslash}p{2.5cm} >{\raggedright\arraybackslash}p{2.5cm} >{\raggedright\arraybackslash}p{2.5cm} >{\raggedright\arraybackslash}p{2.5cm} >{\raggedright\arraybackslash}X}
\toprule
\multicolumn{1}{>{\raggedright\arraybackslash}p{2.5cm}}{\textbf{Fase}} &
\multicolumn{1}{>{\raggedright\arraybackslash}p{2.5cm}}{\textbf{ARR}} &
\multicolumn{1}{>{\raggedright\arraybackslash}p{2.5cm}}{\textbf{NPS}} &
\multicolumn{1}{>{\raggedright\arraybackslash}p{2.5cm}}{\textbf{Uptime}} &
\multicolumn{1}{>{\raggedright\arraybackslash}p{2.5cm}}{\textbf{Churn (Ent.)}} &
\multicolumn{1}{>{\raggedright\arraybackslash}X}{\textbf{Clientes Activos}} \\
\midrule
Q3 2026 & Pre-revenue & --- & $\ge$ 99.5\,\% (staging) & --- & 0 (fundación) \\
\midrule
Q4 2026 & Pre-revenue & --- & $\ge$ 99.9\,\% & --- & $\ge$ 50 Beta Testers \\
\midrule
Q1--Q2 2027 & Primeros \$10K--\$30K MRR & $>$ 20 & $\ge$ 99.95\,\% & $<$ 3\,\% & $\ge$ 100 clientes \\
\midrule
Q3--Q4 2027 & \$100K ARR & $>$ 30 & $\ge$ 99.95\,\% & $<$ 2\,\% & $\ge$ 200 clientes \\
\midrule
Q1--Q2 2028 & \$300K--\$500K ARR & $>$ 40 & $\ge$ 99.99\,\% & $<$ 1\,\% & $\ge$ 350 clientes \\
\midrule
Q3--Q4 2028 & \$1M+ ARR & $>$ 50 (Ent.) & $\ge$ 99.99\,\% & $<$ 0.5\,\% & $\ge$ 500 clientes \\
\bottomrule
\end{tabularx}
\caption{Indicadores clave de progreso por fase del roadmap}
\end{table}

%---------------------------------------------------------------------
\subsection{Alineación Vertical: Estratégico \texorpdfstring{$\rightarrow$}{→} Táctico \texorpdfstring{$\rightarrow$}{→} Operativo}

\begin{table}[H]
\small
\centering
\begin{tabularx}{\linewidth}{>{\raggedright\arraybackslash}p{3.5cm} >{\raggedright\arraybackslash}p{3.5cm} >{\raggedright\arraybackslash}X}
\toprule
\multicolumn{1}{>{\raggedright\arraybackslash}p{3.5cm}}{\textbf{Objetivo Estratégico}} &
\multicolumn{1}{>{\raggedright\arraybackslash}p{3.5cm}}{\textbf{Objetivos Tácticos}} &
\multicolumn{1}{>{\raggedright\arraybackslash}X}{\textbf{Objetivos Operativos}} \\
\midrule
OE1 --- Growth Hacking & OT1, OT2, OT4, OT19 & OP1, OP2, OP25, OP30, OP31 \\
\midrule
OE2 --- APIs y SDKs & OT3, OT4, OT11 & OP24, OP25, OP26 \\
\midrule
OE3 --- Cloud HA & OT5, OT6, OT14, OT15, OT16 & OP8, OP9, OP15, OP16, OP17, OP18, OP19, OP27, OP28 \\
\midrule
OE4 --- BI Centralizada & OT7, OT8, OT15 & OP3, OP4, OP7 \\
\midrule
OE5 --- IA Aeronáutica & OT9, OT13, OT21, OT22 & OP5, OP6, OP32, OP33 \\
\midrule
OE6 --- Compliance & OT10, OT17 & OP10, OP11, OP21, OP22, OP29 \\
\midrule
OE7 --- Customer Success & OT2, OT11, OT19 & OP1, OP12, OP13, OP26, OP30 \\
\midrule
OE8 --- Talento Remote-First & OT12, OT13, OT20 & OP14, OP23, OP29, OP31 \\
\midrule
OE9 --- Calidad de Datos & OT7, OT15, OT18 & OP3, OP4, OP34 \\
\midrule
OE10 --- Gestión de Producto & OT18, OT19, OT20 & OP14, OP20, OP30, OP31 \\
\bottomrule
\end{tabularx}
\caption{Alineación vertical de objetivos: estratégico $\rightarrow$ táctico $\rightarrow$ operativo}
\end{table}

\medskip
\textbf{Correspondencia OP $\rightarrow$ Caso de Uso Operativo.} Los 34 Objetivos Operativos se cubren con 20 Casos de Uso Operativos; varios CUs implementan múltiples OPs relacionados. Cada OP tiene al menos un CU asignado para garantizar su ejecución y trazabilidad.

\small
\setlength{\LTleft}{0pt}
\setlength{\LTright}{0pt}
\begin{longtable}{>{\raggedright\arraybackslash}p{0.8cm} >{\raggedright\arraybackslash}p{2.2cm} >{\raggedright\arraybackslash}p{7.8cm} >{\raggedright\arraybackslash}p{4.8cm}}
\caption{Correspondencia entre los 34 Objetivos Operativos y los Casos de Uso}\\
\toprule
\textbf{OP} & \textbf{CU} & \textbf{Descripción} & \textbf{Justificación} \\
\midrule
\endfirsthead
\multicolumn{4}{c}{\small \textit{...continuación...}}\\
\toprule
\textbf{OP} & \textbf{CU} & \textbf{Descripción} & \textbf{Justificación} \\
\midrule
\endhead
\midrule
\endfoot
\bottomrule
\endlastfoot
OP1 & CU-O01 & Gestionar registro automatizado & Validación de email + generación de API Keys + aprovisionamiento de tenant \\
OP2 & CU-O02 & Monitorear rate limit y auto-scaling & Supervisión de thresholds, throttling graduado y elasticidad de pods \\
OP3 & CU-O03 & Supervisar ejecución de pipelines ETL & Monitoreo de duración, volumen, estado y re-ejecución de DAGs \\
OP4 & CU-O04 & Validar calidad de datos post-ingesta & Suite GE: nulos, duplicados, rangos, frescura y schema \\
OP5 & CU-O05 & Reentrenar modelos ML programados & Entrenamiento semanal, comparación MAPE, auto-promoción \\
OP6 & CU-O05 & Verificar drift de modelos en producción & Monitoreo PSI diario, alerta si $>$ 0.2, reentrenamiento automático \\
OP7 & CU-O06 & Refrescar dashboards BI & Refresh nocturno con verificación de integridad y consistencia \\
OP8 & CU-O07 & Monitorear telemetría de contenedores & Métricas de CPU/mem/red cada 30s, dashboards Grafana, respuesta a alertas \\
OP9 & CU-O07 & Configurar y mantener reglas de alerta & Thresholds por servicio, pruebas en staging, revisión de falsos positivos \\
OP10 & CU-O18 & Ejecutar escaneo SAST/DAST en CI/CD & SonarQube + OWASP ZAP en cada PR, bloqueo por vulnerabilidades críticas \\
OP11 & CU-O18 & Aplicar parches de seguridad críticos & Monitoreo de CVEs, parches $\le$ 24h, inventory de versiones de paquetes \\
OP12 & CU-O08 & Atender tickets de desarrolladores externos & Triaging NLP, primera respuesta $\le$ 2h, resolución e investigación \\
OP13 & CU-O14 & Documentar FAQs en base de conocimientos & Detección de temas recurrentes, borrador automático, publicación y medición \\
OP14 & CU-O15 & Gestionar sprints en Linear & Creación de tickets, standup diario, cierre con retrospectiva documentada \\
OP15 & CU-O19 & Revisar y optimizar costos cloud & Auditoría semanal, anomalías de gasto, recursos ociosos, recomendaciones \\
OP16 & CU-O09 & Ejecutar backup diario automatizado & Backup encriptado con KMS, upload a cold storage, log de integridad \\
OP17 & CU-O09 & Ejecutar prueba de restauración semanal & Restore en instancia aislada, validación de row counts y smoke tests \\
OP18 & CU-O10 & Calcular SLI/SLO y error budget & Cómputo de uptime 30 días, proyección de consumo, recomendación de congelación \\
OP19 & CU-O10 & Ejecutar congelación de deploys & Bloqueo de CI/CD automático, notificación a stakeholders, plan de remediación \\
OP20 & CU-O11 & Publicar changelog semanal SemVer & Recolección de commits, categorización, bump de versión, publicación en portal \\
OP21 & CU-O12 & Rotar credenciales de servicio & Generación en Vault, actualización en K8s, verificación con smoke tests \\
OP22 & CU-O12 & Auditar accesos IAM trimestralmente & Revisión de roles y políticas, revocación de accesos injustificados, reporte \\
OP23 & CU-O13 & Realizar post-mortem blameless $\le$ 72h & Timeline automático, 5 whys, documentación de aprendizajes, tickets de acción \\
OP24 & CU-O16 & Validar especificaciones OpenAPI en CI/CD & Spectral linter en cada PR, detección de breaking changes, bloqueo de merge \\
OP25 & CU-T02 & Monitorear presencia en marketplace RapidAPI & Uptime del listing, respuesta a reviews, actualización trimestral de planes \\
OP26 & CU-T07 & Mantener documentación interactiva del portal & Sincronización con specs, verificación de ejemplos, guías de migración \\
OP27 & CU-T09 & Ejecutar pruebas de carga semanales en staging & k6 con 10K usuarios, medición de latencia p95 y throughput, reporte \\
OP28 & CU-T09 & Ejecutar chaos engineering mensualmente & Litmus Chaos, kill pod, network partition, DB failover, medición MTTR \\
OP29 & CU-E04 & Mantener documentación legal y compliance & Actualización de ToS/DPA, evidencias SOC 2, registros de procesamiento GDPR \\
OP30 & CU-T10 & Recolectar y analizar feedback de clientes & Encuestas NPS + CSAT, transcripción STT, clasificación de sentimiento \\
OP31 & CU-T01 & Publicar contenido técnico regularmente & Artículos en blog, redes sociales, propuestas de charlas, medición de engagement \\
OP32 & CU-O20 & Ejecutar experimentos de Deep Learning OCR & Entrenamiento de modelos CRNN/TrOCR sobre documentos de compliance, registro en MLflow \\
OP33 & CU-O20 & Implementar pipeline de Speech-to-Text & Despliegue de Whisper, medición de WER, extracción de temas y sentimiento \\
OP34 & CU-O17 & Validar data contracts entre productores y consumidores & Schema contracts en Git, validación en CI, alerta de drift, bloqueo de ingesta \\
\end{longtable}

%=====================================================================
% APÉNDICE: CASOS DE USO DETALLADOS
%=====================================================================

\appendix
\newpage


\section{Casos de Uso Estratégicos}

\subsection{CU-E01: Consultar Tablero Balanced Scorecard}

\begin{table}[H]
\footnotesize
\begin{tabularx}{\linewidth}{>{\bfseries\raggedright\arraybackslash}p{2.5cm} >{\raggedright\arraybackslash}X}
\toprule
\multicolumn{2}{l}{\textbf{CU-E01: Consultar Tablero Balanced Scorecard} \hfill \textit{Prioridad: 10 \quad | \quad Estratégico}} \\
\midrule
Actor Principal & CEO \\
\midrule
Propósito & Monitorear en tiempo real el desempeño global contra las metas del BSC. \\
\midrule
Precondición & El CEO accede al sistema con autenticación 2FA. \\
\midrule
Flujo Principal & El sistema carga KPIs de Fact\_Suscripcion, Fact\_API\_Call, Fact\_Satisfaccion y Fact\_Incidente. Renderiza semáforos (verde/amarillo/rojo) contrastando el valor real contra la meta. El CEO puede hacer drill-down a la tendencia histórica de cada indicador. \\
\midrule
Postcondición & vista 360 del estado de la empresa. \\
\bottomrule
\end{tabularx}
\end{table}

\vspace{2pt}
\begin{table}[H]
\footnotesize
\begin{tabularx}{\linewidth}{>{\raggedright\arraybackslash}p{1.5cm} >{\raggedright\arraybackslash}p{2.2cm} >{\raggedright\arraybackslash}p{5.5cm} >{\raggedright\arraybackslash}X}
\toprule
\textbf{ID} & \textbf{Como...} & \textbf{Deseo...} & \textbf{Para...} \\
\midrule
HU-E01.1 & CEO & visualizar el ARR actual comparado contra la meta anual & evaluar el desempeño financiero \\
\midrule
HU-E01.2 & CEO & ver el NPS segmentado por región geográfica & identificar mercados con baja satisfacción \\
\midrule
HU-E01.3 & CEO & comparar el uptime real contra el SLO del 99.99\% & verificar la disponibilidad del servicio \\
\midrule
HU-E01.4 & CEO & recibir alertas automáticas cuando un KPI entre en zona roja & tomar acción correctiva inmediata \\
\midrule
HU-E01.5 & CEO & exportar el tablero BSC completo en formato PDF & presentarlo en reuniones de directorio \\
\bottomrule
\end{tabularx}
\end{table}

\vspace{4pt}
\noindent\rule{\linewidth}{0.4pt}
\vspace{4pt}


\subsection{CU-E02: Analizar ARR y Rentabilidad Mensual}

\begin{table}[H]
\footnotesize
\begin{tabularx}{\linewidth}{>{\bfseries\raggedright\arraybackslash}p{2.5cm} >{\raggedright\arraybackslash}X}
\toprule
\multicolumn{2}{l}{\textbf{CU-E02: Analizar ARR y Rentabilidad Mensual} \hfill \textit{Prioridad: 9 \quad | \quad Estratégico}} \\
\midrule
Actor Principal & CEO \\
\midrule
Propósito & Evaluar ingresos recurrentes, márgenes y tendencias de crecimiento. \\
\midrule
Precondición & El CEO o CFO selecciona un período de análisis. \\
\midrule
Flujo Principal & El sistema carga agregaciones de Fact\_Suscripcion y Fact\_API\_Call. Compara los resultados contra períodos anteriores. Detecta anomalías con un algoritmo de desviación estándar. Genera un reporte financiero detallado. \\
\midrule
Postcondición & diagnóstico financiero actualizado y disponible. \\
\bottomrule
\end{tabularx}
\end{table}

\vspace{2pt}
\begin{table}[H]
\footnotesize
\begin{tabularx}{\linewidth}{>{\raggedright\arraybackslash}p{1.5cm} >{\raggedright\arraybackslash}p{2.2cm} >{\raggedright\arraybackslash}p{5.5cm} >{\raggedright\arraybackslash}X}
\toprule
\textbf{ID} & \textbf{Como...} & \textbf{Deseo...} & \textbf{Para...} \\
\midrule
HU-E02.1 & CEO & visualizar el ARR desglosado por plan de suscripción & entender la composición de ingresos \\
\midrule
HU-E02.2 & CEO & comparar el crecimiento mes a mes del ARR & evaluar la tracción del negocio \\
\midrule
HU-E02.3 & CEO & detectar caídas inesperadas de revenue & investigar posibles causas de fuga \\
\midrule
HU-E02.4 & CEO & analizar la tasa de churn versus la expansión de cuentas existentes & medir la salud de la base de clientes \\
\midrule
HU-E02.5 & CEO & exportar el reporte financiero mensual & compartirlo con inversionistas y el board \\
\bottomrule
\end{tabularx}
\end{table}

\vspace{4pt}
\noindent\rule{\linewidth}{0.4pt}
\vspace{4pt}


\subsection{CU-E03: Evaluar Uptime Global y Cumplimiento SLA}

\begin{table}[H]
\footnotesize
\begin{tabularx}{\linewidth}{>{\bfseries\raggedright\arraybackslash}p{2.5cm} >{\raggedright\arraybackslash}X}
\toprule
\multicolumn{2}{l}{\textbf{CU-E03: Evaluar Uptime Global y Cumplimiento SLA} \hfill \textit{Prioridad: 10 \quad | \quad Estratégico}} \\
\midrule
Actor Principal & CTO \\
\midrule
Propósito & Verificar disponibilidad del sistema en tiempo real en todas las regiones. \\
\midrule
Precondición & El CTO selecciona una ventana temporal. \\
\midrule
Flujo Principal & El sistema carga datos de uptime por región desde Prometheus. Calcula el error budget consumido. Muestra el porcentaje de cumplimiento de SLA desglosado por endpoint. \\
\midrule
Postcondición & estado de infraestructura verificado y documentado. \\
\bottomrule
\end{tabularx}
\end{table}

\vspace{2pt}
\begin{table}[H]
\footnotesize
\begin{tabularx}{\linewidth}{>{\raggedright\arraybackslash}p{1.5cm} >{\raggedright\arraybackslash}p{2.2cm} >{\raggedright\arraybackslash}p{5.5cm} >{\raggedright\arraybackslash}X}
\toprule
\textbf{ID} & \textbf{Como...} & \textbf{Deseo...} & \textbf{Para...} \\
\midrule
HU-E03.1 & CTO & visualizar el uptime por región de despliegue & identificar zonas con degradación \\
\midrule
HU-E03.2 & CTO & comparar el uptime mensual contra el SLO del 99.99\% & verificar el cumplimiento contractual \\
\midrule
HU-E03.3 & CTO & recibir una alerta cuando exista una brecha de disponibilidad & activar el protocolo de respuesta \\
\midrule
HU-E03.4 & CTO & hacer drill-down hasta el nivel de endpoint & identificar la causa raíz de una caída \\
\midrule
HU-E03.5 & CTO & generar automáticamente un reporte de SLA por cliente & entregar transparencia a clientes Enterprise \\
\bottomrule
\end{tabularx}
\end{table}

\vspace{4pt}
\noindent\rule{\linewidth}{0.4pt}
\vspace{4pt}


\subsection{CU-E04: Revisar Cumplimiento Normativo}

\begin{table}[H]
\footnotesize
\begin{tabularx}{\linewidth}{>{\bfseries\raggedright\arraybackslash}p{2.5cm} >{\raggedright\arraybackslash}X}
\toprule
\multicolumn{2}{l}{\textbf{CU-E04: Revisar Cumplimiento Normativo} \hfill \textit{Prioridad: 8 \quad | \quad Estratégico}} \\
\midrule
Actor Principal & CTO / Auditor \\
\midrule
Propósito & Verificar estado de compliance contra marcos SOC 2, ISO 27001, GDPR, IATA. \\
\midrule
Precondición & El CTO o Auditor selecciona el marco normativo a evaluar. \\
\midrule
Flujo Principal & El sistema carga las evidencias de controles implementados. Muestra un score de cumplimiento por control. Identifica gaps con semáforo. Genera un reporte de auditoría exportable. \\
\midrule
Postcondición & diagnóstico de compliance actualizado. \\
\bottomrule
\end{tabularx}
\end{table}

\vspace{2pt}
\begin{table}[H]
\footnotesize
\begin{tabularx}{\linewidth}{>{\raggedright\arraybackslash}p{1.5cm} >{\raggedright\arraybackslash}p{2.2cm} >{\raggedright\arraybackslash}p{5.5cm} >{\raggedright\arraybackslash}X}
\toprule
\textbf{ID} & \textbf{Como...} & \textbf{Deseo...} & \textbf{Para...} \\
\midrule
HU-E04.1 & Auditor & ver el progreso del programa SOC 2 Tipo II & planificar la auditoría externa \\
\midrule
HU-E04.2 & CTO & verificar la residencia de datos por región & cumplir con las exigencias del GDPR \\
\midrule
HU-E04.3 & Auditor & comprobar que el cifrado en tránsito y en reposo está activo & validar el control técnico \\
\midrule
HU-E04.4 & CTO & identificar controles no conformes marcados en rojo & priorizar su remediación \\
\midrule
HU-E04.5 & Auditor & exportar el reporte de cumplimiento en formato PDF & adjuntarlo al expediente de certificación \\
\bottomrule
\end{tabularx}
\end{table}

\vspace{4pt}
\noindent\rule{\linewidth}{0.4pt}
\vspace{4pt}


\subsection{CU-E05: Definir Metas Estratégicas Trimestrales}

\begin{table}[H]
\footnotesize
\begin{tabularx}{\linewidth}{>{\bfseries\raggedright\arraybackslash}p{2.5cm} >{\raggedright\arraybackslash}X}
\toprule
\multicolumn{2}{l}{\textbf{CU-E05: Definir Metas Estratégicas Trimestrales} \hfill \textit{Prioridad: 9 \quad | \quad Estratégico}} \\
\midrule
Actor Principal & CEO / CTO \\
\midrule
Propósito & Establecer OKRs y metas del BSC para el próximo trimestre. \\
\midrule
Precondición & El equipo C-Level revisa los resultados del trimestre anterior. \\
\midrule
Flujo Principal & Proyecta tendencias con un módulo de simulación. Define targets cuantitativos por área. Asigna responsables. Publica las metas a toda la organización. \\
\midrule
Postcondición & metas trimestrales definidas y comunicadas. \\
\bottomrule
\end{tabularx}
\end{table}

\vspace{2pt}
\begin{table}[H]
\footnotesize
\begin{tabularx}{\linewidth}{>{\raggedright\arraybackslash}p{1.5cm} >{\raggedright\arraybackslash}p{2.2cm} >{\raggedright\arraybackslash}p{5.5cm} >{\raggedright\arraybackslash}X}
\toprule
\textbf{ID} & \textbf{Como...} & \textbf{Deseo...} & \textbf{Para...} \\
\midrule
HU-E05.1 & CEO & proyectar el ARR del próximo trimestre con simulación de escenarios & fijar metas realistas \\
\midrule
HU-E05.2 & CEO & definir OKRs por departamento alineados al BSC & asegurar cohesión estratégica \\
\midrule
HU-E05.3 & CTO & ajustar los targets técnicos según las tendencias de rendimiento & evitar metas inalcanzables \\
\midrule
HU-E05.4 & CEO & notificar automáticamente a cada responsable de área & iniciar la cascada táctica \\
\midrule
HU-E05.5 & CEO & hacer seguimiento semanal del progreso contra metas & corregir desviaciones a tiempo \\
\bottomrule
\end{tabularx}
\end{table}

\vspace{4pt}
\noindent\rule{\linewidth}{0.4pt}
\vspace{4pt}


\subsection{CU-E06: Analizar Retención de Talento y eNPS}

\begin{table}[H]
\footnotesize
\begin{tabularx}{\linewidth}{>{\bfseries\raggedright\arraybackslash}p{2.5cm} >{\raggedright\arraybackslash}X}
\toprule
\multicolumn{2}{l}{\textbf{CU-E06: Analizar Retención de Talento y eNPS} \hfill \textit{Prioridad: 7 \quad | \quad Estratégico}} \\
\midrule
Actor Principal & CEO / VP People \\
\midrule
Propósito & Evaluar satisfacción, retención y productividad del equipo. \\
\midrule
Precondición & El CEO carga los resultados de encuestas eNPS trimestrales. \\
\midrule
Flujo Principal & Segmenta por equipo, seniority y región. Compara contra el baseline histórico. Identifica áreas de riesgo de rotación. Recomienda acciones correctivas. \\
\midrule
Postcondición & diagnóstico de clima laboral disponible. \\
\bottomrule
\end{tabularx}
\end{table}

\vspace{2pt}
\begin{table}[H]
\footnotesize
\begin{tabularx}{\linewidth}{>{\raggedright\arraybackslash}p{1.5cm} >{\raggedright\arraybackslash}p{2.2cm} >{\raggedright\arraybackslash}p{5.5cm} >{\raggedright\arraybackslash}X}
\toprule
\textbf{ID} & \textbf{Como...} & \textbf{Deseo...} & \textbf{Para...} \\
\midrule
HU-E06.1 & CEO & visualizar la tendencia del eNPS en los últimos cuatro trimestres & evaluar la evolución del clima laboral \\
\midrule
HU-E06.2 & VP People & identificar equipos con eNPS en zona de riesgo & activar planes de retención \\
\midrule
HU-E06.3 & CEO & comparar la tasa de retención anual contra benchmarks de la industria & evaluar competitividad \\
\midrule
HU-E06.4 & VP People & analizar el time-to-productivity promedio por equipo & mejorar el proceso de onboarding \\
\midrule
HU-E06.5 & CEO & generar un reporte de People Analytics trimestral & presentarlo en la revisión de directorio \\
\bottomrule
\end{tabularx}
\end{table}

\vspace{4pt}
\noindent\rule{\linewidth}{0.4pt}
\vspace{4pt}

\newpage
\section{Casos de Uso Tácticos}

\subsection{CU-T01: Gestionar Campañas de Growth Hacking}

\begin{table}[H]
\footnotesize
\begin{tabularx}{\linewidth}{>{\bfseries\raggedright\arraybackslash}p{2.5cm} >{\raggedright\arraybackslash}X}
\toprule
\multicolumn{2}{l}{\textbf{CU-T01: Gestionar Campañas de Growth Hacking} \hfill \textit{Prioridad: 8 \quad | \quad Táctico}} \\
\midrule
Actor Principal & VP Marketing \\
\midrule
Propósito & Crear y automatizar campañas para captar Beta Testers. \\
\midrule
Precondición & El VP Marketing diseña una campaña en HubSpot. \\
\midrule
Flujo Principal & Segmenta la audiencia por industria y región. Programa envíos automatizados. Monitorea tasas de conversión en tiempo real. Optimiza con A/B testing de contenido. \\
\midrule
Postcondición & campañas activas monitoreadas. \\
\bottomrule
\end{tabularx}
\end{table}

\vspace{2pt}
\begin{table}[H]
\footnotesize
\begin{tabularx}{\linewidth}{>{\raggedright\arraybackslash}p{1.5cm} >{\raggedright\arraybackslash}p{2.2cm} >{\raggedright\arraybackslash}p{5.5cm} >{\raggedright\arraybackslash}X}
\toprule
\textbf{ID} & \textbf{Como...} & \textbf{Deseo...} & \textbf{Para...} \\
\midrule
HU-T01.1 & VP Marketing & crear una secuencia de emails automatizada en HubSpot & nurturar leads de manera escalable \\
\midrule
HU-T01.2 & VP Marketing & segmentar la audiencia por industria aeronáutica & personalizar el mensaje \\
\midrule
HU-T01.3 & VP Marketing & ejecutar un A/B test del asunto del email & maximizar la tasa de apertura \\
\midrule
HU-T01.4 & VP Marketing & medir la conversión desde email hasta registro en la plataforma & calcular el ROI de la campaña \\
\midrule
HU-T01.5 & VP Marketing & visualizar el ROI de cada campaña en un dashboard & asignar presupuesto a las más efectivas \\
\bottomrule
\end{tabularx}
\end{table}

\vspace{4pt}
\noindent\rule{\linewidth}{0.4pt}
\vspace{4pt}


\subsection{CU-T02: Configurar y Publicar API en RapidAPI}

\begin{table}[H]
\footnotesize
\begin{tabularx}{\linewidth}{>{\bfseries\raggedright\arraybackslash}p{2.5cm} >{\raggedright\arraybackslash}X}
\toprule
\multicolumn{2}{l}{\textbf{CU-T02: Configurar y Publicar API en RapidAPI} \hfill \textit{Prioridad: 9 \quad | \quad Táctico}} \\
\midrule
Actor Principal & CTO / Developer Advocate \\
\midrule
Propósito & Publicar API Beta en marketplaces con planes y rate limits. \\
\midrule
Precondición & El CTO valida la spec OpenAPI 3.1. \\
\midrule
Flujo Principal & Configura el listing en RapidAPI. Define los planes (Freemium/Pro/Enterprise) con sus respectivos rate limits. Configura autenticación OAuth2. Publica y monitorea la adopción. \\
\midrule
Postcondición & API disponible en marketplace. \\
\bottomrule
\end{tabularx}
\end{table}

\vspace{2pt}
\begin{table}[H]
\footnotesize
\begin{tabularx}{\linewidth}{>{\raggedright\arraybackslash}p{1.5cm} >{\raggedright\arraybackslash}p{2.2cm} >{\raggedright\arraybackslash}p{5.5cm} >{\raggedright\arraybackslash}X}
\toprule
\textbf{ID} & \textbf{Como...} & \textbf{Deseo...} & \textbf{Para...} \\
\midrule
HU-T02.1 & CTO & validar la especificación OpenAPI 3.1 antes de publicar & garantizar conformidad con el estándar \\
\midrule
HU-T02.2 & Developer Advocate & configurar rate limits diferenciados por plan & proteger la infraestructura \\
\midrule
HU-T02.3 & CTO & configurar OAuth2 con flujo client\_credentials & asegurar el acceso a la API \\
\midrule
HU-T02.4 & Developer Advocate & publicar la API en el marketplace de RapidAPI & alcanzar desarrolladores globales \\
\midrule
HU-T02.5 & VP Marketing & ver las analíticas de adopción desde el marketplace & medir el alcance orgánico \\
\bottomrule
\end{tabularx}
\end{table}

\vspace{4pt}
\noindent\rule{\linewidth}{0.4pt}
\vspace{4pt}


\subsection{CU-T03: Gestionar Infraestructura Cloud con Terraform}

\begin{table}[H]
\footnotesize
\begin{tabularx}{\linewidth}{>{\bfseries\raggedright\arraybackslash}p{2.5cm} >{\raggedright\arraybackslash}X}
\toprule
\multicolumn{2}{l}{\textbf{CU-T03: Gestionar Infraestructura Cloud con Terraform} \hfill \textit{Prioridad: 10 \quad | \quad Táctico}} \\
\midrule
Actor Principal & SRE / DevOps \\
\midrule
Propósito & Administrar todos los recursos cloud como código versionado. \\
\midrule
Precondición & El SRE escribe un módulo Terraform para el recurso requerido. \\
\midrule
Flujo Principal & Somete el cambio a PR review por el equipo. Ejecuta terraform plan y terraform apply. Verifica que los recursos se hayan creado correctamente. Versiona todo en Git. \\
\midrule
Postcondición & infraestructura 100\% como código. \\
\bottomrule
\end{tabularx}
\end{table}

\vspace{2pt}
\begin{table}[H]
\footnotesize
\begin{tabularx}{\linewidth}{>{\raggedright\arraybackslash}p{1.5cm} >{\raggedright\arraybackslash}p{2.2cm} >{\raggedright\arraybackslash}p{5.5cm} >{\raggedright\arraybackslash}X}
\toprule
\textbf{ID} & \textbf{Como...} & \textbf{Deseo...} & \textbf{Para...} \\
\midrule
HU-T03.1 & SRE & definir un módulo Terraform para el clúster K8s & replicar entornos idénticos \\
\midrule
HU-T03.2 & SRE & configurar la CDN multi-proveedor vía Terraform & gestionarla como código \\
\midrule
HU-T03.3 & SRE & gestionar los registros DNS como código & evitar configuraciones manuales propensas a error \\
\midrule
HU-T03.4 & SRE & versionar todos los cambios de infraestructura en Git & tener trazabilidad completa \\
\midrule
HU-T03.5 & SRE & ejecutar rollback de infraestructura vía git revert & recuperar el último estado estable \\
\bottomrule
\end{tabularx}
\end{table}

\vspace{4pt}
\noindent\rule{\linewidth}{0.4pt}
\vspace{4pt}


\subsection{CU-T04: Monitorear Pipelines ETL y Calidad de Datos}

\begin{table}[H]
\footnotesize
\begin{tabularx}{\linewidth}{>{\bfseries\raggedright\arraybackslash}p{2.5cm} >{\raggedright\arraybackslash}X}
\toprule
\multicolumn{2}{l}{\textbf{CU-T04: Monitorear Pipelines ETL y Calidad de Datos} \hfill \textit{Prioridad: 9 \quad | \quad Táctico}} \\
\midrule
Actor Principal & Data Engineer \\
\midrule
Propósito & Asegurar que los flujos de datos estén saludables y cumplan estándares. \\
\midrule
Precondición & El Data Engineer ejecuta la suite de Great Expectations programada. \\
\midrule
Flujo Principal & Verifica reglas de frescura, volumen, nulos y schema. Genera un reporte de calidad consolidado. Recibe alertas en Slack ante cualquier fallo. \\
\midrule
Postcondición & calidad de datos verificada y documentada. \\
\bottomrule
\end{tabularx}
\end{table}

\vspace{2pt}
\begin{table}[H]
\footnotesize
\begin{tabularx}{\linewidth}{>{\raggedright\arraybackslash}p{1.5cm} >{\raggedright\arraybackslash}p{2.2cm} >{\raggedright\arraybackslash}p{5.5cm} >{\raggedright\arraybackslash}X}
\toprule
\textbf{ID} & \textbf{Como...} & \textbf{Deseo...} & \textbf{Para...} \\
\midrule
HU-T04.1 & Data Engineer & visualizar el estado de ejecución de cada pipeline ETL & detectar fallos rápidamente \\
\midrule
HU-T04.2 & Data Engineer & verificar la frescura de los datos por fuente & asegurar que no haya rezagos \\
\midrule
HU-T04.3 & Data Engineer & validar la integridad del schema contra el contrato de datos & prevenir errores en downstream \\
\midrule
HU-T04.4 & Data Engineer & recibir una alerta inmediata ante un fallo de validación & iniciar la corrección \\
\midrule
HU-T04.5 & CTO & ver un scorecard mensual de calidad de datos & monitorear la salud del data pipeline \\
\bottomrule
\end{tabularx}
\end{table}

\vspace{4pt}
\noindent\rule{\linewidth}{0.4pt}
\vspace{4pt}


\subsection{CU-T05: Entrenar y Versionar Modelos ML}

\begin{table}[H]
\footnotesize
\begin{tabularx}{\linewidth}{>{\bfseries\raggedright\arraybackslash}p{2.5cm} >{\raggedright\arraybackslash}X}
\toprule
\multicolumn{2}{l}{\textbf{CU-T05: Entrenar y Versionar Modelos ML} \hfill \textit{Prioridad: 9 \quad | \quad Táctico}} \\
\midrule
Actor Principal & ML Engineer \\
\midrule
Propósito & Gestionar ciclo de vida completo de modelos ML. \\
\midrule
Precondición & El ML Engineer crea un experimento en MLflow. \\
\midrule
Flujo Principal & Registra hiperparámetros y métricas. Compara versiones. Registra el modelo campeón. Promueve a staging. Valida comportamiento y despliega a producción. \\
\midrule
Postcondición & modelo versionado en producción. \\
\bottomrule
\end{tabularx}
\end{table}

\vspace{2pt}
\begin{table}[H]
\footnotesize
\begin{tabularx}{\linewidth}{>{\raggedright\arraybackslash}p{1.5cm} >{\raggedright\arraybackslash}p{2.2cm} >{\raggedright\arraybackslash}p{5.5cm} >{\raggedright\arraybackslash}X}
\toprule
\textbf{ID} & \textbf{Como...} & \textbf{Deseo...} & \textbf{Para...} \\
\midrule
HU-T05.1 & ML Engineer & loguear cada experimento con sus hiperparámetros en MLflow & mantener trazabilidad \\
\midrule
HU-T05.2 & ML Engineer & comparar dos versiones de modelo visualmente & seleccionar la de mejor desempeño \\
\midrule
HU-T05.3 & ML Engineer & registrar el champion model en el Model Registry & gestionar su ciclo de vida \\
\midrule
HU-T05.4 & ML Engineer & promover un modelo a producción vía API de MLflow & automatizar el deployment \\
\midrule
HU-T05.5 & ML Engineer & ejecutar rollback a la versión anterior del modelo & revertir una degradación \\
\bottomrule
\end{tabularx}
\end{table}

\vspace{4pt}
\noindent\rule{\linewidth}{0.4pt}
\vspace{4pt}


\subsection{CU-T06: Configurar Alertas de Seguridad y Compliance}

\begin{table}[H]
\footnotesize
\begin{tabularx}{\linewidth}{>{\bfseries\raggedright\arraybackslash}p{2.5cm} >{\raggedright\arraybackslash}X}
\toprule
\multicolumn{2}{l}{\textbf{CU-T06: Configurar Alertas de Seguridad y Compliance} \hfill \textit{Prioridad: 9 \quad | \quad Táctico}} \\
\midrule
Actor Principal & SRE / DevOps \\
\midrule
Propósito & Definir monitoreo de seguridad proactivo. \\
\midrule
Precondición & El SRE define umbrales de alerta en el sistema de monitoreo. \\
\midrule
Flujo Principal & Integra escaneos SAST/DAST en el pipeline de CI/CD. Configura notificaciones a PagerDuty. Prueba el disparo de alertas en staging. Activa el monitoreo en producción. \\
\midrule
Postcondición & seguridad monitoreada 24/7. \\
\bottomrule
\end{tabularx}
\end{table}

\vspace{2pt}
\begin{table}[H]
\footnotesize
\begin{tabularx}{\linewidth}{>{\raggedright\arraybackslash}p{1.5cm} >{\raggedright\arraybackslash}p{2.2cm} >{\raggedright\arraybackslash}p{5.5cm} >{\raggedright\arraybackslash}X}
\toprule
\textbf{ID} & \textbf{Como...} & \textbf{Deseo...} & \textbf{Para...} \\
\midrule
HU-T06.1 & SRE & definir una alerta para vulnerabilidades críticas con SLA de parcheo < 24h & cumplir la política de seguridad \\
\midrule
HU-T06.2 & SRE & configurar una alerta de rate limit breach & detectar abusos o ataques \\
\midrule
HU-T06.3 & SRE & configurar una alerta de expiración de certificados SSL & renovarlos antes del vencimiento \\
\midrule
HU-T06.4 & SRE & configurar una alerta de acceso no autorizado & activar respuesta ante intrusiones \\
\midrule
HU-T06.5 & SRE & probar la entrega de alertas a PagerDuty con un incidente simulado & validar el flujo completo \\
\bottomrule
\end{tabularx}
\end{table}

\vspace{4pt}
\noindent\rule{\linewidth}{0.4pt}
\vspace{4pt}


\subsection{CU-T07: Mantener Developer Portal y SDKs}

\begin{table}[H]
\footnotesize
\begin{tabularx}{\linewidth}{>{\bfseries\raggedright\arraybackslash}p{2.5cm} >{\raggedright\arraybackslash}X}
\toprule
\multicolumn{2}{l}{\textbf{CU-T07: Mantener Developer Portal y SDKs} \hfill \textit{Prioridad: 8 \quad | \quad Táctico}} \\
\midrule
Actor Principal & Developer Advocate \\
\midrule
Propósito & Mantener documentación, SDKs y sandbox actualizados. \\
\midrule
Precondición & El Developer Advocate genera documentación interactiva desde la spec OpenAPI. \\
\midrule
Flujo Principal & Publica SDKs auto-generados para Python, JavaScript y Java. Actualiza el changelog. Monitorea analíticas de páginas. Mantiene el sandbox con datos sintéticos. \\
\midrule
Postcondición & portal actualizado y operativo. \\
\bottomrule
\end{tabularx}
\end{table}

\vspace{2pt}
\begin{table}[H]
\footnotesize
\begin{tabularx}{\linewidth}{>{\raggedright\arraybackslash}p{1.5cm} >{\raggedright\arraybackslash}p{2.2cm} >{\raggedright\arraybackslash}p{5.5cm} >{\raggedright\arraybackslash}X}
\toprule
\textbf{ID} & \textbf{Como...} & \textbf{Deseo...} & \textbf{Para...} \\
\midrule
HU-T07.1 & Developer Advocate & generar documentación interactiva desde la spec OpenAPI & mantenerla sincronizada con el código \\
\midrule
HU-T07.2 & Developer Advocate & publicar una nueva versión del SDK de Python & facilitar la integración de los clientes \\
\midrule
HU-T07.3 & Developer Advocate & ver las analíticas de tráfico en las páginas de documentación & identificar endpoints populares \\
\midrule
HU-T07.4 & Developer Advocate & gestionar el sandbox con datos sintéticos renovados & ofrecer una experiencia de prueba realista \\
\midrule
HU-T07.5 & Developer Advocate & sincronizar automáticamente el changelog con cada release & mantener informada a la comunidad \\
\bottomrule
\end{tabularx}
\end{table}

\vspace{4pt}
\noindent\rule{\linewidth}{0.4pt}
\vspace{4pt}


\subsection{CU-T08: Analizar Costos Cloud y Optimizar Recursos}

\begin{table}[H]
\footnotesize
\begin{tabularx}{\linewidth}{>{\bfseries\raggedright\arraybackslash}p{2.5cm} >{\raggedright\arraybackslash}X}
\toprule
\multicolumn{2}{l}{\textbf{CU-T08: Analizar Costos Cloud y Optimizar Recursos} \hfill \textit{Prioridad: 8 \quad | \quad Táctico}} \\
\midrule
Actor Principal & CTO / SRE \\
\midrule
Propósito & Monitorear y optimizar gasto cloud. \\
\midrule
Precondición & El CTO revisa el dashboard de costos por servicio cloud. \\
\midrule
Flujo Principal & Compara el gasto real contra el presupuesto. Detecta anomalías de gasto con detección de outliers. Identifica recursos subutilizados. Genera recomendaciones de optimización. \\
\midrule
Postcondición & gasto cloud optimizado. \\
\bottomrule
\end{tabularx}
\end{table}

\vspace{2pt}
\begin{table}[H]
\footnotesize
\begin{tabularx}{\linewidth}{>{\raggedright\arraybackslash}p{1.5cm} >{\raggedright\arraybackslash}p{2.2cm} >{\raggedright\arraybackslash}p{5.5cm} >{\raggedright\arraybackslash}X}
\toprule
\textbf{ID} & \textbf{Como...} & \textbf{Deseo...} & \textbf{Para...} \\
\midrule
HU-T08.1 & CTO & visualizar los costos desglosados por servicio cloud & identificar los mayores drivers de gasto \\
\midrule
HU-T08.2 & SRE & detectar una anomalía de gasto diario & investigar posibles fugas \\
\midrule
HU-T08.3 & CTO & comparar la tendencia mensual de costos contra el presupuesto & controlar desviaciones \\
\midrule
HU-T08.4 & CTO & identificar recursos cloud ociosos & eliminarlos o redimensionarlos \\
\midrule
HU-T08.5 & SRE & generar recomendaciones automáticas de ahorro & reducir el costo mensual \\
\bottomrule
\end{tabularx}
\end{table}

\vspace{4pt}
\noindent\rule{\linewidth}{0.4pt}
\vspace{4pt}


\subsection{CU-T09: Ejecutar Pruebas de Carga y Chaos Engineering}

\begin{table}[H]
\footnotesize
\begin{tabularx}{\linewidth}{>{\bfseries\raggedright\arraybackslash}p{2.5cm} >{\raggedright\arraybackslash}X}
\toprule
\multicolumn{2}{l}{\textbf{CU-T09: Ejecutar Pruebas de Carga y Chaos Engineering} \hfill \textit{Prioridad: 8 \quad | \quad Táctico}} \\
\midrule
Actor Principal & SRE / DevOps \\
\midrule
Propósito & Validar resiliencia de la plataforma bajo estrés. \\
\midrule
Precondición & El SRE define escenarios de prueba en k6. \\
\midrule
Flujo Principal & Ejecuta la simulación con 10K usuarios concurrentes. Inyecta caos controlado (matar un pod aleatorio). Mide el tiempo de recuperación. Documenta hallazgos en un reporte. \\
\midrule
Postcondición & resiliencia validada. \\
\bottomrule
\end{tabularx}
\end{table}

\vspace{2pt}
\begin{table}[H]
\footnotesize
\begin{tabularx}{\linewidth}{>{\raggedright\arraybackslash}p{1.5cm} >{\raggedright\arraybackslash}p{2.2cm} >{\raggedright\arraybackslash}p{5.5cm} >{\raggedright\arraybackslash}X}
\toprule
\textbf{ID} & \textbf{Como...} & \textbf{Deseo...} & \textbf{Para...} \\
\midrule
HU-T09.1 & SRE & ejecutar k6 con un escenario de 10K usuarios concurrentes & validar la capacidad del sistema \\
\midrule
HU-T09.2 & SRE & ejecutar un experimento de chaos engineering & verificar que el sistema se auto-recupera \\
\midrule
HU-T09.3 & SRE & medir el tiempo de recuperación tras la inyección de fallo & compararlo contra el RTO \\
\midrule
HU-T09.4 & SRE & validar que el auto-scaling responde en menos de 2 minutos & garantizar elasticidad \\
\midrule
HU-T09.5 & CTO & revisar el reporte de resiliencia trimestral & tomar decisiones de inversión en infraestructura \\
\bottomrule
\end{tabularx}
\end{table}

\vspace{4pt}
\noindent\rule{\linewidth}{0.4pt}
\vspace{4pt}


\subsection{CU-T10: Validar Estrategia de Pricing con A/B Testing}

\begin{table}[H]
\footnotesize
\begin{tabularx}{\linewidth}{>{\bfseries\raggedright\arraybackslash}p{2.5cm} >{\raggedright\arraybackslash}X}
\toprule
\multicolumn{2}{l}{\textbf{CU-T10: Validar Estrategia de Pricing con A/B Testing} \hfill \textit{Prioridad: 8 \quad | \quad Táctico}} \\
\midrule
Actor Principal & VP Marketing / CEO \\
\midrule
Propósito & Probar hipótesis de precios antes del lanzamiento. \\
\midrule
Precondición & El VP Marketing define una hipótesis de pricing. \\
\midrule
Flujo Principal & Configura un A/B test en la landing page. Dirige tráfico 50/50 a cada variante. Mide la conversión. Analiza los resultados estadísticamente. Selecciona el pricing óptimo. \\
\midrule
Postcondición & pricing validado con datos. \\
\bottomrule
\end{tabularx}
\end{table}

\vspace{2pt}
\begin{table}[H]
\footnotesize
\begin{tabularx}{\linewidth}{>{\raggedright\arraybackslash}p{1.5cm} >{\raggedright\arraybackslash}p{2.2cm} >{\raggedright\arraybackslash}p{5.5cm} >{\raggedright\arraybackslash}X}
\toprule
\textbf{ID} & \textbf{Como...} & \textbf{Deseo...} & \textbf{Para...} \\
\midrule
HU-T10.1 & VP Marketing & configurar un A/B test en la landing page & comparar dos estructuras de precios \\
\midrule
HU-T10.2 & VP Marketing & comparar las tasas de conversión entre variantes & determinar cuál es más efectiva \\
\midrule
HU-T10.3 & CEO & medir la willingness-to-pay de los prospectos & definir el precio óptimo \\
\midrule
HU-T10.4 & VP Marketing & analizar la tasa de churn por nivel de precio & evitar precios que ahuyenten clientes \\
\midrule
HU-T10.5 & CEO & seleccionar el pricing que maximice el ratio LTV/CAC & asegurar rentabilidad a largo plazo \\
\bottomrule
\end{tabularx}
\end{table}

\vspace{4pt}
\noindent\rule{\linewidth}{0.4pt}
\vspace{4pt}

\newpage
\section{Casos de Uso Operativos}

\subsection{CU-O01: Registrar Tenant y Autogenerar API Keys}

\begin{table}[H]
\footnotesize
\begin{tabularx}{\linewidth}{>{\bfseries\raggedright\arraybackslash}p{2.5cm} >{\raggedright\arraybackslash}X}
\toprule
\multicolumn{2}{l}{\textbf{CU-O01: Registrar Tenant y Autogenerar API Keys} \hfill \textit{Prioridad: 10 \quad | \quad Operativo}} \\
\midrule
Actor Principal & Cliente B2B / Sistema \\
\midrule
Propósito & Permitir registro self-service de nuevos clientes. \\
\midrule
Precondición & El cliente completa un formulario de registro con email y datos de empresa. \\
\midrule
Flujo Principal & El sistema valida que el email sea único. Genera una API Key automáticamente. Crea un tenant con plan Freemium. Envía un email de bienvenida con enlaces a documentación. \\
\midrule
Postcondición & tenant activo. \\
\bottomrule
\end{tabularx}
\end{table}

\vspace{2pt}
\begin{table}[H]
\footnotesize
\begin{tabularx}{\linewidth}{>{\raggedright\arraybackslash}p{1.5cm} >{\raggedright\arraybackslash}p{2.2cm} >{\raggedright\arraybackslash}p{5.5cm} >{\raggedright\arraybackslash}X}
\toprule
\textbf{ID} & \textbf{Como...} & \textbf{Deseo...} & \textbf{Para...} \\
\midrule
HU-O01.1 & Cliente B2B & registrarme con mi email corporativo & crear una cuenta rápidamente \\
\midrule
HU-O01.2 & Sistema & autogenerar una API Key única por tenant & habilitar el acceso inmediato \\
\midrule
HU-O01.3 & Cliente B2B & verificar mi email mediante un link de confirmación & activar mi cuenta \\
\midrule
HU-O01.4 & Cliente B2B & recibir un email de bienvenida con guías rápidas & empezar a usar la API \\
\midrule
HU-O01.5 & Sistema & provisionar automáticamente un sandbox para el nuevo tenant & reducir el tiempo hasta el primer insight \\
\bottomrule
\end{tabularx}
\end{table}

\vspace{4pt}
\noindent\rule{\linewidth}{0.4pt}
\vspace{4pt}


\subsection{CU-O02: Consultar Datos de Vuelo vía API REST/GraphQL}

\begin{table}[H]
\footnotesize
\begin{tabularx}{\linewidth}{>{\bfseries\raggedright\arraybackslash}p{2.5cm} >{\raggedright\arraybackslash}X}
\toprule
\multicolumn{2}{l}{\textbf{CU-O02: Consultar Datos de Vuelo vía API REST/GraphQL} \hfill \textit{Prioridad: 10 \quad | \quad Operativo}} \\
\midrule
Actor Principal & Cliente B2B \\
\midrule
Propósito & Consultar datos aeronáuticos mediante API. \\
\midrule
Precondición & El cliente envía un request HTTP con API Key en header. \\
\midrule
Flujo Principal & El sistema autentica vía OAuth2. Valida parámetros y rate limit del plan. Ejecuta la query contra el Data Warehouse. Retorna respuesta en formato JSON. Loggea el consumo para facturación. \\
\midrule
Postcondición & datos entregados. \\
\bottomrule
\end{tabularx}
\end{table}

\vspace{2pt}
\begin{table}[H]
\footnotesize
\begin{tabularx}{\linewidth}{>{\raggedright\arraybackslash}p{1.5cm} >{\raggedright\arraybackslash}p{2.2cm} >{\raggedright\arraybackslash}p{5.5cm} >{\raggedright\arraybackslash}X}
\toprule
\textbf{ID} & \textbf{Como...} & \textbf{Deseo...} & \textbf{Para...} \\
\midrule
HU-O02.1 & Cliente B2B & consultar vuelos filtrando por aerolínea y rango de fechas & obtener datos específicos \\
\midrule
HU-O02.2 & Cliente B2B & filtrar vuelos con retraso mayor a 30 minutos & analizar patrones de impuntualidad \\
\midrule
HU-O02.3 & Cliente B2B & paginar resultados grandes con cursores & manejar eficientemente grandes volúmenes \\
\midrule
HU-O02.4 & Cliente B2B & consultar vía GraphQL con campos anidados & reducir el número de llamadas \\
\midrule
HU-O02.5 & Cliente B2B & recibir headers X-RateLimit-Remaining en cada respuesta & gestionar mi cuota de consumo \\
\bottomrule
\end{tabularx}
\end{table}

\vspace{4pt}
\noindent\rule{\linewidth}{0.4pt}
\vspace{4pt}


\subsection{CU-O03: Ejecutar Ingesta Diaria de Datos Aeronáuticos}

\begin{table}[H]
\footnotesize
\begin{tabularx}{\linewidth}{>{\bfseries\raggedright\arraybackslash}p{2.5cm} >{\raggedright\arraybackslash}X}
\toprule
\multicolumn{2}{l}{\textbf{CU-O03: Ejecutar Ingesta Diaria de Datos Aeronáuticos} \hfill \textit{Prioridad: 10 \quad | \quad Operativo}} \\
\midrule
Actor Principal & Data Engineer / Sistema \\
\midrule
Propósito & Ejecutar ingesta programada de datos desde proveedores. \\
\midrule
Precondición & El sistema descarga archivos desde el SFTP del proveedor de datos. \\
\midrule
Flujo Principal & Valida el schema CSV/Parquet. Transforma timestamps a UTC. Carga los datos a la tabla staging. Verifica que el row count coincida con lo esperado. Loggea el resultado de la ingesta. \\
\midrule
Postcondición & datos en staging. \\
\bottomrule
\end{tabularx}
\end{table}

\vspace{2pt}
\begin{table}[H]
\footnotesize
\begin{tabularx}{\linewidth}{>{\raggedright\arraybackslash}p{1.5cm} >{\raggedright\arraybackslash}p{2.2cm} >{\raggedright\arraybackslash}p{5.5cm} >{\raggedright\arraybackslash}X}
\toprule
\textbf{ID} & \textbf{Como...} & \textbf{Deseo...} & \textbf{Para...} \\
\midrule
HU-O03.1 & Data Engineer & descargar los datos del proveedor vía SFTP programado & automatizar la ingesta \\
\midrule
HU-O03.2 & Data Engineer & validar el schema de los archivos antes de cargar & evitar corrupción en staging \\
\midrule
HU-O03.3 & Sistema & transformar todas las zonas horarias a UTC & estandarizar los timestamps \\
\midrule
HU-O03.4 & Sistema & cargar los datos validados a la tabla de staging & dejarlos disponibles para ETL \\
\midrule
HU-O03.5 & Data Engineer & verificar que el conteo de filas coincida con lo esperado & detectar pérdida de datos \\
\bottomrule
\end{tabularx}
\end{table}

\vspace{4pt}
\noindent\rule{\linewidth}{0.4pt}
\vspace{4pt}


\subsection{CU-O04: Validar Calidad de Datos en Pipeline ETL}

\begin{table}[H]
\footnotesize
\begin{tabularx}{\linewidth}{>{\bfseries\raggedright\arraybackslash}p{2.5cm} >{\raggedright\arraybackslash}X}
\toprule
\multicolumn{2}{l}{\textbf{CU-O04: Validar Calidad de Datos en Pipeline ETL} \hfill \textit{Prioridad: 10 \quad | \quad Operativo}} \\
\midrule
Actor Principal & Data Engineer / Sistema \\
\midrule
Propósito & Ejecutar controles de calidad tras la ingesta. \\
\midrule
Precondición & El sistema ejecuta la suite de Great Expectations sobre los datos en staging. \\
\midrule
Flujo Principal & Verifica nulos en columnas críticas. Detecta duplicados. Valida rangos plausibles de valores. Controla la frescura de los datos. Genera un reporte pass/fail. \\
\midrule
Postcondición & calidad verificada. \\
\bottomrule
\end{tabularx}
\end{table}

\vspace{2pt}
\begin{table}[H]
\footnotesize
\begin{tabularx}{\linewidth}{>{\raggedright\arraybackslash}p{1.5cm} >{\raggedright\arraybackslash}p{2.2cm} >{\raggedright\arraybackslash}p{5.5cm} >{\raggedright\arraybackslash}X}
\toprule
\textbf{ID} & \textbf{Como...} & \textbf{Deseo...} & \textbf{Para...} \\
\midrule
HU-O04.1 & Data Engineer & verificar que no haya nulos en flight\_id y departure\_time & garantizar integridad referencial \\
\midrule
HU-O04.2 & Sistema & detectar vuelos duplicados por flight\_id y fecha & prevenir métricas infladas \\
\midrule
HU-O04.3 & Data Engineer & validar que los valores de retraso estén en un rango plausible & detectar outliers \\
\midrule
HU-O04.4 & Sistema & verificar que la frescura de los datos sea menor a 5 minutos & cumplir el SLO de calidad \\
\midrule
HU-O04.5 & CTO & visualizar un dashboard pass/fail de calidad por fuente & monitorear la salud del pipeline \\
\bottomrule
\end{tabularx}
\end{table}

\vspace{4pt}
\noindent\rule{\linewidth}{0.4pt}
\vspace{4pt}


\subsection{CU-O05: Reentrenar Modelo de Predicción de Retrasos}

\begin{table}[H]
\footnotesize
\begin{tabularx}{\linewidth}{>{\bfseries\raggedright\arraybackslash}p{2.5cm} >{\raggedright\arraybackslash}X}
\toprule
\multicolumn{2}{l}{\textbf{CU-O05: Reentrenar Modelo de Predicción de Retrasos} \hfill \textit{Prioridad: 9 \quad | \quad Operativo}} \\
\midrule
Actor Principal & ML Engineer \\
\midrule
Propósito & Reentrenar modelo con datos frescos. \\
\midrule
Precondición & El sistema obtiene features desde Feast. \\
\midrule
Flujo Principal & Entrena un modelo XGBoost con los nuevos datos. Loggea métricas en MLflow. Compara el MAPE contra la versión anterior. Auto-promueve si mejora. Alerta al ML Engineer si degrada. \\
\midrule
Postcondición & modelo actualizado o alerta emitida. \\
\bottomrule
\end{tabularx}
\end{table}

\vspace{2pt}
\begin{table}[H]
\footnotesize
\begin{tabularx}{\linewidth}{>{\raggedright\arraybackslash}p{1.5cm} >{\raggedright\arraybackslash}p{2.2cm} >{\raggedright\arraybackslash}p{5.5cm} >{\raggedright\arraybackslash}X}
\toprule
\textbf{ID} & \textbf{Como...} & \textbf{Deseo...} & \textbf{Para...} \\
\midrule
HU-O05.1 & ML Engineer & obtener las features desde Feast automáticamente & evitar duplicación de lógica \\
\midrule
HU-O05.2 & Sistema & entrenar el modelo XGBoost con los datos frescos & mantenerlo actualizado \\
\midrule
HU-O05.3 & ML Engineer & comparar el MAPE del nuevo modelo contra el anterior & decidir si vale la pena promoverlo \\
\midrule
HU-O05.4 & Sistema & auto-promover el modelo si supera el umbral de mejora & acelerar el ciclo de deploy \\
\midrule
HU-O05.5 & ML Engineer & recibir una alerta si el modelo degrada por drift & investigar y corregir \\
\bottomrule
\end{tabularx}
\end{table}

\vspace{4pt}
\noindent\rule{\linewidth}{0.4pt}
\vspace{4pt}


\subsection{CU-O06: Refrescar Dashboards BI para Clientes}

\begin{table}[H]
\footnotesize
\begin{tabularx}{\linewidth}{>{\bfseries\raggedright\arraybackslash}p{2.5cm} >{\raggedright\arraybackslash}X}
\toprule
\multicolumn{2}{l}{\textbf{CU-O06: Refrescar Dashboards BI para Clientes} \hfill \textit{Prioridad: 9 \quad | \quad Operativo}} \\
\midrule
Actor Principal & Data Engineer / Sistema \\
\midrule
Propósito & Actualizar dashboards self-service con datos recientes. \\
\midrule
Precondición & El sistema dispara un refresh programado nocturno. \\
\midrule
Flujo Principal & Ejecuta agregaciones precalculadas. Actualiza la caché de dashboards. Verifica el renderizado correcto de visualizaciones. Loggea el tiempo total de refresh. \\
\midrule
Postcondición & dashboards actualizados. \\
\bottomrule
\end{tabularx}
\end{table}

\vspace{2pt}
\begin{table}[H]
\footnotesize
\begin{tabularx}{\linewidth}{>{\raggedright\arraybackslash}p{1.5cm} >{\raggedright\arraybackslash}p{2.2cm} >{\raggedright\arraybackslash}p{5.5cm} >{\raggedright\arraybackslash}X}
\toprule
\textbf{ID} & \textbf{Como...} & \textbf{Deseo...} & \textbf{Para...} \\
\midrule
HU-O06.1 & Data Engineer & programar el refresh nocturno de dashboards & que los clientes vean datos frescos al iniciar \\
\midrule
HU-O06.2 & Sistema & verificar que todos los KPIs se hayan cargado correctamente & garantizar la integridad visual \\
\midrule
HU-O06.3 & Data Engineer & validar que los datos del dashboard coincidan con la fuente & asegurar consistencia \\
\midrule
HU-O06.4 & Sistema & trackear la duración del refresh en cada ejecución & detectar degradaciones \\
\midrule
HU-O06.5 & SRE & recibir una alerta si el refresh de dashboards falla & reiniciar el proceso \\
\bottomrule
\end{tabularx}
\end{table}

\vspace{4pt}
\noindent\rule{\linewidth}{0.4pt}
\vspace{4pt}


\subsection{CU-O07: Monitorear Telemetría de Contenedores K8s}

\begin{table}[H]
\footnotesize
\begin{tabularx}{\linewidth}{>{\bfseries\raggedright\arraybackslash}p{2.5cm} >{\raggedright\arraybackslash}X}
\toprule
\multicolumn{2}{l}{\textbf{CU-O07: Monitorear Telemetría de Contenedores K8s} \hfill \textit{Prioridad: 10 \quad | \quad Operativo}} \\
\midrule
Actor Principal & SRE / DevOps \\
\midrule
Propósito & Monitoreo en tiempo real de infraestructura. \\
\midrule
Precondición & Prometheus recolecta métricas de CPU, memoria y red de los pods. \\
\midrule
Flujo Principal & Evalúa reglas de alerta cada 30 segundos. Dispara notificación a PagerDuty si un umbral es superado. Loggea el incidente automáticamente. El SRE ejecuta el runbook correspondiente. \\
\midrule
Postcondición & infraestructura monitoreada. \\
\bottomrule
\end{tabularx}
\end{table}

\vspace{2pt}
\begin{table}[H]
\footnotesize
\begin{tabularx}{\linewidth}{>{\raggedright\arraybackslash}p{1.5cm} >{\raggedright\arraybackslash}p{2.2cm} >{\raggedright\arraybackslash}p{5.5cm} >{\raggedright\arraybackslash}X}
\toprule
\textbf{ID} & \textbf{Como...} & \textbf{Deseo...} & \textbf{Para...} \\
\midrule
HU-O07.1 & SRE & monitorear CPU y memoria de cada pod en tiempo real & detectar saturación \\
\midrule
HU-O07.2 & SRE & trackear la tasa de errores 500 por endpoint & identificar servicios degradados \\
\midrule
HU-O07.3 & SRE & monitorear el pool de conexiones de la base de datos & prevenir agotamiento \\
\midrule
HU-O07.4 & Sistema & alertar automáticamente si la memoria supera el 80\% & activar respuesta proactiva \\
\midrule
HU-O07.5 & Sistema & auto-escalar los pods cuando la CPU supere el 70\% & mantener la disponibilidad \\
\bottomrule
\end{tabularx}
\end{table}

\vspace{4pt}
\noindent\rule{\linewidth}{0.4pt}
\vspace{4pt}


\subsection{CU-O08: Atender Ticket de Bug Reportado por Desarrollador}

\begin{table}[H]
\footnotesize
\begin{tabularx}{\linewidth}{>{\bfseries\raggedright\arraybackslash}p{2.5cm} >{\raggedright\arraybackslash}X}
\toprule
\multicolumn{2}{l}{\textbf{CU-O08: Atender Ticket de Bug Reportado por Desarrollador} \hfill \textit{Prioridad: 9 \quad | \quad Operativo}} \\
\midrule
Actor Principal & Developer Advocate \\
\midrule
Propósito & Triar y resolver bugs reportados por la comunidad. \\
\midrule
Precondición & Un desarrollador externo reporta un bug en el portal. \\
\midrule
Flujo Principal & El sistema auto-clasifica la severidad con NLP. Asigna prioridad. Inicia el SLA tracker. El Developer Advocate investiga y corrige. Notifica la resolución al reportante. \\
\midrule
Postcondición & bug resuelto. \\
\bottomrule
\end{tabularx}
\end{table}

\vspace{2pt}
\begin{table}[H]
\footnotesize
\begin{tabularx}{\linewidth}{>{\raggedright\arraybackslash}p{1.5cm} >{\raggedright\arraybackslash}p{2.2cm} >{\raggedright\arraybackslash}p{5.5cm} >{\raggedright\arraybackslash}X}
\toprule
\textbf{ID} & \textbf{Como...} & \textbf{Deseo...} & \textbf{Para...} \\
\midrule
HU-O08.1 & Desarrollador externo & reportar un bug con pasos para reproducirlo & facilitar el diagnóstico \\
\midrule
HU-O08.2 & Sistema & auto-clasificar la severidad del bug con NLP & priorizar los críticos \\
\midrule
HU-O08.3 & Developer Advocate & ver el SLA tracker indicando el tiempo restante & cumplir con la respuesta en menos de 4 horas \\
\midrule
HU-O08.4 & Developer Advocate & linkear el commit de fix al ticket original & mantener trazabilidad \\
\midrule
HU-O08.5 & Sistema & notificar al reportante automáticamente cuando el bug esté resuelto & cerrar el ciclo \\
\bottomrule
\end{tabularx}
\end{table}

\vspace{4pt}
\noindent\rule{\linewidth}{0.4pt}
\vspace{4pt}


\subsection{CU-O09: Ejecutar Backup y Prueba de Restauración}

\begin{table}[H]
\footnotesize
\begin{tabularx}{\linewidth}{>{\bfseries\raggedright\arraybackslash}p{2.5cm} >{\raggedright\arraybackslash}X}
\toprule
\multicolumn{2}{l}{\textbf{CU-O09: Ejecutar Backup y Prueba de Restauración} \hfill \textit{Prioridad: 10 \quad | \quad Operativo}} \\
\midrule
Actor Principal & SRE / DevOps \\
\midrule
Propósito & Backup diario y prueba de restore semanal. \\
\midrule
Precondición & El sistema ejecuta pg\_dump programado diariamente. \\
\midrule
Flujo Principal & Encripta el backup con KMS. Lo sube a cold storage (S3 Glacier). Semanalmente restaura el backup a una instancia de prueba. Valida integridad de datos. Loggea el resultado. \\
\midrule
Postcondición & backup verificado. \\
\bottomrule
\end{tabularx}
\end{table}

\vspace{2pt}
\begin{table}[H]
\footnotesize
\begin{tabularx}{\linewidth}{>{\raggedright\arraybackslash}p{1.5cm} >{\raggedright\arraybackslash}p{2.2cm} >{\raggedright\arraybackslash}p{5.5cm} >{\raggedright\arraybackslash}X}
\toprule
\textbf{ID} & \textbf{Como...} & \textbf{Deseo...} & \textbf{Para...} \\
\midrule
HU-O09.1 & SRE & programar el backup diario a las 03:00 UTC & no impactar operaciones \\
\midrule
HU-O09.2 & Sistema & encriptar el backup con KMS antes de almacenarlo & cumplir con requisitos de seguridad \\
\midrule
HU-O09.3 & Sistema & subir el backup encriptado a cold storage & minimizar el costo de almacenamiento \\
\midrule
HU-O09.4 & SRE & ejecutar una restauración de prueba semanal & validar que el backup es íntegro \\
\midrule
HU-O09.5 & SRE & registrar un log de éxito o fallo de cada restore & auditar la efectividad del proceso \\
\bottomrule
\end{tabularx}
\end{table}

\vspace{4pt}
\noindent\rule{\linewidth}{0.4pt}
\vspace{4pt}


\subsection{CU-O10: Revisar Error Budget y Congelar Deploys}

\begin{table}[H]
\footnotesize
\begin{tabularx}{\linewidth}{>{\bfseries\raggedright\arraybackslash}p{2.5cm} >{\raggedright\arraybackslash}X}
\toprule
\multicolumn{2}{l}{\textbf{CU-O10: Revisar Error Budget y Congelar Deploys} \hfill \textit{Prioridad: 9 \quad | \quad Operativo}} \\
\midrule
Actor Principal & SRE / DevOps \\
\midrule
Propósito & Revisión semanal de SLI/SLO. \\
\midrule
Precondición & El sistema calcula el SLI de uptime de los últimos 30 días. \\
\midrule
Flujo Principal & Compara contra el SLO del 99.99\%. Computa el error budget consumido. Si supera el 80\% congela automáticamente el pipeline CI/CD. Notifica al CTO y team leads. \\
\midrule
Postcondición & riesgo controlado. \\
\bottomrule
\end{tabularx}
\end{table}

\vspace{2pt}
\begin{table}[H]
\footnotesize
\begin{tabularx}{\linewidth}{>{\raggedright\arraybackslash}p{1.5cm} >{\raggedright\arraybackslash}p{2.2cm} >{\raggedright\arraybackslash}p{5.5cm} >{\raggedright\arraybackslash}X}
\toprule
\textbf{ID} & \textbf{Como...} & \textbf{Deseo...} & \textbf{Para...} \\
\midrule
HU-O10.1 & SRE & calcular el SLI de uptime de los últimos 30 días & evaluar la confiabilidad real \\
\midrule
HU-O10.2 & Sistema & computar el error budget consumido en minutos & compararlo contra el umbral \\
\midrule
HU-O10.3 & SRE & comparar el consumo contra el umbral del 80\% & decidir si se debe congelar \\
\midrule
HU-O10.4 & Sistema & congelar automáticamente el pipeline de CI/CD & evitar nuevos deploys que aumenten el riesgo \\
\midrule
HU-O10.5 & Sistema & notificar a todos los stakeholders sobre la congelación & coordinar las acciones de remediación \\
\bottomrule
\end{tabularx}
\end{table}

\vspace{4pt}
\noindent\rule{\linewidth}{0.4pt}
\vspace{4pt}


\subsection{CU-O11: Publicar Changelog Semanal en Developer Portal}

\begin{table}[H]
\footnotesize
\begin{tabularx}{\linewidth}{>{\bfseries\raggedright\arraybackslash}p{2.5cm} >{\raggedright\arraybackslash}X}
\toprule
\multicolumn{2}{l}{\textbf{CU-O11: Publicar Changelog Semanal en Developer Portal} \hfill \textit{Prioridad: 7 \quad | \quad Operativo}} \\
\midrule
Actor Principal & Developer Advocate \\
\midrule
Propósito & Comunicar cambios de API a desarrolladores. \\
\midrule
Precondición & El sistema recolecta commits desde el último release. \\
\midrule
Flujo Principal & Categoriza cada cambio como feature, fix o breaking. Genera un changelog con Semantic Versioning. Publica en el portal. Envía email a suscriptores. \\
\midrule
Postcondición & changelog publicado. \\
\bottomrule
\end{tabularx}
\end{table}

\vspace{2pt}
\begin{table}[H]
\footnotesize
\begin{tabularx}{\linewidth}{>{\raggedright\arraybackslash}p{1.5cm} >{\raggedright\arraybackslash}p{2.2cm} >{\raggedright\arraybackslash}p{5.5cm} >{\raggedright\arraybackslash}X}
\toprule
\textbf{ID} & \textbf{Como...} & \textbf{Deseo...} & \textbf{Para...} \\
\midrule
HU-O11.1 & Developer Advocate & recolectar commits automáticamente desde el último release & ahorrar tiempo manual \\
\midrule
HU-O11.2 & Sistema & categorizar cada commit por tipo & estructurar el changelog \\
\midrule
HU-O11.3 & Developer Advocate & generar la nueva versión siguiendo Semantic Versioning & comunicar claramente el impacto \\
\midrule
HU-O11.4 & Sistema & publicar el changelog en el Developer Portal & que los desarrolladores lo consulten \\
\midrule
HU-O11.5 & Sistema & enviar un email con el changelog a los suscriptores & mantenerlos informados proactivamente \\
\bottomrule
\end{tabularx}
\end{table}

\vspace{4pt}
\noindent\rule{\linewidth}{0.4pt}
\vspace{4pt}


\subsection{CU-O12: Rotar Credenciales y Auditar Accesos IAM}

\begin{table}[H]
\footnotesize
\begin{tabularx}{\linewidth}{>{\bfseries\raggedright\arraybackslash}p{2.5cm} >{\raggedright\arraybackslash}X}
\toprule
\multicolumn{2}{l}{\textbf{CU-O12: Rotar Credenciales y Auditar Accesos IAM} \hfill \textit{Prioridad: 9 \quad | \quad Operativo}} \\
\midrule
Actor Principal & SRE / DevOps \\
\midrule
Propósito & Rotación mensual de secretos y auditoría trimestral. \\
\midrule
Precondición & El sistema identifica credenciales próximas a expirar. \\
\midrule
Flujo Principal & Genera nuevas credenciales en el vault. Actualiza los servicios con las nuevas. Verifica que todo funcione. Desactiva las credenciales antiguas tras un período de gracia de 24 horas. Loggea toda la operación. \\
\midrule
Postcondición & credenciales rotadas. \\
\bottomrule
\end{tabularx}
\end{table}

\vspace{2pt}
\begin{table}[H]
\footnotesize
\begin{tabularx}{\linewidth}{>{\raggedright\arraybackslash}p{1.5cm} >{\raggedright\arraybackslash}p{2.2cm} >{\raggedright\arraybackslash}p{5.5cm} >{\raggedright\arraybackslash}X}
\toprule
\textbf{ID} & \textbf{Como...} & \textbf{Deseo...} & \textbf{Para...} \\
\midrule
HU-O12.1 & SRE & rotar las API Keys de los clientes mensualmente & reducir la ventana de exposición \\
\midrule
HU-O12.2 & Sistema & actualizar los secrets en K8s automáticamente & que los pods usen las nuevas credenciales \\
\midrule
HU-O12.3 & SRE & verificar que todos los servicios funcionen con las nuevas credenciales & evitar disrupciones \\
\midrule
HU-O12.4 & Sistema & desactivar las credenciales antiguas tras 24 horas & completar la rotación \\
\midrule
HU-O12.5 & CTO & generar una auditoría IAM trimestral con todos los accesos & revisar y revocar los innecesarios \\
\bottomrule
\end{tabularx}
\end{table}

\vspace{4pt}
\noindent\rule{\linewidth}{0.4pt}
\vspace{4pt}


\subsection{CU-O13: Realizar Post-Mortem de Incidente}

\begin{table}[H]
\footnotesize
\begin{tabularx}{\linewidth}{>{\bfseries\raggedright\arraybackslash}p{2.5cm} >{\raggedright\arraybackslash}X}
\toprule
\multicolumn{2}{l}{\textbf{CU-O13: Realizar Post-Mortem de Incidente} \hfill \textit{Prioridad: 8 \quad | \quad Operativo}} \\
\midrule
Actor Principal & SRE / DevOps \\
\midrule
Propósito & Análisis blameless post-incidente Sev1/Sev2. \\
\midrule
Precondición & Tras resolver un incidente, el SRE agenda una reunión blameless. \\
\midrule
Flujo Principal & El sistema genera automáticamente un timeline completo. Se identifica la root cause con la técnica de los 5 whys. Se documentan los aprendizajes. Se crean tickets de acción. Se publica el reporte en la wiki. \\
\midrule
Postcondición & aprendizajes documentados. \\
\bottomrule
\end{tabularx}
\end{table}

\vspace{2pt}
\begin{table}[H]
\footnotesize
\begin{tabularx}{\linewidth}{>{\raggedright\arraybackslash}p{1.5cm} >{\raggedright\arraybackslash}p{2.2cm} >{\raggedright\arraybackslash}p{5.5cm} >{\raggedright\arraybackslash}X}
\toprule
\textbf{ID} & \textbf{Como...} & \textbf{Deseo...} & \textbf{Para...} \\
\midrule
HU-O13.1 & Sistema & generar un timeline automático del incidente & reconstruir la secuencia exacta \\
\midrule
HU-O13.2 & SRE & documentar qué funcionó bien y qué no durante la respuesta & mejorar los runbooks \\
\midrule
HU-O13.3 & SRE & identificar la root cause utilizando la técnica de los 5 whys & prevenir recurrencia \\
\midrule
HU-O13.4 & SRE & crear tickets de acción en Linear desde el post-mortem & asegurar seguimiento \\
\midrule
HU-O13.5 & SRE & publicar el post-mortem en la wiki interna & compartir el aprendizaje con toda la organización \\
\bottomrule
\end{tabularx}
\end{table}

\vspace{4pt}
\noindent\rule{\linewidth}{0.4pt}
\vspace{4pt}


\subsection{CU-O14: Documentar FAQ en Base de Conocimientos del Chatbot}

\begin{table}[H]
\footnotesize
\begin{tabularx}{\linewidth}{>{\bfseries\raggedright\arraybackslash}p{2.5cm} >{\raggedright\arraybackslash}X}
\toprule
\multicolumn{2}{l}{\textbf{CU-O14: Documentar FAQ en Base de Conocimientos del Chatbot} \hfill \textit{Prioridad: 7 \quad | \quad Operativo}} \\
\midrule
Actor Principal & Developer Advocate \\
\midrule
Propósito & Alimentar knowledge base desde tickets resueltos. \\
\midrule
Precondición & El sistema detecta un tema recurrente al acumularse 3 o más tickets similares. \\
\midrule
Flujo Principal & Sugiere un borrador de FAQ. El Developer Advocate revisa, edita y publica con screenshots. Actualiza los embeddings del chatbot. Mide la tasa de deflexión resultante. \\
\midrule
Postcondición & KB enriquecida. \\
\bottomrule
\end{tabularx}
\end{table}

\vspace{2pt}
\begin{table}[H]
\footnotesize
\begin{tabularx}{\linewidth}{>{\raggedright\arraybackslash}p{1.5cm} >{\raggedright\arraybackslash}p{2.2cm} >{\raggedright\arraybackslash}p{5.5cm} >{\raggedright\arraybackslash}X}
\toprule
\textbf{ID} & \textbf{Como...} & \textbf{Deseo...} & \textbf{Para...} \\
\midrule
HU-O14.1 & Sistema & detectar temas recurrentes cuando aparecen 3 o más tickets similares & proponer nueva documentación \\
\midrule
HU-O14.2 & Sistema & generar un borrador de FAQ automáticamente & agilizar la redacción \\
\midrule
HU-O14.3 & Developer Advocate & revisar el borrador y publicarlo & garantizar calidad editorial \\
\midrule
HU-O14.4 & Sistema & actualizar los embeddings del chatbot con la nueva FAQ & mejorar la precisión de respuestas \\
\midrule
HU-O14.5 & Developer Advocate & medir la tasa de deflexión después de publicar la FAQ & evaluar su efectividad \\
\bottomrule
\end{tabularx}
\end{table}

\vspace{4pt}
\noindent\rule{\linewidth}{0.4pt}
\vspace{4pt}

\newpage
\section{Casos de Uso Operativos Adicionales}

\subsection{CU-O15: Gestionar Sprint en Linear}

\begin{table}[H]
\footnotesize
\begin{tabularx}{\linewidth}{>{\bfseries\raggedright\arraybackslash}p{2.5cm} >{\raggedright\arraybackslash}X}
\toprule
\multicolumn{2}{l}{\textbf{CU-O15: Gestionar Sprint en Linear} \hfill \textit{Prioridad: 8 \quad | \quad Operativo}} \\
\midrule
Actor Principal & Todo el equipo de ingeniería \\
\midrule
Propósito & Gestionar el backlog, planificar sprints y medir la productividad del equipo. \\
\midrule
Precondición & Los leads de equipo crean tickets de trabajo en Linear desde el roadmap priorizado. \\
\midrule
Flujo Principal & Cada sprint de 2 semanas inicia con planning meeting para asignar estimaciones. Diariamente cada miembro actualiza el estado de sus tickets en Slack standup. Al cierre del sprint se realiza retrospectiva con revisión de velocity y burndown. Los aprendizajes se documentan y se ajusta el backlog del siguiente sprint. \\
\midrule
Postcondición & sprint cerrado con métricas y aprendizajes documentados. \\
\bottomrule
\end{tabularx}
\end{table}

\vspace{2pt}
\begin{table}[H]
\footnotesize
\begin{tabularx}{\linewidth}{>{\raggedright\arraybackslash}p{1.5cm} >{\raggedright\arraybackslash}p{2.2cm} >{\raggedright\arraybackslash}p{5.5cm} >{\raggedright\arraybackslash}X}
\toprule
\textbf{ID} & \textbf{Como...} & \textbf{Deseo...} & \textbf{Para...} \\
\midrule
HU-O15.1 & Tech Lead & crear tickets de trabajo con título, descripción y prioridad RICE & que el equipo tenga claridad sobre qué construir \\
\midrule
HU-O15.2 & Desarrollador & actualizar el estado de mis tickets desde Slack sin abrir Linear & reducir el cambio de contexto \\
\midrule
HU-O15.3 & Tech Lead & visualizar el burndown chart en tiempo real durante el sprint & identificar riesgos de entrega a tiempo \\
\midrule
HU-O15.4 & Tech Lead & documentar los aprendizajes de la retrospectiva en la wiki & que el equipo no repita los mismos errores \\
\midrule
HU-O15.5 & CTO & ver un dashboard de velocity y cycle time por equipo & evaluar la productividad y predecir entregables \\
\bottomrule
\end{tabularx}
\end{table}

\vspace{4pt}
\noindent\rule{\linewidth}{0.4pt}
\vspace{4pt}


\subsection{CU-O16: Validar Especificaciones OpenAPI en CI/CD}

\begin{table}[H]
\footnotesize
\begin{tabularx}{\linewidth}{>{\bfseries\raggedright\arraybackslash}p{2.5cm} >{\raggedright\arraybackslash}X}
\toprule
\multicolumn{2}{l}{\textbf{CU-O16: Validar Especificaciones OpenAPI en CI/CD} \hfill \textit{Prioridad: 9 \quad | \quad Operativo}} \\
\midrule
Actor Principal & CTO / SRE \\
\midrule
Propósito & Asegurar que todas las especificaciones de API cumplan el estándar OpenAPI 3.1. \\
\midrule
Precondición & Al abrir un PR que modifique archivos .yaml de especificación, el pipeline de CI/CD ejecuta Spectral linter automáticamente. \\
\midrule
Flujo Principal & Verifica reglas: todos los endpoints tienen descripción y ejemplos, no hay breaking changes sin bump de versión major, los schemas de request/response son válidos. Si hay errores, el PR se bloquea hasta que se corrijan. \\
\midrule
Postcondición & spec validada y conforme al estándar. \\
\bottomrule
\end{tabularx}
\end{table}

\vspace{2pt}
\begin{table}[H]
\footnotesize
\begin{tabularx}{\linewidth}{>{\raggedright\arraybackslash}p{1.5cm} >{\raggedright\arraybackslash}p{2.2cm} >{\raggedright\arraybackslash}p{5.5cm} >{\raggedright\arraybackslash}X}
\toprule
\textbf{ID} & \textbf{Como...} & \textbf{Deseo...} & \textbf{Para...} \\
\midrule
HU-O16.1 & CTO & ejecutar un linter automático sobre cada spec OpenAPI en el PR & evitar que specs no conformes lleguen a producción \\
\midrule
HU-O16.2 & Desarrollador & ver los errores de lint en el diff del PR & corregirlos antes de pedir revisión \\
\midrule
HU-O16.3 & Sistema & detectar automáticamente breaking changes comparando specs & exigir un bump de versión major \\
\midrule
HU-O16.4 & Developer Advocate & verificar que todos los endpoints tengan ejemplos request/response & garantizar documentación completa \\
\midrule
HU-O16.5 & CTO & bloquear el merge si el linter encuentra errores & mantener la calidad del contrato de API \\
\bottomrule
\end{tabularx}
\end{table}

\vspace{4pt}
\noindent\rule{\linewidth}{0.4pt}
\vspace{4pt}


\subsection{CU-O17: Validar Data Contracts en Pipeline ETL}

\begin{table}[H]
\footnotesize
\begin{tabularx}{\linewidth}{>{\bfseries\raggedright\arraybackslash}p{2.5cm} >{\raggedright\arraybackslash}X}
\toprule
\multicolumn{2}{l}{\textbf{CU-O17: Validar Data Contracts en Pipeline ETL} \hfill \textit{Prioridad: 10 \quad | \quad Operativo}} \\
\midrule
Actor Principal & Data Engineer \\
\midrule
Propósito & Asegurar que los datos externos cumplan los contratos de schema definidos. \\
\midrule
Precondición & El Data Engineer define y versiona en Git los schema contracts para cada fuente de datos (proveedor aeronáutico, meteorología, etc.). \\
\midrule
Flujo Principal & El pipeline de ingesta ejecuta validación automática contra el contrato antes de cargar a staging: tipos de datos, nulabilidad de columnas críticas, rangos de valores plausibles y estructura del archivo. Si hay schema drift del proveedor, se alerta inmediatamente y se pausa la ingesta. \\
\midrule
Postcondición & datos validados contra contrato o ingesta pausada con alerta. \\
\bottomrule
\end{tabularx}
\end{table}

\vspace{2pt}
\begin{table}[H]
\footnotesize
\begin{tabularx}{\linewidth}{>{\raggedright\arraybackslash}p{1.5cm} >{\raggedright\arraybackslash}p{2.2cm} >{\raggedright\arraybackslash}p{5.5cm} >{\raggedright\arraybackslash}X}
\toprule
\textbf{ID} & \textbf{Como...} & \textbf{Deseo...} & \textbf{Para...} \\
\midrule
HU-O17.1 & Data Engineer & definir el schema contract de cada fuente en un archivo YAML versionado & tener trazabilidad de cambios \\
\midrule
HU-O17.2 & Sistema & validar que los datos entrantes cumplan el contrato antes de cargar a staging & evitar corrupción downstream \\
\midrule
HU-O17.3 & Data Engineer & recibir una alerta si el proveedor cambió el schema sin notificar & actualizar el pipeline oportunamente \\
\midrule
HU-O17.4 & Sistema & pausar la ingesta automáticamente si la validación de contrato falla & prevenir cascada de errores \\
\midrule
HU-O17.5 & CTO & ver un dashboard de conformidad de data contracts por fuente & monitorear la estabilidad del pipeline \\
\bottomrule
\end{tabularx}
\end{table}

\vspace{4pt}
\noindent\rule{\linewidth}{0.4pt}
\vspace{4pt}


\subsection{CU-O18: Ejecutar Pipeline de Seguridad SAST/DAST en CI/CD}

\begin{table}[H]
\footnotesize
\begin{tabularx}{\linewidth}{>{\bfseries\raggedright\arraybackslash}p{2.5cm} >{\raggedright\arraybackslash}X}
\toprule
\multicolumn{2}{l}{\textbf{CU-O18: Ejecutar Pipeline de Seguridad SAST/DAST en CI/CD} \hfill \textit{Prioridad: 9 \quad | \quad Operativo}} \\
\midrule
Actor Principal & SRE / DevOps \\
\midrule
Propósito & Detectar vulnerabilidades de seguridad de forma automática en cada cambio de código. \\
\midrule
Precondición & En cada PR, el pipeline de CI/CD ejecuta SonarQube para análisis estático de código (SAST) y lanza un escaneo con OWASP ZAP contra una instancia efímera del servicio modificado (DAST). \\
\midrule
Flujo Principal & Si se detectan vulnerabilidades con severidad critical o high, el PR se bloquea y se notifica al autor. El reporte de seguridad se adjunta al release. \\
\midrule
Postcondición & código validado sin vulnerabilidades críticas en el diff. \\
\bottomrule
\end{tabularx}
\end{table}

\vspace{2pt}
\begin{table}[H]
\footnotesize
\begin{tabularx}{\linewidth}{>{\raggedright\arraybackslash}p{1.5cm} >{\raggedright\arraybackslash}p{2.2cm} >{\raggedright\arraybackslash}p{5.5cm} >{\raggedright\arraybackslash}X}
\toprule
\textbf{ID} & \textbf{Como...} & \textbf{Deseo...} & \textbf{Para...} \\
\midrule
HU-O18.1 & SRE & ejecutar SonarQube automáticamente en cada PR & detectar vulnerabilidades antes del merge \\
\midrule
HU-O18.2 & SRE & lanzar un escaneo OWASP ZAP contra una instancia efímera del servicio & detectar vulnerabilidades en tiempo de ejecución \\
\midrule
HU-O18.3 & Desarrollador & ver el resultado del escaneo directamente en el PR & corregir issues antes de que se bloquee el merge \\
\midrule
HU-O18.4 & Sistema & bloquear automáticamente el merge si hay vulnerabilidades críticas & evitar que el código inseguro llegue a producción \\
\midrule
HU-O18.5 & CTO & generar un reporte de seguridad firmado por release & presentarlo como evidencia en auditorías SOC 2 \\
\bottomrule
\end{tabularx}
\end{table}

\vspace{4pt}
\noindent\rule{\linewidth}{0.4pt}
\vspace{4pt}


\subsection{CU-O19: Analizar y Optimizar Costos Cloud}

\begin{table}[H]
\footnotesize
\begin{tabularx}{\linewidth}{>{\bfseries\raggedright\arraybackslash}p{2.5cm} >{\raggedright\arraybackslash}X}
\toprule
\multicolumn{2}{l}{\textbf{CU-O19: Analizar y Optimizar Costos Cloud} \hfill \textit{Prioridad: 8 \quad | \quad Operativo}} \\
\midrule
Actor Principal & CTO / SRE \\
\midrule
Propósito & Monitorear, analizar y reducir el gasto en infraestructura cloud. \\
\midrule
Precondición & Semanalmente, el SRE audita la factura cloud desglosada por servicio. \\
\midrule
Flujo Principal & El sistema compara el gasto real contra el forecast mensual y detecta anomalías con un algoritmo de desviación estándar. Identifica recursos ociosos (instancias sin tráfico, volúmenes no adjuntos, IPs no asignadas) y recursos sobre-aprovisionados. Genera un reporte con recomendaciones de rightsizing, reserved instances y eliminación de waste. El CTO revisa y aprueba las acciones. \\
\midrule
Postcondición & gasto cloud optimizado y documentado. \\
\bottomrule
\end{tabularx}
\end{table}

\vspace{2pt}
\begin{table}[H]
\footnotesize
\begin{tabularx}{\linewidth}{>{\raggedright\arraybackslash}p{1.5cm} >{\raggedright\arraybackslash}p{2.2cm} >{\raggedright\arraybackslash}p{5.5cm} >{\raggedright\arraybackslash}X}
\toprule
\textbf{ID} & \textbf{Como...} & \textbf{Deseo...} & \textbf{Para...} \\
\midrule
HU-O19.1 & CTO & visualizar los costos desglosados por servicio cloud en tiempo real & identificar los mayores drivers de gasto \\
\midrule
HU-O19.2 & SRE & recibir una alerta si el gasto diario se desvía más del 20\% del forecast & investigar fugas a tiempo \\
\midrule
HU-O19.3 & Sistema & detectar automáticamente recursos cloud ociosos & proponer su eliminación o redimensionamiento \\
\midrule
HU-O19.4 & SRE & generar recomendaciones de ahorro con estimación de impacto & priorizar las de mayor retorno \\
\midrule
HU-O19.5 & CTO & comparar la tendencia mensual de costos contra el presupuesto & tomar decisiones de inversión en infraestructura \\
\bottomrule
\end{tabularx}
\end{table}

\vspace{4pt}
\noindent\rule{\linewidth}{0.4pt}
\vspace{4pt}


\subsection{CU-O20: Ejecutar Experimentos de Deep Learning}

\begin{table}[H]
\footnotesize
\begin{tabularx}{\linewidth}{>{\bfseries\raggedright\arraybackslash}p{2.5cm} >{\raggedright\arraybackslash}X}
\toprule
\multicolumn{2}{l}{\textbf{CU-O20: Ejecutar Experimentos de Deep Learning} \hfill \textit{Prioridad: 7 \quad | \quad Operativo}} \\
\midrule
Actor Principal & ML Engineer \\
\midrule
Propósito & Explorar y validar técnicas de Deep Learning para OCR documental y Speech-to-Text. \\
\midrule
Precondición & El ML Engineer diseña experimentos de DL en sprints de experimentación de 4 semanas. \\
\midrule
Flujo Principal & Para OCR: entrena modelos con redes neuronales (CRNN/TrOCR) sobre documentos de compliance, facturas de proveedores y logs de auditoría escaneados. Para Speech-to-Text: despliega modelos Whisper/DeepSpeech, integra con el pipeline de grabación de entrevistas de feedback y mide Word Error Rate. Todos los experimentos se registran en MLflow con hiperparámetros, métricas y artefactos. Se comparan contra baseline y se decide si promover a desarrollo de producto. \\
\midrule
Postcondición & experimentos documentados con recomendación go/no-go. \\
\bottomrule
\end{tabularx}
\end{table}

\vspace{2pt}
\begin{table}[H]
\footnotesize
\begin{tabularx}{\linewidth}{>{\raggedright\arraybackslash}p{1.5cm} >{\raggedright\arraybackslash}p{2.2cm} >{\raggedright\arraybackslash}p{5.5cm} >{\raggedright\arraybackslash}X}
\toprule
\textbf{ID} & \textbf{Como...} & \textbf{Deseo...} & \textbf{Para...} \\
\midrule
HU-O20.1 & ML Engineer & entrenar un modelo OCR sobre documentos de compliance escaneados & automatizar la extracción de datos para auditoría \\
\midrule
HU-O20.2 & ML Engineer & comparar la precisión del OCR automático contra la entrada manual & decidir si es viable para producción \\
\midrule
HU-O20.3 & ML Engineer & desplegar un modelo Whisper para transcripción de entrevistas de feedback & extraer texto de las grabaciones \\
\midrule
HU-O20.4 & ML Engineer & medir el Word Error Rate del transcriptor contra transcripciones manuales & validar la calidad del modelo \\
\midrule
HU-O20.5 & ML Engineer & loguear métricas, hiperparámetros y artefactos de cada experimento en MLflow & tener trazabilidad completa del proceso de I+D \\
\bottomrule
\end{tabularx}
\end{table}

\vspace{4pt}
\noindent\rule{\linewidth}{0.4pt}
\vspace{4pt}


\vfill
\begin{center}
    {\small SkyAnalytics Inc. -- Perfil Estratégico Corporativo · Documento Integral · 2026}
\end{center}

\end{document}
